% Recovered from the previously completed PDF; no new mathematical content.
% Prose is editable; formula boxes retain original symbol positions.
% Build from repository root: xelatex -no-shell-escape -halt-on-error -output-directory=... print/fudan/NAME.tex
\documentclass[11pt,a4paper]{article}
\usepackage[margin=23mm,headheight=14pt,headsep=8.5mm,footskip=10mm]{geometry}
\usepackage{amsmath,amssymb}
\DeclareMathAlphabet{\mathsf}{OT1}{cmss}{m}{n}
\usepackage[no-math]{fontspec}
\usepackage{xeCJK}
\setmainfont{Latin Modern Roman}
\setCJKmainfont{FandolSong-Regular.otf}[Path=assets/fonts/fandol/,BoldFont=FandolSong-Bold.otf,ItalicFont=FandolSong-Regular.otf,Script=CustomDefault]
\xeCJKsetup{PunctStyle=quanjiao,CheckSingle=true}
\usepackage{setspace}\setstretch{1.12}
\setlength{\parindent}{2em}\setlength{\parskip}{3pt}
\emergencystretch=6em
\hfuzz=9pt
\usepackage{titlesec}
\titleformat{\section}{\centering\normalfont\Large\bfseries}{}{0pt}{}
\titlespacing*{\section}{0pt}{1.8\baselineskip}{1.0\baselineskip}
\usepackage{fancyhdr}\pagestyle{fancy}\fancyhf{}
\fancyhead[L]{\small 数学手记·高等代数}\fancyhead[R]{\small 17S09–17S15}
\fancyfoot[C]{\thepage}\renewcommand{\headrulewidth}{0.4pt}
\usepackage[hidelinks,unicode]{hyperref}
\hypersetup{pdftitle={17S09–17S15：题目、解答与小结},pdfauthor={Vincent Arqain}}
\newcommand{\RB}[4]{\raisebox{0pt}[#2bp][#3bp]{\hbox to #1bp{#4\hss}}}
\newcommand{\RG}[5]{\hbox to 0pt{\kern#1bp\raisebox{#2bp}[0pt][0pt]{\fontsize{#3bp}{#3bp}\csname font#4\endcsname\char#5}\hss}}
\newcommand{\RT}[4]{\hbox to 0pt{\kern#1bp\raisebox{#2bp}[0pt][0pt]{\fontsize{#3bp}{#3bp}\selectfont#4}\hss}}
\newcommand{\RS}[4]{\hbox to 0pt{\kern#1bp\raisebox{#2bp}[0pt][0pt]{\fontsize{#3bp}{#3bp}\selectfont\ensuremath{#4}}\hss}}
\newcommand{\RR}[4]{\hbox to 0pt{\kern#1bp\raisebox{#2bp}[0pt][0pt]{\rule{#3bp}{#4bp}}\hss}}
\newcommand{\fontmi}{\usefont{OML}{cmm}{m}{it}}
\newcommand{\fontsy}{\usefont{OMS}{cmsy}{m}{n}}
\newcommand{\fontex}{\usefont{OMX}{cmex}{m}{n}}
\newcommand{\fontrm}{\usefont{OT1}{cmr}{m}{n}}
\newcommand{\fontbf}{\usefont{OT1}{cmr}{bx}{n}}
\newcommand{\fontam}{\usefont{U}{msa}{m}{n}}
\newcommand{\fontbb}{\usefont{U}{msb}{m}{n}}
\widowpenalty=10000\clubpenalty=10000\raggedbottom
\begin{document}
\begin{center}{\LARGE\bfseries 17S09–17S15\par}\vspace{8pt}复旦大学高等代数习题集·题目与解答\end{center}
% Original page 1.
复旦大学高等代数习题集题目与证明

题目选自谢启鸿《复旦大学高等代数习题集》（2023版），§1.7，第45–46页。

记号：\allowbreak{}\hskip.1em \RB{5.244}{7.023}{0.500}{%
\RG{0.000}{0.004}{7.9701}{mi}{73}%
}\hskip.1em\allowbreak{} 表示恒等算子或单位矩阵；\allowbreak{}\hskip.1em \RB{12.809}{8.558}{0.500}{%
\RG{0.000}{0.004}{7.9701}{mi}{66}\RG{7.719}{3.468}{7.9701}{rm}{84}%
}\hskip.1em\allowbreak{} 表示转置，\allowbreak{}\hskip.1em \RB{11.534}{7.075}{0.500}{%
\RG{0.000}{0.004}{7.9701}{mi}{66}\RG{7.720}{3.468}{7.9701}{sy}{3}%
}\hskip.1em\allowbreak{} 表示共轭转置。

\section*{2017S09}
设\allowbreak{}\hskip.1em \RB{8.715}{7.953}{0.727}{%
\RG{0.000}{0.002}{10.9091}{mi}{86}%
}\hskip.1em\allowbreak{} 是数域\allowbreak{}\hskip.1em \RB{8.889}{7.975}{0.500}{%
\RG{0.000}{0.002}{10.9091}{bb}{75}%
}\hskip.1em\allowbreak{} 上的\allowbreak{}\hskip.1em \RB{6.555}{5.313}{0.607}{%
\RG{0.000}{0.002}{10.9091}{mi}{110}%
}\hskip.1em\allowbreak{} 维线性空间，\allowbreak{}\hskip.1em \RB{7.122}{5.313}{2.855}{%
\RG{0.000}{0.002}{10.9091}{mi}{39}%
}\hskip.1em\allowbreak{} 是\allowbreak{}\hskip.1em \RB{8.715}{7.953}{0.727}{%
\RG{0.000}{0.002}{10.9091}{mi}{86}%
}\hskip.1em\allowbreak{} 上的线性变换，\allowbreak{}\hskip.1em \RB{44.312}{8.684}{3.236}{%
\RG{0.000}{0.002}{10.9091}{mi}{102}\RG{6.430}{0.002}{10.9091}{rm}{40}\RG{10.584}{0.002}{10.9091}{mi}{116}\RG{14.441}{0.002}{10.9091}{rm}{41}\RG{18.595}{0.002}{10.9091}{mi}{59}\RG{23.373}{0.002}{10.9091}{mi}{109}\RG{32.725}{0.002}{10.9091}{rm}{40}\RG{36.878}{0.002}{10.9091}{mi}{116}\RG{40.735}{0.002}{10.9091}{rm}{41}%
}\hskip.1em\allowbreak{} 分别是\allowbreak{}\hskip.1em \RB{7.122}{5.313}{2.855}{%
\RG{0.000}{0.002}{10.9091}{mi}{39}%
}\hskip.1em\allowbreak{} 的特征多项式和极小多项式。以下各小问的假设是独立的。

(i)设\allowbreak{}\hskip.1em \RB{109.841}{8.686}{3.235}{%
\RG{0.000}{0.004}{10.9091}{mi}{102}\RG{6.430}{0.004}{10.9091}{rm}{40}\RG{10.583}{0.004}{10.9091}{mi}{116}\RG{14.441}{0.004}{10.9091}{rm}{41}\RG{21.620}{0.004}{10.9091}{rm}{61}\RG{32.964}{0.004}{10.9091}{mi}{109}\RG{42.308}{0.004}{10.9091}{rm}{40}\RG{46.461}{0.004}{10.9091}{mi}{116}\RG{50.319}{0.004}{10.9091}{rm}{41}\RG{57.497}{0.004}{10.9091}{rm}{61}\RG{68.843}{0.004}{10.9091}{rm}{40}\RG{72.996}{0.004}{10.9091}{mi}{116}\RG{79.277}{0.004}{10.9091}{sy}{0}\RG{90.186}{0.004}{10.9091}{mi}{21}\RG{96.418}{-1.632}{7.9701}{rm}{48}\RG{100.736}{0.004}{10.9091}{rm}{41}\RG{104.890}{3.962}{7.9701}{mi}{110}%
}\hskip.1em\allowbreak{} ，求\allowbreak{}\hskip.1em \RB{8.715}{7.955}{0.725}{%
\RG{0.000}{0.004}{10.9091}{mi}{86}%
}\hskip.1em\allowbreak{} 的所有\allowbreak{}\hskip.1em \RB{7.122}{5.315}{2.853}{%
\RG{0.000}{0.004}{10.9091}{mi}{39}%
}\hskip.1em\allowbreak{} -不变子空间。

(ii)设\allowbreak{}\hskip.1em \RB{96.215}{8.686}{3.235}{%
\RG{0.000}{0.004}{10.9091}{mi}{102}\RG{6.430}{0.004}{10.9091}{rm}{40}\RG{10.583}{0.004}{10.9091}{mi}{116}\RG{14.441}{0.004}{10.9091}{rm}{41}\RG{23.431}{0.004}{10.9091}{rm}{61}\RG{36.578}{0.004}{10.9091}{mi}{102}\RG{41.843}{-1.633}{7.9701}{rm}{49}\RG{46.161}{0.004}{10.9091}{rm}{40}\RG{50.315}{0.004}{10.9091}{mi}{116}\RG{54.172}{0.004}{10.9091}{rm}{41}\RG{60.145}{0.004}{10.9091}{sy}{1}\RG{64.988}{0.004}{10.9091}{sy}{1}\RG{69.843}{0.004}{10.9091}{sy}{1}\RG{74.690}{0.004}{10.9091}{mi}{102}\RG{79.955}{-1.633}{7.9701}{mi}{107}\RG{84.627}{0.004}{10.9091}{rm}{40}\RG{88.781}{0.004}{10.9091}{mi}{116}\RG{92.638}{0.004}{10.9091}{rm}{41}%
}\hskip.1em\allowbreak{} ，其中各\allowbreak{}\hskip.1em \RB{42.862}{8.686}{3.224}{%
\RG{0.000}{0.004}{10.9091}{mi}{102}\RG{5.265}{-1.633}{7.9701}{mi}{105}\RG{13.188}{0.004}{10.9091}{sy}{50}\RG{25.298}{0.004}{10.9091}{bb}{75}\RG{33.783}{0.004}{10.9091}{rm}{91}\RG{36.813}{0.004}{10.9091}{mi}{116}\RG{40.671}{0.004}{10.9091}{rm}{93}%
}\hskip.1em\allowbreak{} 两两互素。令\allowbreak{}\hskip.1em \RB{66.435}{8.686}{3.235}{%
\RG{0.000}{0.004}{10.9091}{mi}{86}\RG{6.227}{-1.633}{7.9701}{mi}{105}\RG{14.149}{0.004}{10.9091}{rm}{61}\RG{27.295}{0.004}{10.9091}{rm}{107}\RG{32.652}{0.004}{10.9091}{rm}{101}\RG{37.397}{0.004}{10.9091}{rm}{114}\RG{43.356}{0.004}{10.9091}{mi}{102}\RG{48.621}{-1.633}{7.9701}{mi}{105}\RG{51.707}{0.004}{10.9091}{rm}{40}\RG{55.860}{0.004}{10.9091}{mi}{39}\RG{62.858}{0.004}{10.9091}{rm}{41}%
}\hskip.1em\allowbreak{} ，则\allowbreak{}\hskip.1em \RB{21.651}{7.955}{0.725}{%
\RG{0.000}{0.004}{10.9091}{mi}{86}\RG{13.439}{0.004}{10.9091}{rm}{61}%
}\hskip.1em\allowbreak{}

\allowbreak{}\hskip.1em \RB{60.915}{7.948}{2.212}{%
\RG{0.000}{-0.003}{10.9091}{mi}{86}\RG{6.227}{-1.639}{7.9701}{rm}{49}\RG{13.063}{-0.003}{10.9091}{sy}{8}\RG{24.059}{-0.003}{10.9091}{sy}{1}\RG{28.914}{-0.003}{10.9091}{sy}{1}\RG{33.768}{-0.003}{10.9091}{sy}{1}\RG{39.310}{-0.003}{10.9091}{sy}{8}\RG{50.311}{-0.003}{10.9091}{mi}{86}\RG{56.538}{-1.639}{7.9701}{mi}{107}%
}\hskip.1em\allowbreak{} 。任取\allowbreak{}\hskip.1em \RB{8.715}{7.948}{0.732}{%
\RG{0.000}{-0.003}{10.9091}{mi}{86}%
}\hskip.1em\allowbreak{} 的\allowbreak{}\hskip.1em \RB{7.122}{5.308}{2.859}{%
\RG{0.000}{-0.003}{10.9091}{mi}{39}%
}\hskip.1em\allowbreak{} -不变子空间\allowbreak{}\hskip.1em \RB{8.562}{7.948}{0.732}{%
\RG{0.000}{-0.003}{10.9091}{mi}{85}%
}\hskip.1em\allowbreak{} ，证明\allowbreak{}\hskip.1em \RB{86.197}{7.948}{2.212}{%
\RG{0.000}{-0.003}{10.9091}{mi}{85}\RG{11.656}{-0.003}{10.9091}{rm}{61}\RG{23.227}{-0.003}{10.9091}{mi}{85}\RG{30.482}{-1.639}{7.9701}{rm}{49}\RG{37.317}{-0.003}{10.9091}{sy}{8}\RG{48.314}{-0.003}{10.9091}{sy}{1}\RG{53.168}{-0.003}{10.9091}{sy}{1}\RG{58.023}{-0.003}{10.9091}{sy}{1}\RG{63.565}{-0.003}{10.9091}{sy}{8}\RG{74.566}{-0.003}{10.9091}{mi}{85}\RG{81.820}{-1.639}{7.9701}{mi}{107}%
}\hskip.1em\allowbreak{} ，其中\allowbreak{}\hskip.1em \RB{10.060}{7.948}{2.212}{%
\RG{0.000}{-0.003}{10.9091}{mi}{85}\RG{7.254}{-1.639}{7.9701}{mi}{105}%
}\hskip.1em\allowbreak{} 是\allowbreak{}\hskip.1em \RB{9.033}{7.948}{2.212}{%
\RG{0.000}{-0.003}{10.9091}{mi}{86}\RG{6.227}{-1.639}{7.9701}{mi}{105}%
}\hskip.1em\allowbreak{}的\allowbreak{}\hskip.1em \RB{7.122}{5.311}{2.856}{%
\RG{0.000}{-0.000}{10.9091}{mi}{39}%
}\hskip.1em\allowbreak{} -不变子空间。

(iii)设

\[
\RB{125.617}{17.366}{11.442}{%
\RG{72.986}{11.819}{7.9701}{mi}{107}\RG{69.286}{9.091}{9.0910}{ex}{89}\RG{116.162}{4.504}{7.9701}{mi}{114}\RG{119.650}{3.395}{5.9776}{mi}{105}\RG{0.000}{-0.000}{10.9091}{mi}{102}\RG{6.430}{-0.000}{10.9091}{rm}{40}\RG{10.584}{-0.000}{10.9091}{mi}{116}\RG{14.441}{-0.000}{10.9091}{rm}{41}\RG{21.620}{-0.000}{10.9091}{rm}{61}\RG{32.965}{-0.000}{10.9091}{mi}{109}\RG{42.308}{-0.000}{10.9091}{rm}{40}\RG{46.462}{-0.000}{10.9091}{mi}{116}\RG{50.320}{-0.000}{10.9091}{rm}{41}\RG{57.498}{-0.000}{10.9091}{rm}{61}\RG{81.347}{-0.000}{10.9091}{rm}{40}\RG{85.501}{-0.000}{10.9091}{mi}{116}\RG{91.782}{-0.000}{10.9091}{sy}{0}\RG{102.691}{-0.000}{10.9091}{mi}{21}\RG{112.008}{-0.000}{10.9091}{rm}{41}\RG{122.989}{-0.000}{10.9091}{mi}{59}\RG{108.923}{-1.636}{7.9701}{mi}{105}\RG{68.842}{-10.869}{7.9701}{mi}{105}\RG{71.473}{-10.869}{7.9701}{rm}{61}\RG{77.481}{-10.869}{7.9701}{rm}{49}%
}
\]
其中\allowbreak{}\hskip.1em \RB{67.601}{8.072}{2.605}{%
\RG{0.000}{0.001}{10.9091}{mi}{21}\RG{6.231}{-1.636}{7.9701}{rm}{49}\RG{10.550}{0.001}{10.9091}{mi}{59}\RG{15.328}{0.001}{10.9091}{mi}{58}\RG{20.117}{0.001}{10.9091}{mi}{58}\RG{24.895}{0.001}{10.9091}{mi}{58}\RG{29.684}{0.001}{10.9091}{mi}{59}\RG{34.474}{0.001}{10.9091}{mi}{21}\RG{40.707}{-1.636}{7.9701}{mi}{107}\RG{48.409}{0.001}{10.9091}{sy}{50}\RG{58.712}{0.001}{10.9091}{bb}{75}%
}\hskip.1em\allowbreak{} 两两不同，\allowbreak{}\hskip.1em \RB{27.418}{7.756}{2.209}{%
\RG{0.000}{0.001}{10.9091}{mi}{114}\RG{4.813}{-1.636}{7.9701}{mi}{105}\RG{10.930}{0.001}{10.9091}{sy}{21}\RG{22.445}{0.001}{10.9091}{rm}{49}%
}\hskip.1em\allowbreak{} 。求所有\allowbreak{}\hskip.1em \RB{7.122}{5.312}{2.856}{%
\RG{0.000}{0.001}{10.9091}{mi}{39}%
}\hskip.1em\allowbreak{} -不变子空间。

(iv)设\allowbreak{}\hskip.1em \RB{7.122}{5.312}{2.855}{%
\RG{0.000}{0.001}{10.9091}{mi}{39}%
}\hskip.1em\allowbreak{} 在某组基\allowbreak{}\hskip.1em \RB{43.097}{5.312}{2.605}{%
\RG{0.000}{0.001}{10.9091}{mi}{101}\RG{4.950}{-1.636}{7.9701}{rm}{49}\RG{9.268}{0.001}{10.9091}{mi}{59}\RG{14.047}{0.001}{10.9091}{mi}{58}\RG{18.836}{0.001}{10.9091}{mi}{58}\RG{23.625}{0.001}{10.9091}{mi}{58}\RG{28.403}{0.001}{10.9091}{mi}{59}\RG{33.192}{0.001}{10.9091}{mi}{101}\RG{38.146}{-1.636}{7.9701}{mi}{110}%
}\hskip.1em\allowbreak{} 下的矩阵为

\[
\RB{100.357}{12.270}{21.440}{%
\RG{25.842}{11.416}{9.0910}{ex}{48}\RG{87.956}{11.416}{9.0910}{ex}{49}\RG{33.797}{0.003}{10.9091}{mi}{70}\RG{42.195}{0.003}{10.9091}{rm}{40}\RG{46.349}{0.003}{10.9091}{mi}{103}\RG{51.856}{0.003}{10.9091}{rm}{41}\RG{72.316}{0.003}{10.9091}{mi}{67}\RG{25.842}{-4.949}{9.0910}{ex}{64}\RG{87.956}{-4.949}{9.0910}{ex}{65}\RG{0.000}{-7.312}{10.9091}{mi}{77}\RG{14.503}{-7.312}{10.9091}{rm}{61}\RG{97.729}{-7.312}{10.9091}{mi}{59}\RG{42.233}{-15.288}{10.9091}{rm}{48}\RG{65.100}{-15.288}{10.9091}{mi}{70}\RG{73.499}{-15.288}{10.9091}{rm}{40}\RG{77.653}{-15.288}{10.9091}{mi}{104}\RG{83.802}{-15.288}{10.9091}{rm}{41}%
}
\]
其中\allowbreak{}\hskip.1em \RB{47.338}{8.682}{3.227}{%
\RG{0.000}{0.000}{10.9091}{mi}{103}\RG{5.498}{0.000}{10.9091}{mi}{59}\RG{10.287}{0.000}{10.9091}{mi}{104}\RG{19.471}{0.000}{10.9091}{sy}{50}\RG{29.774}{0.000}{10.9091}{bb}{75}\RG{38.259}{0.000}{10.9091}{rm}{91}\RG{41.289}{0.000}{10.9091}{mi}{116}\RG{45.147}{0.000}{10.9091}{rm}{93}%
}\hskip.1em\allowbreak{} 为首一不可约多项式，\allowbreak{}\hskip.1em \RB{49.276}{8.682}{3.238}{%
\RG{0.000}{0.000}{10.9091}{mi}{70}\RG{8.399}{0.000}{10.9091}{rm}{40}\RG{12.553}{0.000}{10.9091}{mi}{103}\RG{18.060}{0.000}{10.9091}{rm}{41}\RG{22.213}{0.000}{10.9091}{mi}{59}\RG{26.992}{0.000}{10.9091}{mi}{70}\RG{35.397}{0.000}{10.9091}{rm}{40}\RG{39.551}{0.000}{10.9091}{mi}{104}\RG{45.700}{0.000}{10.9091}{rm}{41}%
}\hskip.1em\allowbreak{} 为相应的Frobenius块，即友阵的转置；\allowbreak{}\hskip.1em \RB{8.649}{8.183}{0.726}{%
\RG{0.000}{0.003}{10.9091}{mi}{67}%
}\hskip.1em\allowbreak{} 仅左下角元素为\allowbreak{}\hskip.1em \RB{4.973}{7.757}{0.500}{%
\RG{0.000}{0.003}{10.9091}{rm}{49}%
}\hskip.1em\allowbreak{} ，其余元素均为\allowbreak{}\hskip.1em \RB{5.387}{7.757}{0.726}{%
\RG{0.000}{0.003}{10.9091}{rm}{48}%
}\hskip.1em\allowbreak{} 。求所有\allowbreak{}\hskip.1em \RB{7.122}{5.314}{2.854}{%
\RG{0.000}{0.003}{10.9091}{mi}{39}%
}\hskip.1em\allowbreak{} -不变子空间。

(v)设

\[
\RB{114.754}{17.369}{11.438}{%
\RG{72.986}{11.822}{7.9701}{mi}{107}\RG{69.286}{9.094}{9.0910}{ex}{89}\RG{105.300}{4.508}{7.9701}{mi}{114}\RG{108.787}{3.399}{5.9776}{mi}{105}\RG{0.000}{0.004}{10.9091}{mi}{102}\RG{6.430}{0.004}{10.9091}{rm}{40}\RG{10.583}{0.004}{10.9091}{mi}{116}\RG{14.441}{0.004}{10.9091}{rm}{41}\RG{21.619}{0.004}{10.9091}{rm}{61}\RG{32.964}{0.004}{10.9091}{mi}{109}\RG{42.308}{0.004}{10.9091}{rm}{40}\RG{46.462}{0.004}{10.9091}{mi}{116}\RG{50.319}{0.004}{10.9091}{rm}{41}\RG{57.497}{0.004}{10.9091}{rm}{61}\RG{83.165}{0.004}{10.9091}{mi}{80}\RG{93.135}{0.004}{10.9091}{rm}{40}\RG{97.289}{0.004}{10.9091}{mi}{116}\RG{101.147}{0.004}{10.9091}{rm}{41}\RG{112.127}{0.004}{10.9091}{mi}{59}\RG{90.050}{-1.633}{7.9701}{mi}{105}\RG{68.842}{-10.865}{7.9701}{mi}{105}\RG{71.473}{-10.865}{7.9701}{rm}{61}\RG{77.480}{-10.865}{7.9701}{rm}{49}%
}
\]
其中\allowbreak{}\hskip.1em \RB{77.585}{8.686}{3.224}{%
\RG{0.000}{0.004}{10.9091}{mi}{80}\RG{6.885}{-1.632}{7.9701}{rm}{49}\RG{11.203}{0.004}{10.9091}{mi}{59}\RG{15.981}{0.004}{10.9091}{mi}{58}\RG{20.770}{0.004}{10.9091}{mi}{58}\RG{25.560}{0.004}{10.9091}{mi}{58}\RG{30.338}{0.004}{10.9091}{mi}{59}\RG{35.127}{0.004}{10.9091}{mi}{80}\RG{42.015}{-1.632}{7.9701}{mi}{107}\RG{49.717}{0.004}{10.9091}{sy}{50}\RG{60.021}{0.004}{10.9091}{bb}{75}\RG{68.506}{0.004}{10.9091}{rm}{91}\RG{71.536}{0.004}{10.9091}{mi}{116}\RG{75.394}{0.004}{10.9091}{rm}{93}%
}\hskip.1em\allowbreak{} 为互异的首一不可约多项式，\allowbreak{}\hskip.1em \RB{27.418}{7.758}{2.205}{%
\RG{0.000}{0.004}{10.9091}{mi}{114}\RG{4.813}{-1.632}{7.9701}{mi}{105}\RG{10.929}{0.004}{10.9091}{sy}{21}\RG{22.445}{0.004}{10.9091}{rm}{49}%
}\hskip.1em\allowbreak{} 。求所有\allowbreak{}\hskip.1em \RB{7.122}{5.315}{2.853}{%
\RG{0.000}{0.004}{10.9091}{mi}{39}%
}\hskip.1em\allowbreak{} -不变子空间。

证明：记\allowbreak{}\hskip.1em \RB{29.236}{7.899}{2.854}{%
\RG{0.000}{0.003}{10.9091}{mi}{84}\RG{10.775}{0.003}{10.9091}{rm}{61}\RG{22.114}{0.003}{10.9091}{mi}{39}%
}\hskip.1em\allowbreak{} 。先建立本题所需的分解与分类结论。

设两两互素的非常数多项式\allowbreak{}\hskip.1em \RB{41.410}{5.312}{2.615}{%
\RG{0.000}{0.001}{10.9091}{mi}{113}\RG{4.736}{-1.636}{7.9701}{rm}{49}\RG{9.054}{0.001}{10.9091}{mi}{59}\RG{13.832}{0.001}{10.9091}{mi}{58}\RG{18.621}{0.001}{10.9091}{mi}{58}\RG{23.411}{0.001}{10.9091}{mi}{58}\RG{28.189}{0.001}{10.9091}{mi}{59}\RG{32.978}{0.001}{10.9091}{mi}{113}\RG{37.717}{-1.636}{7.9701}{mi}{115}%
}\hskip.1em\allowbreak{} 满足\allowbreak{}\hskip.1em \RB{45.502}{8.683}{3.237}{%
\RG{0.000}{0.001}{10.9091}{mi}{72}\RG{9.693}{0.001}{10.9091}{rm}{40}\RG{13.846}{0.001}{10.9091}{mi}{84}\RG{21.591}{0.001}{10.9091}{rm}{41}\RG{28.769}{0.001}{10.9091}{rm}{61}\RG{40.115}{0.001}{10.9091}{rm}{48}%
}\hskip.1em\allowbreak{} ，其中

\[
\RB{123.884}{15.526}{11.441}{%
\RG{28.528}{11.819}{7.9701}{mi}{115}\RG{24.506}{9.092}{9.0910}{ex}{89}\RG{0.000}{0.001}{10.9091}{mi}{72}\RG{12.723}{0.001}{10.9091}{rm}{61}\RG{38.385}{0.001}{10.9091}{mi}{113}\RG{46.206}{0.001}{10.9091}{mi}{59}\RG{72.803}{0.001}{10.9091}{mi}{72}\RG{87.776}{0.001}{10.9091}{rm}{61}\RG{99.115}{0.001}{10.9091}{mi}{72}\RG{108.203}{0.001}{10.9091}{mi}{61}\RG{113.548}{0.001}{10.9091}{mi}{113}\RG{121.377}{0.001}{10.9091}{mi}{58}\RG{43.121}{-1.635}{7.9701}{mi}{105}\RG{81.661}{-1.635}{7.9701}{mi}{105}\RG{118.292}{-1.635}{7.9701}{mi}{105}\RG{24.062}{-10.868}{7.9701}{mi}{105}\RG{26.693}{-10.868}{7.9701}{rm}{61}\RG{32.700}{-10.868}{7.9701}{rm}{49}%
}
\]
由多项式Euclid算法，存在\allowbreak{}\hskip.1em \RB{52.002}{8.682}{3.227}{%
\RG{0.000}{-0.000}{10.9091}{mi}{97}\RG{5.607}{-1.636}{7.9701}{mi}{105}\RG{8.694}{-0.000}{10.9091}{mi}{59}\RG{13.472}{-0.000}{10.9091}{mi}{98}\RG{18.020}{-1.636}{7.9701}{mi}{105}\RG{24.136}{-0.000}{10.9091}{sy}{50}\RG{34.439}{-0.000}{10.9091}{bb}{75}\RG{42.924}{-0.000}{10.9091}{rm}{91}\RG{45.954}{-0.000}{10.9091}{mi}{116}\RG{49.811}{-0.000}{10.9091}{rm}{93}%
}\hskip.1em\allowbreak{} ，使

\[
\RB{71.451}{8.069}{2.618}{%
\RG{0.000}{-0.003}{10.9091}{mi}{97}\RG{8.693}{-0.002}{10.9091}{mi}{72}\RG{23.053}{-0.002}{10.9091}{rm}{43}\RG{33.786}{-0.002}{10.9091}{mi}{98}\RG{41.413}{-0.002}{10.9091}{mi}{113}\RG{52.264}{-0.002}{10.9091}{rm}{61}\RG{63.599}{-0.002}{10.9091}{rm}{49}\RG{68.944}{-0.002}{10.9091}{mi}{58}\RG{5.607}{-1.639}{7.9701}{mi}{105}\RG{17.543}{-1.639}{7.9701}{mi}{105}\RG{38.327}{-1.639}{7.9701}{mi}{105}\RG{46.149}{-1.639}{7.9701}{mi}{105}%
}
\]
% Original page 2.
\noindent \RB{429.891}{9.226}{18.913}{%
\RT{0.000}{-0.002}{10.9091}{令}\RG{14.590}{-0.002}{10.9091}{mi}{69}\RG{28.824}{-0.002}{10.9091}{rm}{61}\RG{40.375}{-0.002}{10.9091}{mi}{97}\RG{49.068}{-0.002}{10.9091}{rm}{40}\RG{53.221}{-0.002}{10.9091}{mi}{84}\RG{60.966}{-0.002}{10.9091}{rm}{41}\RG{65.121}{-0.002}{10.9091}{mi}{72}\RG{77.056}{-0.002}{10.9091}{rm}{40}\RG{81.210}{-0.002}{10.9091}{mi}{84}\RG{88.954}{-0.002}{10.9091}{rm}{41}\RT{93.108}{-0.002}{10.9091}{。}\RT{104.094}{-0.002}{10.9091}{由}\RT{115.090}{-0.002}{10.9091}{于}\RG{129.687}{-0.002}{10.9091}{mi}{113}\RG{137.509}{-0.002}{10.9091}{rm}{40}\RG{141.663}{-0.002}{10.9091}{mi}{84}\RG{149.408}{-0.002}{10.9091}{rm}{41}\RG{153.561}{-0.002}{10.9091}{mi}{69}\RG{167.796}{-0.002}{10.9091}{rm}{61}\RG{179.349}{-0.002}{10.9091}{rm}{48}\RT{184.687}{-0.002}{10.9091}{，}\RT{195.673}{-0.002}{10.9091}{有}\RG{210.272}{-0.002}{10.9091}{rm}{73}\RG{214.123}{-0.002}{10.9091}{rm}{109}\RG{224.898}{-0.002}{10.9091}{mi}{69}\RG{239.132}{-0.002}{10.9091}{sy}{18}\RG{250.860}{-0.002}{10.9091}{rm}{107}\RG{256.205}{-0.002}{10.9091}{rm}{101}\RG{260.951}{-0.002}{10.9091}{rm}{114}\RG{266.920}{-0.002}{10.9091}{mi}{113}\RG{274.742}{-0.002}{10.9091}{rm}{40}\RG{278.896}{-0.002}{10.9091}{mi}{84}\RG{286.640}{-0.002}{10.9091}{rm}{41}\RT{290.794}{-0.002}{10.9091}{；}\RT{301.780}{-0.002}{10.9091}{在}\RG{316.378}{-0.002}{10.9091}{rm}{107}\RG{321.724}{-0.002}{10.9091}{rm}{101}\RG{326.469}{-0.002}{10.9091}{rm}{114}\RG{332.439}{-0.002}{10.9091}{mi}{113}\RG{340.260}{-0.002}{10.9091}{rm}{40}\RG{344.415}{-0.002}{10.9091}{mi}{84}\RG{352.159}{-0.002}{10.9091}{rm}{41}\RT{359.994}{-0.002}{10.9091}{上}\RT{370.903}{-0.002}{10.9091}{，}\RG{381.813}{-0.002}{10.9091}{mi}{69}\RT{396.486}{-0.002}{10.9091}{为}\RT{407.472}{-0.002}{10.9091}{恒}\RT{418.468}{-0.002}{10.9091}{等}\RG{22.497}{-1.638}{7.9701}{mi}{105}\RG{45.982}{-1.638}{7.9701}{mi}{105}\RG{73.970}{-1.638}{7.9701}{mi}{105}\RG{134.424}{-1.638}{7.9701}{mi}{105}\RG{161.469}{-1.638}{7.9701}{mi}{105}\RG{232.805}{-1.638}{7.9701}{mi}{105}\RG{271.656}{-1.638}{7.9701}{mi}{105}\RG{337.175}{-1.638}{7.9701}{mi}{105}\RG{389.720}{-1.638}{7.9701}{mi}{105}\RG{211.609}{-8.020}{9.0910}{ex}{80}\RT{0.000}{-15.293}{10.9091}{算}\RT{10.909}{-15.293}{10.9091}{子}\RT{21.819}{-15.293}{10.9091}{，}\RT{31.942}{-15.293}{10.9091}{而}\RG{46.229}{-15.293}{10.9091}{mi}{69}\RT{61.520}{-15.293}{10.9091}{在}\RT{72.429}{-15.293}{10.9091}{该}\RT{83.338}{-15.293}{10.9091}{空}\RT{94.248}{-15.293}{10.9091}{间}\RT{105.157}{-15.293}{10.9091}{上}\RT{116.066}{-15.293}{10.9091}{为}\RT{126.975}{-15.293}{10.9091}{零}\RT{137.099}{-15.293}{10.9091}{（}\RG{148.001}{-15.293}{10.9091}{mi}{106}\RG{156.064}{-15.293}{10.9091}{sy}{54}\RG{156.064}{-15.293}{10.9091}{rm}{61}\RG{167.403}{-15.293}{10.9091}{mi}{105}\RT{171.047}{-15.293}{10.9091}{）}\RT{176.294}{-15.293}{10.9091}{。}\RT{186.407}{-15.293}{10.9091}{又}\RT{197.316}{-15.293}{10.9091}{因}\RG{226.109}{-15.293}{10.9091}{mi}{97}\RG{234.802}{-15.293}{10.9091}{mi}{72}\RG{248.810}{-15.293}{10.9091}{sy}{0}\RG{259.367}{-15.293}{10.9091}{rm}{49}\RT{268.093}{-15.293}{10.9091}{被}\RT{279.002}{-15.293}{10.9091}{每}\RT{289.911}{-15.293}{10.9091}{个}\RG{304.204}{-15.293}{10.9091}{mi}{113}\RT{316.324}{-15.293}{10.9091}{整}\RT{327.233}{-15.293}{10.9091}{除}\RT{338.142}{-15.293}{10.9091}{，}\RT{348.266}{-15.293}{10.9091}{也}\RT{359.175}{-15.293}{10.9091}{被}\RG{373.462}{-15.293}{10.9091}{mi}{72}\RT{386.540}{-15.293}{10.9091}{整}\RT{397.449}{-15.293}{10.9091}{除}\RT{408.358}{-15.293}{10.9091}{，}\RT{418.482}{-15.293}{10.9091}{故}\RG{54.137}{-16.929}{7.9701}{mi}{106}\RG{231.716}{-16.929}{7.9701}{mi}{105}\RG{243.653}{-16.929}{7.9701}{mi}{105}\RG{308.940}{-16.929}{7.9701}{mi}{106}\RG{221.206}{-17.566}{7.9701}{mi}{105}%
}

\allowbreak{}\hskip.1em \RB{45.620}{7.955}{2.842}{%
\RG{0.000}{7.276}{9.0910}{ex}{80}\RG{9.596}{-2.269}{7.9701}{mi}{105}\RG{14.500}{0.004}{10.9091}{mi}{69}\RG{22.407}{-1.633}{7.9701}{mi}{105}\RG{28.523}{0.004}{10.9091}{rm}{61}\RG{39.861}{0.004}{10.9091}{mi}{73}%
}\hskip.1em\allowbreak{} 。于是

\[
\RB{349.834}{15.528}{11.440}{%
\RG{28.052}{11.821}{7.9701}{mi}{115}\RG{189.322}{11.821}{7.9701}{mi}{115}\RG{197.977}{10.548}{9.0910}{ex}{16}\RG{263.857}{10.548}{9.0910}{ex}{17}\RG{22.970}{9.093}{9.0910}{ex}{77}\RG{184.239}{9.093}{9.0910}{ex}{77}\RG{0.000}{0.002}{10.9091}{mi}{86}\RG{11.631}{0.002}{10.9091}{rm}{61}\RG{38.525}{0.002}{10.9091}{rm}{107}\RG{43.871}{0.002}{10.9091}{rm}{101}\RG{48.616}{0.002}{10.9091}{rm}{114}\RG{54.587}{0.002}{10.9091}{mi}{113}\RG{62.408}{0.002}{10.9091}{rm}{40}\RG{66.561}{0.002}{10.9091}{mi}{84}\RG{74.306}{0.002}{10.9091}{rm}{41}\RG{78.460}{0.002}{10.9091}{mi}{59}\RG{161.476}{0.002}{10.9091}{mi}{85}\RG{172.901}{0.002}{10.9091}{rm}{61}\RG{203.406}{0.002}{10.9091}{mi}{85}\RG{214.225}{0.002}{10.9091}{sy}{92}\RG{223.922}{0.002}{10.9091}{rm}{107}\RG{229.268}{0.002}{10.9091}{rm}{101}\RG{234.013}{0.002}{10.9091}{rm}{114}\RG{239.983}{0.002}{10.9091}{mi}{113}\RG{247.804}{0.002}{10.9091}{rm}{40}\RG{251.958}{0.002}{10.9091}{mi}{84}\RG{259.703}{0.002}{10.9091}{rm}{41}\RT{335.676}{0.002}{10.9091}{(}\RT{339.832}{0.002}{10.9091}{1}\RT{345.178}{0.002}{10.9091}{)}\RG{59.322}{-1.634}{7.9701}{mi}{105}\RG{244.719}{-1.634}{7.9701}{mi}{105}\RG{23.587}{-10.867}{7.9701}{mi}{105}\RG{26.217}{-10.867}{7.9701}{rm}{61}\RG{32.225}{-10.867}{7.9701}{rm}{49}\RG{184.856}{-10.867}{7.9701}{mi}{105}\RG{187.487}{-10.867}{7.9701}{rm}{61}\RG{193.494}{-10.867}{7.9701}{rm}{49}%
}
\]
对每个\allowbreak{}\hskip.1em \RB{8.049}{7.901}{0.500}{%
\RG{0.000}{0.005}{10.9091}{mi}{84}%
}\hskip.1em\allowbreak{} -不变子空间\allowbreak{}\hskip.1em \RB{8.562}{7.956}{0.725}{%
\RG{0.000}{0.005}{10.9091}{mi}{85}%
}\hskip.1em\allowbreak{} 成立。这里第二式利用了\allowbreak{}\hskip.1em \RB{10.712}{7.934}{2.204}{%
\RG{0.000}{0.005}{10.9091}{mi}{69}\RG{7.907}{-1.632}{7.9701}{mi}{105}%
}\hskip.1em\allowbreak{} 是\allowbreak{}\hskip.1em \RB{8.049}{7.901}{0.500}{%
\RG{0.000}{0.005}{10.9091}{mi}{84}%
}\hskip.1em\allowbreak{} 的多项式，因而\allowbreak{}\hskip.1em \RB{42.997}{7.956}{2.204}{%
\RG{0.000}{0.005}{10.9091}{mi}{69}\RG{7.908}{-1.632}{7.9701}{mi}{105}\RG{10.993}{0.005}{10.9091}{mi}{85}\RG{22.669}{0.005}{10.9091}{sy}{18}\RG{34.435}{0.005}{10.9091}{mi}{85}%
}\hskip.1em\allowbreak{} 。若某个\allowbreak{}\hskip.1em \RB{7.541}{5.309}{2.618}{%
\RG{0.000}{-0.002}{10.9091}{mi}{113}\RG{4.736}{-1.639}{7.9701}{mi}{105}%
}\hskip.1em\allowbreak{} 为非零常数，其核为零，删除该因子即可。

再设\allowbreak{}\hskip.1em \RB{48.252}{8.066}{2.741}{%
\RG{0.000}{-0.005}{10.9091}{rm}{100}\RG{5.935}{-0.005}{10.9091}{rm}{101}\RG{10.680}{-0.005}{10.9091}{rm}{103}\RG{17.986}{-0.005}{10.9091}{mi}{109}\RG{30.359}{-0.005}{10.9091}{rm}{61}\RG{41.697}{-0.005}{10.9091}{mi}{110}%
}\hskip.1em\allowbreak{} ，将\allowbreak{}\hskip.1em \RB{9.522}{5.306}{0.614}{%
\RG{0.000}{-0.005}{10.9091}{mi}{109}%
}\hskip.1em\allowbreak{} 分解为

\[
\RB{53.467}{15.526}{11.442}{%
\RG{28.177}{11.819}{7.9701}{mi}{115}\RG{24.156}{9.091}{9.0910}{ex}{89}\RG{43.396}{4.866}{7.9701}{mi}{97}\RG{47.501}{3.758}{5.9776}{mi}{105}\RG{0.000}{0.000}{10.9091}{mi}{109}\RG{12.374}{0.000}{10.9091}{rm}{61}\RG{38.034}{0.000}{10.9091}{mi}{112}\RG{50.840}{-0.000}{10.9091}{mi}{59}\RG{43.396}{-2.505}{7.9701}{mi}{105}\RG{23.712}{-10.869}{7.9701}{mi}{105}\RG{26.342}{-10.869}{7.9701}{rm}{61}\RG{32.350}{-10.869}{7.9701}{rm}{49}%
}
\]
其中\allowbreak{}\hskip.1em \RB{8.166}{5.310}{2.617}{%
\RG{0.000}{-0.001}{10.9091}{mi}{112}\RG{5.361}{-1.637}{7.9701}{mi}{105}%
}\hskip.1em\allowbreak{} 首一、不可约且互异。由极小多项式的定义，\allowbreak{}\hskip.1em \RB{67.243}{8.681}{3.239}{%
\RG{0.000}{-0.001}{10.9091}{rm}{40}\RG{4.154}{-0.001}{10.9091}{mi}{109}\RG{13.503}{-0.001}{10.9091}{mi}{61}\RG{18.849}{-0.001}{10.9091}{mi}{112}\RG{24.200}{-1.637}{7.9701}{mi}{105}\RG{27.285}{-0.001}{10.9091}{rm}{41}\RG{31.441}{-0.001}{10.9091}{rm}{40}\RG{35.593}{-0.001}{10.9091}{mi}{84}\RG{43.337}{-0.001}{10.9091}{rm}{41}\RG{50.521}{-0.001}{10.9091}{sy}{54}\RG{50.521}{-0.001}{10.9091}{rm}{61}\RG{61.856}{-0.001}{10.9091}{rm}{48}%
}\hskip.1em\allowbreak{} ，故可取\allowbreak{}\hskip.1em \RB{31.204}{7.950}{2.210}{%
\RG{0.000}{-0.001}{10.9091}{mi}{120}\RG{6.070}{-1.637}{7.9701}{mi}{105}\RG{12.186}{-0.001}{10.9091}{sy}{50}\RG{22.489}{-0.001}{10.9091}{mi}{86}%
}\hskip.1em\allowbreak{} ，使

\[
\RB{78.865}{8.679}{3.241}{%
\RG{0.000}{-0.003}{10.9091}{rm}{40}\RG{4.154}{-0.003}{10.9091}{mi}{109}\RG{13.503}{-0.003}{10.9091}{mi}{61}\RG{18.849}{-0.003}{10.9091}{mi}{112}\RG{27.285}{-0.003}{10.9091}{rm}{41}\RG{31.441}{-0.003}{10.9091}{rm}{40}\RG{35.593}{-0.003}{10.9091}{mi}{84}\RG{43.338}{-0.003}{10.9091}{rm}{41}\RG{47.492}{-0.003}{10.9091}{mi}{120}\RG{59.677}{-0.003}{10.9091}{sy}{54}\RG{59.677}{-0.003}{10.9091}{rm}{61}\RG{71.012}{-0.003}{10.9091}{rm}{48}\RG{76.357}{-0.003}{10.9091}{mi}{58}\RG{24.200}{-1.640}{7.9701}{mi}{105}\RG{53.562}{-1.640}{7.9701}{mi}{105}%
}
\]
令\allowbreak{}\hskip.1em \RB{83.344}{8.686}{3.235}{%
\RG{0.000}{0.004}{10.9091}{mi}{118}\RG{5.164}{-1.632}{7.9701}{mi}{105}\RG{11.280}{0.004}{10.9091}{rm}{61}\RG{22.614}{0.004}{10.9091}{rm}{40}\RG{26.772}{0.004}{10.9091}{mi}{109}\RG{36.121}{0.004}{10.9091}{mi}{61}\RG{41.467}{0.004}{10.9091}{mi}{112}\RG{46.817}{4.870}{7.9701}{mi}{97}\RG{46.817}{-2.501}{7.9701}{mi}{105}\RG{50.922}{3.761}{5.9776}{mi}{105}\RG{54.262}{0.004}{10.9091}{rm}{41}\RG{58.418}{0.004}{10.9091}{rm}{40}\RG{62.570}{0.004}{10.9091}{mi}{84}\RG{70.315}{0.004}{10.9091}{rm}{41}\RG{74.468}{0.004}{10.9091}{mi}{120}\RG{80.538}{-1.632}{7.9701}{mi}{105}%
}\hskip.1em\allowbreak{} ，则

\[
\RB{233.077}{9.842}{3.236}{%
\RG{24.499}{4.506}{7.9701}{mi}{97}\RG{185.292}{4.506}{7.9701}{mi}{97}\RG{192.282}{4.506}{7.9701}{sy}{0}\RG{198.292}{4.506}{7.9701}{rm}{49}\RG{28.604}{3.397}{5.9776}{mi}{105}\RG{189.397}{3.397}{5.9776}{mi}{105}\RG{0.000}{0.002}{10.9091}{mi}{112}\RG{8.446}{0.002}{10.9091}{rm}{40}\RG{12.601}{0.002}{10.9091}{mi}{84}\RG{20.345}{0.002}{10.9091}{rm}{41}\RG{31.943}{0.002}{10.9091}{mi}{118}\RG{43.223}{0.002}{10.9091}{rm}{61}\RG{54.558}{0.002}{10.9091}{rm}{48}\RG{59.903}{0.002}{10.9091}{mi}{59}\RG{160.793}{0.002}{10.9091}{mi}{112}\RG{169.240}{0.002}{10.9091}{rm}{40}\RG{173.394}{0.002}{10.9091}{mi}{84}\RG{181.138}{0.002}{10.9091}{rm}{41}\RG{202.611}{0.002}{10.9091}{mi}{118}\RG{213.890}{0.002}{10.9091}{sy}{54}\RG{213.890}{0.002}{10.9091}{rm}{61}\RG{225.225}{0.002}{10.9091}{rm}{48}\RG{230.570}{0.002}{10.9091}{mi}{58}\RG{5.361}{-1.634}{7.9701}{mi}{105}\RG{37.107}{-1.634}{7.9701}{mi}{105}\RG{166.155}{-1.634}{7.9701}{mi}{105}\RG{207.775}{-1.634}{7.9701}{mi}{105}%
}
\]
所以\allowbreak{}\hskip.1em \RB{7.969}{5.310}{2.210}{%
\RG{0.000}{-0.001}{10.9091}{mi}{118}\RG{5.164}{-1.637}{7.9701}{mi}{105}%
}\hskip.1em\allowbreak{} 的首一零化多项式中次数最小者恰为\allowbreak{}\hskip.1em \RB{12.061}{8.573}{3.079}{%
\RG{0.000}{-0.001}{10.9091}{mi}{112}\RG{5.362}{4.866}{7.9701}{mi}{97}\RG{5.362}{-2.506}{7.9701}{mi}{105}\RG{9.466}{3.757}{5.9776}{mi}{105}%
}\hskip.1em\allowbreak{} 。置\allowbreak{}\hskip.1em \RB{42.394}{7.773}{2.846}{%
\RG{0.000}{-0.001}{10.9091}{mi}{118}\RG{8.586}{-0.001}{10.9091}{rm}{61}\RG{19.924}{7.273}{9.0910}{ex}{80}\RG{29.521}{-2.273}{7.9701}{mi}{105}\RG{34.424}{-0.001}{10.9091}{mi}{118}\RG{39.588}{-1.637}{7.9701}{mi}{105}%
}\hskip.1em\allowbreak{} 。若\allowbreak{}\hskip.1em \RB{46.488}{8.681}{3.239}{%
\RG{0.000}{-0.001}{10.9091}{mi}{113}\RG{5.127}{-0.001}{10.9091}{rm}{40}\RG{9.282}{-0.001}{10.9091}{mi}{84}\RG{17.027}{-0.001}{10.9091}{rm}{41}\RG{21.180}{-0.001}{10.9091}{mi}{118}\RG{29.766}{-0.001}{10.9091}{rm}{61}\RG{41.101}{-0.001}{10.9091}{rm}{48}%
}\hskip.1em\allowbreak{} ，将式(1)中对应于\allowbreak{}\hskip.1em \RB{34.251}{8.576}{3.076}{%
\RG{0.000}{0.002}{10.9091}{mi}{113}\RG{4.736}{-1.634}{7.9701}{mi}{105}\RG{10.852}{0.002}{10.9091}{rm}{61}\RG{22.190}{0.002}{10.9091}{mi}{112}\RG{27.551}{4.868}{7.9701}{mi}{97}\RG{27.551}{-2.503}{7.9701}{mi}{105}\RG{31.656}{3.760}{5.9776}{mi}{105}%
}\hskip.1em\allowbreak{} 的投影作用于该式，得\allowbreak{}\hskip.1em \RB{49.182}{8.684}{3.236}{%
\RG{0.000}{0.002}{10.9091}{mi}{113}\RG{5.127}{0.002}{10.9091}{rm}{40}\RG{9.281}{0.002}{10.9091}{mi}{84}\RG{17.026}{0.002}{10.9091}{rm}{41}\RG{21.179}{0.002}{10.9091}{mi}{118}\RG{26.344}{-1.634}{7.9701}{mi}{105}\RG{32.460}{0.002}{10.9091}{rm}{61}\RG{43.795}{0.002}{10.9091}{rm}{48}%
}\hskip.1em\allowbreak{} ，故每个\allowbreak{}\hskip.1em \RB{12.061}{8.576}{3.076}{%
\RG{0.000}{0.002}{10.9091}{mi}{112}\RG{5.362}{4.868}{7.9701}{mi}{97}\RG{5.362}{-2.503}{7.9701}{mi}{105}\RG{9.466}{3.760}{5.9776}{mi}{105}%
}\hskip.1em\allowbreak{} 均整除\allowbreak{}\hskip.1em \RB{5.311}{5.313}{2.615}{%
\RG{0.000}{0.002}{10.9091}{mi}{113}%
}\hskip.1em\allowbreak{} ，从而\allowbreak{}\hskip.1em \RB{23.746}{8.684}{3.225}{%
\RG{0.000}{0.002}{10.9091}{mi}{109}\RG{12.374}{0.002}{10.9091}{sy}{106}\RG{18.435}{0.002}{10.9091}{mi}{113}%
}\hskip.1em\allowbreak{} 。因此\allowbreak{}\hskip.1em \RB{5.486}{5.306}{0.614}{%
\RG{0.000}{-0.005}{10.9091}{mi}{118}%
}\hskip.1em\allowbreak{} 的极小多项式为\allowbreak{}\hskip.1em \RB{9.522}{5.306}{0.614}{%
\RG{0.000}{-0.005}{10.9091}{mi}{109}%
}\hskip.1em\allowbreak{} ，且

\[
\RB{75.815}{9.839}{2.607}{%
\RG{55.312}{4.502}{7.9701}{mi}{110}\RG{60.000}{4.502}{7.9701}{sy}{0}\RG{66.010}{4.502}{7.9701}{rm}{49}\RG{0.000}{-0.002}{10.9091}{mi}{118}\RG{5.542}{-0.002}{10.9091}{mi}{59}\RG{10.331}{-0.002}{10.9091}{mi}{84}\RG{18.077}{-0.002}{10.9091}{mi}{118}\RG{23.629}{-0.002}{10.9091}{mi}{59}\RG{28.419}{-0.002}{10.9091}{mi}{58}\RG{33.197}{-0.002}{10.9091}{mi}{58}\RG{37.986}{-0.002}{10.9091}{mi}{58}\RG{42.775}{-0.002}{10.9091}{mi}{59}\RG{47.554}{-0.002}{10.9091}{mi}{84}\RG{70.329}{-0.002}{10.9091}{mi}{118}%
}
\]
线性无关，构成\allowbreak{}\hskip.1em \RB{8.715}{7.947}{0.734}{%
\RG{0.000}{-0.005}{10.9091}{mi}{86}%
}\hskip.1em\allowbreak{} 的一组基。

对任一\allowbreak{}\hskip.1em \RB{8.049}{7.899}{0.500}{%
\RG{0.000}{0.003}{10.9091}{mi}{84}%
}\hskip.1em\allowbreak{} -不变子空间\allowbreak{}\hskip.1em \RB{8.562}{7.954}{0.726}{%
\RG{0.000}{0.003}{10.9091}{mi}{85}%
}\hskip.1em\allowbreak{} ，集合

\[
\RB{131.634}{8.682}{3.238}{%
\RG{0.000}{-0.000}{10.9091}{mi}{74}\RG{15.391}{-0.000}{10.9091}{rm}{61}\RG{26.729}{-0.000}{10.9091}{sy}{102}\RG{32.184}{-0.000}{10.9091}{mi}{113}\RG{40.342}{-0.000}{10.9091}{sy}{50}\RG{50.645}{-0.000}{10.9091}{bb}{75}\RG{59.131}{-0.000}{10.9091}{rm}{91}\RG{62.160}{-0.000}{10.9091}{mi}{116}\RG{66.018}{-0.000}{10.9091}{rm}{93}\RG{72.008}{-0.000}{10.9091}{rm}{58}\RG{78.013}{-0.000}{10.9091}{mi}{113}\RG{83.141}{-0.000}{10.9091}{rm}{40}\RG{87.295}{-0.000}{10.9091}{mi}{84}\RG{95.039}{-0.000}{10.9091}{rm}{41}\RG{99.193}{-0.000}{10.9091}{mi}{118}\RG{107.779}{-0.000}{10.9091}{sy}{50}\RG{118.082}{-0.000}{10.9091}{mi}{85}\RG{126.476}{-0.000}{10.9091}{sy}{103}\RG{5.904}{-1.636}{7.9701}{mi}{85}%
}
\]
是包含\allowbreak{}\hskip.1em \RB{9.522}{5.309}{0.611}{%
\RG{0.000}{-0.002}{10.9091}{mi}{109}%
}\hskip.1em\allowbreak{} 的理想。取其中次数最小的首一非零多项式\allowbreak{}\hskip.1em \RB{6.020}{8.069}{0.611}{%
\RG{0.000}{-0.002}{10.9091}{mi}{100}%
}\hskip.1em\allowbreak{} ，用带余除法可得\allowbreak{}\hskip.1em \RB{40.011}{8.680}{3.240}{%
\RG{0.000}{-0.002}{10.9091}{mi}{74}\RG{5.905}{-1.638}{7.9701}{mi}{85}\RG{15.391}{-0.002}{10.9091}{rm}{61}\RG{26.726}{-0.002}{10.9091}{rm}{40}\RG{30.884}{-0.002}{10.9091}{mi}{100}\RG{36.435}{-0.002}{10.9091}{rm}{41}%
}\hskip.1em\allowbreak{} ，并且\allowbreak{}\hskip.1em \RB{24.163}{8.683}{3.226}{%
\RG{0.000}{0.001}{10.9091}{mi}{100}\RG{8.581}{0.001}{10.9091}{sy}{106}\RG{14.641}{0.001}{10.9091}{mi}{109}%
}\hskip.1em\allowbreak{} 。于是\allowbreak{}\hskip.1em \RB{53.080}{8.683}{3.237}{%
\RG{0.000}{0.001}{10.9091}{mi}{85}\RG{11.424}{0.001}{10.9091}{rm}{61}\RG{22.763}{0.001}{10.9091}{mi}{100}\RG{28.313}{0.001}{10.9091}{rm}{40}\RG{32.467}{0.001}{10.9091}{mi}{84}\RG{40.212}{0.001}{10.9091}{rm}{41}\RG{44.365}{0.001}{10.9091}{mi}{86}%
}\hskip.1em\allowbreak{} 。写\allowbreak{}\hskip.1em \RB{35.601}{8.072}{0.608}{%
\RG{0.000}{0.001}{10.9091}{mi}{109}\RG{12.374}{0.001}{10.9091}{rm}{61}\RG{23.712}{0.001}{10.9091}{mi}{100}\RG{29.264}{0.001}{10.9091}{mi}{104}%
}\hskip.1em\allowbreak{} ，则

\[
\RB{194.670}{8.680}{3.240}{%
\RG{0.000}{-0.002}{10.9091}{mi}{113}\RG{5.128}{-0.002}{10.9091}{rm}{40}\RG{9.282}{-0.002}{10.9091}{mi}{84}\RG{17.027}{-0.002}{10.9091}{rm}{41}\RG{21.180}{-0.002}{10.9091}{mi}{118}\RG{29.766}{-0.002}{10.9091}{sy}{50}\RG{40.069}{-0.002}{10.9091}{rm}{107}\RG{45.415}{-0.002}{10.9091}{rm}{101}\RG{50.160}{-0.002}{10.9091}{rm}{114}\RG{56.130}{-0.002}{10.9091}{mi}{104}\RG{62.279}{-0.002}{10.9091}{rm}{40}\RG{66.433}{-0.002}{10.9091}{mi}{84}\RG{74.177}{-0.002}{10.9091}{rm}{41}\RG{84.392}{-0.002}{10.9091}{sy}{40}\RG{93.490}{-0.002}{10.9091}{sy}{41}\RG{110.452}{-0.002}{10.9091}{mi}{109}\RG{122.826}{-0.002}{10.9091}{sy}{106}\RG{128.887}{-0.002}{10.9091}{mi}{104}\RG{135.040}{-0.002}{10.9091}{mi}{113}\RG{146.225}{-0.002}{10.9091}{sy}{40}\RG{155.323}{-0.002}{10.9091}{sy}{41}\RG{172.285}{-0.002}{10.9091}{mi}{100}\RG{180.866}{-0.002}{10.9091}{sy}{106}\RG{186.926}{-0.002}{10.9091}{mi}{113}\RG{192.043}{-0.002}{10.9091}{mi}{59}%
}
\]
故\allowbreak{}\hskip.1em \RB{60.448}{8.678}{3.242}{%
\RG{0.000}{-0.004}{10.9091}{mi}{85}\RG{11.425}{-0.004}{10.9091}{rm}{61}\RG{22.759}{-0.004}{10.9091}{rm}{107}\RG{28.116}{-0.004}{10.9091}{rm}{101}\RG{32.861}{-0.004}{10.9091}{rm}{114}\RG{38.824}{-0.004}{10.9091}{mi}{104}\RG{44.973}{-0.004}{10.9091}{rm}{40}\RG{49.127}{-0.004}{10.9091}{mi}{84}\RG{56.871}{-0.004}{10.9091}{rm}{41}%
}\hskip.1em\allowbreak{} 。映射\allowbreak{}\hskip.1em \RB{70.365}{8.678}{3.242}{%
\RG{0.000}{-0.004}{10.9091}{mi}{113}\RG{8.158}{-0.004}{10.9091}{sy}{55}\RG{8.158}{-0.004}{10.9091}{sy}{33}\RG{22.097}{-0.004}{10.9091}{mi}{113}\RG{27.225}{-0.004}{10.9091}{rm}{40}\RG{31.379}{-0.004}{10.9091}{mi}{84}\RG{39.123}{-0.004}{10.9091}{rm}{41}\RG{43.277}{-0.004}{10.9091}{mi}{100}\RG{48.828}{-0.004}{10.9091}{rm}{40}\RG{52.981}{-0.004}{10.9091}{mi}{84}\RG{60.726}{-0.004}{10.9091}{rm}{41}\RG{64.880}{-0.004}{10.9091}{mi}{118}%
}\hskip.1em\allowbreak{} 的核为理想\allowbreak{}\hskip.1em \RB{13.880}{8.678}{3.242}{%
\RG{0.000}{-0.004}{10.9091}{rm}{40}\RG{4.154}{-0.004}{10.9091}{mi}{104}\RG{10.303}{-0.004}{10.9091}{rm}{41}%
}\hskip.1em\allowbreak{} ，因此其像的维数为\allowbreak{}\hskip.1em \RB{24.321}{8.067}{2.740}{%
\RG{0.000}{-0.004}{10.9091}{rm}{100}\RG{5.935}{-0.004}{10.9091}{rm}{101}\RG{10.680}{-0.004}{10.9091}{rm}{103}\RG{17.985}{-0.004}{10.9091}{mi}{104}%
}\hskip.1em\allowbreak{} 。这证明：若\allowbreak{}\hskip.1em \RB{30.168}{8.179}{2.726}{%
\RG{0.000}{-0.001}{10.9091}{mi}{109}\RG{12.374}{-0.001}{10.9091}{rm}{61}\RG{23.712}{-0.001}{10.9091}{mi}{102}%
}\hskip.1em\allowbreak{} ，则所有不变子空间恰为

\allowbreak{}\hskip.1em \RB{247.689}{8.678}{3.242}{%
\RG{0.000}{-0.004}{10.9091}{mi}{85}\RG{7.254}{-1.640}{7.9701}{mi}{104}\RG{15.192}{-0.004}{10.9091}{rm}{61}\RG{26.527}{-0.004}{10.9091}{rm}{107}\RG{31.884}{-0.004}{10.9091}{rm}{101}\RG{36.629}{-0.004}{10.9091}{rm}{114}\RG{42.592}{-0.004}{10.9091}{mi}{104}\RG{48.741}{-0.004}{10.9091}{rm}{40}\RG{52.895}{-0.004}{10.9091}{mi}{84}\RG{60.639}{-0.004}{10.9091}{rm}{41}\RG{67.817}{-0.004}{10.9091}{rm}{61}\RG{79.163}{-0.004}{10.9091}{rm}{40}\RG{83.316}{-0.004}{10.9091}{mi}{109}\RG{92.665}{-0.004}{10.9091}{mi}{61}\RG{98.011}{-0.004}{10.9091}{mi}{104}\RG{104.150}{-0.004}{10.9091}{rm}{41}\RG{108.306}{-0.004}{10.9091}{rm}{40}\RG{112.457}{-0.004}{10.9091}{mi}{84}\RG{120.202}{-0.004}{10.9091}{rm}{41}\RG{124.355}{-0.004}{10.9091}{mi}{86}\RG{131.174}{-0.004}{10.9091}{mi}{59}\RG{201.076}{-0.004}{10.9091}{mi}{104}\RG{210.256}{-0.004}{10.9091}{sy}{106}\RG{216.316}{-0.004}{10.9091}{mi}{109}\RG{225.665}{-0.004}{10.9091}{mi}{59}\RG{241.352}{-0.004}{10.9091}{mi}{104}%
}\hskip.1em\allowbreak{} 首一\allowbreak{}\hskip.1em \RB{2.627}{1.555}{2.609}{%
\RG{0.000}{-0.004}{10.9091}{mi}{59}%
}\hskip.1em\allowbreak{}

\[
\RB{362.951}{9.296}{2.734}{%
\RG{0.000}{0.003}{10.9091}{rm}{100}\RG{5.935}{0.003}{10.9091}{rm}{105}\RG{8.902}{0.003}{10.9091}{rm}{109}\RG{19.674}{0.003}{10.9091}{mi}{85}\RG{34.866}{0.003}{10.9091}{rm}{61}\RG{46.201}{0.003}{10.9091}{rm}{100}\RG{52.136}{0.003}{10.9091}{rm}{101}\RG{56.881}{0.003}{10.9091}{rm}{103}\RG{64.190}{0.003}{10.9091}{mi}{104}\RG{70.343}{0.003}{10.9091}{mi}{58}\RT{348.792}{0.003}{10.9091}{(}\RT{352.949}{0.003}{10.9091}{2}\RT{358.295}{0.003}{10.9091}{)}\RG{26.928}{-1.633}{7.9701}{mi}{104}%
}
\]
不同的\allowbreak{}\hskip.1em \RB{6.336}{8.071}{0.609}{%
\RG{0.000}{-0.000}{10.9091}{mi}{104}%
}\hskip.1em\allowbreak{} 对应不同的子空间，因为\allowbreak{}\hskip.1em \RB{58.932}{8.682}{3.474}{%
\RG{0.000}{-0.000}{10.9091}{mi}{74}\RG{5.904}{-1.636}{7.9701}{mi}{85}\RG{11.132}{-2.919}{5.9776}{mi}{104}\RG{19.030}{-0.000}{10.9091}{rm}{61}\RG{30.364}{-0.000}{10.9091}{rm}{40}\RG{34.522}{-0.000}{10.9091}{mi}{109}\RG{43.871}{-0.000}{10.9091}{mi}{61}\RG{49.217}{-0.000}{10.9091}{mi}{104}\RG{55.356}{-0.000}{10.9091}{rm}{41}%
}\hskip.1em\allowbreak{} 唯一确定。

% Original page 3.
(i)令\allowbreak{}\hskip.1em \RB{61.448}{8.069}{2.291}{%
\RG{0.000}{-0.002}{10.9091}{mi}{78}\RG{12.723}{-0.002}{10.9091}{rm}{61}\RG{24.062}{-0.002}{10.9091}{mi}{84}\RG{34.231}{-0.002}{10.9091}{sy}{0}\RG{45.140}{-0.002}{10.9091}{mi}{21}\RG{51.371}{-1.638}{7.9701}{rm}{48}\RG{55.690}{-0.002}{10.9091}{mi}{73}%
}\hskip.1em\allowbreak{} 。由式(2)，所求子空间为

\[
\RB{158.097}{9.817}{2.727}{%
\RG{25.754}{4.502}{7.9701}{mi}{106}\RG{68.440}{4.502}{7.9701}{mi}{110}\RG{73.130}{4.502}{7.9701}{sy}{0}\RG{79.139}{4.502}{7.9701}{mi}{106}\RG{0.000}{-0.002}{10.9091}{rm}{107}\RG{5.346}{-0.002}{10.9091}{rm}{101}\RG{10.091}{-0.002}{10.9091}{rm}{114}\RG{16.060}{-0.002}{10.9091}{mi}{78}\RG{32.783}{-0.002}{10.9091}{rm}{61}\RG{44.118}{-0.002}{10.9091}{rm}{73}\RG{47.969}{-0.002}{10.9091}{rm}{109}\RG{58.748}{-0.002}{10.9091}{mi}{78}\RG{83.138}{-0.002}{10.9091}{mi}{59}\RG{109.742}{-0.002}{10.9091}{rm}{48}\RG{118.114}{-0.002}{10.9091}{sy}{20}\RG{129.629}{-0.002}{10.9091}{mi}{106}\RG{137.692}{-0.002}{10.9091}{sy}{20}\RG{149.208}{-0.002}{10.9091}{mi}{110}\RG{155.590}{-0.002}{10.9091}{mi}{58}%
}
\]
它们构成严格递增链，每个空间的维数为\allowbreak{}\hskip.1em \RB{4.766}{7.665}{2.728}{%
\RG{0.000}{-0.003}{10.9091}{mi}{106}%
}\hskip.1em\allowbreak{} 。其中\allowbreak{}\hskip.1em \RB{24.785}{7.752}{2.728}{%
\RG{0.000}{-0.003}{10.9091}{mi}{106}\RG{8.063}{-0.003}{10.9091}{rm}{61}\RG{19.398}{-0.003}{10.9091}{rm}{48}%
}\hskip.1em\allowbreak{} 时为\allowbreak{}\hskip.1em \RB{15.954}{8.679}{3.230}{%
\RG{0.000}{-0.003}{10.9091}{sy}{102}\RG{5.455}{-0.003}{10.9091}{rm}{48}\RG{10.796}{-0.003}{10.9091}{sy}{103}%
}\hskip.1em\allowbreak{} ，\allowbreak{}\hskip.1em \RB{25.956}{7.665}{2.728}{%
\RG{0.000}{-0.003}{10.9091}{mi}{106}\RG{8.063}{-0.003}{10.9091}{rm}{61}\RG{19.402}{-0.003}{10.9091}{mi}{110}%
}\hskip.1em\allowbreak{} 时为\allowbreak{}\hskip.1em \RB{8.715}{7.948}{0.732}{%
\RG{0.000}{-0.003}{10.9091}{mi}{86}%
}\hskip.1em\allowbreak{} 。(ii) Cayley–Hamilton定理给出\allowbreak{}\hskip.1em \RB{43.668}{8.680}{3.240}{%
\RG{0.000}{-0.002}{10.9091}{mi}{102}\RG{6.430}{-0.002}{10.9091}{rm}{40}\RG{10.584}{-0.002}{10.9091}{mi}{84}\RG{18.328}{-0.002}{10.9091}{rm}{41}\RG{26.226}{-0.002}{10.9091}{rm}{61}\RG{38.281}{-0.002}{10.9091}{rm}{48}%
}\hskip.1em\allowbreak{} （定理的简证见文末）。在式(1)中取\allowbreak{}\hskip.1em \RB{31.942}{8.178}{2.727}{%
\RG{0.000}{-0.002}{10.9091}{mi}{72}\RG{13.436}{-0.002}{10.9091}{rm}{61}\RG{25.486}{-0.002}{10.9091}{mi}{102}%
}\hskip.1em\allowbreak{} 、\allowbreak{}\hskip.1em \RB{30.261}{8.181}{2.725}{%
\RG{0.000}{0.001}{10.9091}{mi}{113}\RG{4.736}{-1.635}{7.9701}{mi}{105}\RG{10.853}{0.001}{10.9091}{rm}{61}\RG{22.190}{0.001}{10.9091}{mi}{102}\RG{27.456}{-1.635}{7.9701}{mi}{105}%
}\hskip.1em\allowbreak{} ，得到

\[
\RB{132.308}{17.363}{11.444}{%
\RG{27.523}{11.816}{7.9701}{mi}{107}\RG{22.763}{9.089}{9.0910}{ex}{77}\RG{0.000}{-0.003}{10.9091}{mi}{85}\RG{11.424}{-0.003}{10.9091}{rm}{61}\RG{38.318}{-0.003}{10.9091}{mi}{85}\RG{48.659}{-0.003}{10.9091}{mi}{59}\RG{75.256}{-0.003}{10.9091}{mi}{85}\RG{88.633}{-0.003}{10.9091}{rm}{61}\RG{99.971}{-0.003}{10.9091}{mi}{85}\RG{110.790}{-0.003}{10.9091}{sy}{92}\RG{120.487}{-0.003}{10.9091}{mi}{86}\RG{129.800}{-0.003}{10.9091}{mi}{58}\RG{45.574}{-1.639}{7.9701}{mi}{105}\RG{82.518}{-1.639}{7.9701}{mi}{105}\RG{126.714}{-1.639}{7.9701}{mi}{105}\RG{23.379}{-10.871}{7.9701}{mi}{105}\RG{26.010}{-10.871}{7.9701}{rm}{61}\RG{32.018}{-10.871}{7.9701}{rm}{49}%
}
\]
因为\allowbreak{}\hskip.1em \RB{8.049}{7.899}{0.500}{%
\RG{0.000}{0.003}{10.9091}{mi}{84}%
}\hskip.1em\allowbreak{} 与\allowbreak{}\hskip.1em \RB{23.826}{8.685}{3.235}{%
\RG{0.000}{0.003}{10.9091}{mi}{102}\RG{5.265}{-1.633}{7.9701}{mi}{105}\RG{8.351}{0.003}{10.9091}{rm}{40}\RG{12.505}{0.003}{10.9091}{mi}{84}\RG{20.250}{0.003}{10.9091}{rm}{41}%
}\hskip.1em\allowbreak{} 可交换，\allowbreak{}\hskip.1em \RB{9.032}{7.954}{2.206}{%
\RG{0.000}{0.003}{10.9091}{mi}{86}\RG{6.227}{-1.633}{7.9701}{mi}{105}%
}\hskip.1em\allowbreak{} 为\allowbreak{}\hskip.1em \RB{8.049}{7.899}{0.500}{%
\RG{0.000}{0.003}{10.9091}{mi}{84}%
}\hskip.1em\allowbreak{} -不变子空间；因此\allowbreak{}\hskip.1em \RB{10.060}{7.954}{2.206}{%
\RG{0.000}{0.003}{10.9091}{mi}{85}\RG{7.254}{-1.633}{7.9701}{mi}{105}%
}\hskip.1em\allowbreak{} 也为\allowbreak{}\hskip.1em \RB{8.049}{7.899}{0.500}{%
\RG{0.000}{0.003}{10.9091}{mi}{84}%
}\hskip.1em\allowbreak{} -不变子空间。上述分解同时表明\allowbreak{}\hskip.1em \RB{10.060}{7.947}{2.213}{%
\RG{0.000}{-0.004}{10.9091}{mi}{85}\RG{7.254}{-1.641}{7.9701}{mi}{105}%
}\hskip.1em\allowbreak{} 由\allowbreak{}\hskip.1em \RB{8.562}{7.947}{0.733}{%
\RG{0.000}{-0.004}{10.9091}{mi}{85}%
}\hskip.1em\allowbreak{} 唯一确定。(iii)所有不变子空间恰为

\[
\RB{205.057}{17.369}{11.438}{%
\RG{51.113}{11.822}{7.9701}{mi}{107}\RG{46.353}{9.094}{9.0910}{ex}{77}\RG{120.422}{4.508}{7.9701}{mi}{116}\RG{123.212}{3.399}{5.9776}{mi}{105}\RG{0.000}{0.004}{10.9091}{mi}{85}\RG{35.014}{0.004}{10.9091}{rm}{61}\RG{61.908}{0.004}{10.9091}{rm}{107}\RG{67.254}{0.004}{10.9091}{rm}{101}\RG{72.000}{0.004}{10.9091}{rm}{114}\RG{76.156}{0.004}{10.9091}{rm}{40}\RG{80.305}{0.004}{10.9091}{mi}{84}\RG{90.473}{0.004}{10.9091}{sy}{0}\RG{101.383}{0.004}{10.9091}{mi}{21}\RG{110.699}{0.004}{10.9091}{mi}{73}\RG{116.268}{0.004}{10.9091}{rm}{41}\RG{126.552}{0.004}{10.9091}{mi}{59}\RG{153.156}{0.004}{10.9091}{rm}{48}\RG{161.528}{0.004}{10.9091}{sy}{20}\RG{173.043}{0.004}{10.9091}{mi}{116}\RG{183.015}{0.004}{10.9091}{sy}{20}\RG{194.531}{0.004}{10.9091}{mi}{114}\RG{202.430}{0.004}{10.9091}{mi}{59}\RG{7.254}{-1.633}{7.9701}{mi}{116}\RG{13.833}{-1.633}{7.9701}{mi}{59}\RG{15.978}{-1.633}{7.9701}{mi}{58}\RG{18.124}{-1.633}{7.9701}{mi}{58}\RG{20.269}{-1.633}{7.9701}{mi}{58}\RG{22.415}{-1.633}{7.9701}{mi}{59}\RG{24.560}{-1.633}{7.9701}{mi}{116}\RG{107.614}{-1.633}{7.9701}{mi}{105}\RG{176.900}{-1.633}{7.9701}{mi}{105}\RG{199.345}{-1.633}{7.9701}{mi}{105}\RG{10.045}{-2.644}{5.9776}{rm}{49}\RG{27.356}{-2.915}{5.9776}{mi}{107}\RG{46.969}{-10.865}{7.9701}{mi}{105}\RG{49.600}{-10.865}{7.9701}{rm}{61}\RG{55.607}{-10.865}{7.9701}{rm}{49}%
}
\]
并且

\[
\RB{90.426}{19.004}{9.803}{%
\RG{70.483}{13.457}{7.9701}{mi}{107}\RG{66.027}{10.730}{9.0910}{ex}{88}\RG{0.000}{1.639}{10.9091}{rm}{100}\RG{5.935}{1.639}{10.9091}{rm}{105}\RG{8.902}{1.639}{10.9091}{rm}{109}\RG{19.674}{1.639}{10.9091}{mi}{85}\RG{54.688}{1.639}{10.9091}{rm}{61}\RG{80.976}{1.639}{10.9091}{mi}{116}\RG{87.919}{1.639}{10.9091}{mi}{58}\RG{26.928}{0.003}{7.9701}{mi}{116}\RG{33.507}{0.003}{7.9701}{mi}{59}\RG{35.652}{0.003}{7.9701}{mi}{58}\RG{37.798}{0.003}{7.9701}{mi}{58}\RG{39.943}{0.003}{7.9701}{mi}{58}\RG{42.088}{0.003}{7.9701}{mi}{59}\RG{44.234}{0.003}{7.9701}{mi}{116}\RG{84.834}{0.003}{7.9701}{mi}{105}\RG{29.719}{-1.007}{5.9776}{rm}{49}\RG{47.030}{-1.280}{5.9776}{mi}{107}\RG{66.340}{-9.230}{7.9701}{mi}{105}\RG{68.970}{-9.230}{7.9701}{rm}{61}\RG{74.978}{-9.230}{7.9701}{rm}{49}%
}
\]
事实上，式(2)中的\allowbreak{}\hskip.1em \RB{6.336}{8.069}{0.611}{%
\RG{0.000}{-0.002}{10.9091}{mi}{104}%
}\hskip.1em\allowbreak{} 恰为\allowbreak{}\hskip.1em \RB{51.140}{9.002}{3.240}{%
\RG{0.000}{7.271}{9.0910}{ex}{81}\RG{8.586}{-2.275}{7.9701}{mi}{105}\RG{11.672}{-0.002}{10.9091}{rm}{40}\RG{15.826}{-0.002}{10.9091}{mi}{116}\RG{21.742}{-0.002}{10.9091}{sy}{0}\RG{32.285}{-0.002}{10.9091}{mi}{21}\RG{38.516}{-1.639}{7.9701}{mi}{105}\RG{41.601}{-0.002}{10.9091}{rm}{41}\RG{45.755}{3.957}{7.9701}{mi}{116}\RG{48.546}{2.848}{5.9776}{mi}{105}%
}\hskip.1em\allowbreak{} 。在每个\allowbreak{}\hskip.1em \RB{63.863}{8.680}{3.240}{%
\RG{0.000}{-0.002}{10.9091}{rm}{107}\RG{5.346}{-0.002}{10.9091}{rm}{101}\RG{10.091}{-0.002}{10.9091}{rm}{114}\RG{14.248}{-0.002}{10.9091}{rm}{40}\RG{18.396}{-0.002}{10.9091}{mi}{84}\RG{28.198}{-0.002}{10.9091}{sy}{0}\RG{38.743}{-0.002}{10.9091}{mi}{21}\RG{44.973}{-1.639}{7.9701}{mi}{105}\RG{48.059}{-0.002}{10.9091}{mi}{73}\RG{53.627}{-0.002}{10.9091}{rm}{41}\RG{57.782}{3.957}{7.9701}{mi}{114}\RG{61.269}{2.848}{5.9776}{mi}{105}%
}\hskip.1em\allowbreak{} 上，其余因子均可逆：对与\allowbreak{}\hskip.1em \RB{41.193}{8.682}{3.238}{%
\RG{0.000}{0.000}{10.9091}{rm}{40}\RG{4.154}{0.000}{10.9091}{mi}{116}\RG{10.583}{0.000}{10.9091}{sy}{0}\RG{21.640}{0.000}{10.9091}{mi}{21}\RG{27.872}{-1.636}{7.9701}{mi}{105}\RG{30.957}{0.000}{10.9091}{rm}{41}\RG{35.111}{3.959}{7.9701}{mi}{114}\RG{38.598}{2.851}{5.9776}{mi}{105}%
}\hskip.1em\allowbreak{} 互素的多项式作Bézout等式即可得到逆算子。因此\allowbreak{}\hskip.1em \RB{37.684}{8.682}{3.238}{%
\RG{0.000}{0.000}{10.9091}{rm}{107}\RG{5.346}{0.000}{10.9091}{rm}{101}\RG{10.091}{0.000}{10.9091}{rm}{114}\RG{16.060}{0.000}{10.9091}{mi}{104}\RG{22.210}{0.000}{10.9091}{rm}{40}\RG{26.363}{0.000}{10.9091}{mi}{84}\RG{34.108}{0.000}{10.9091}{rm}{41}%
}\hskip.1em\allowbreak{} 正是上述直和。不同指标组对应不同的\allowbreak{}\hskip.1em \RB{6.336}{8.074}{0.606}{%
\RG{0.000}{0.003}{10.9091}{mi}{104}%
}\hskip.1em\allowbreak{} ，故共有\allowbreak{}\hskip.1em \RB{45.806}{8.685}{3.235}{%
\RG{0.000}{7.276}{9.0910}{ex}{81}\RG{8.586}{-2.269}{7.9701}{mi}{105}\RG{11.671}{0.003}{10.9091}{rm}{40}\RG{15.825}{0.003}{10.9091}{mi}{114}\RG{20.639}{-1.633}{7.9701}{mi}{105}\RG{26.149}{0.003}{10.9091}{rm}{43}\RG{36.884}{0.003}{10.9091}{rm}{49}\RG{42.229}{0.003}{10.9091}{rm}{41}%
}\hskip.1em\allowbreak{} 个不变子空间。(iv)令\allowbreak{}\hskip.1em \RB{43.310}{8.075}{2.733}{%
\RG{0.000}{0.004}{10.9091}{mi}{112}\RG{8.391}{0.004}{10.9091}{rm}{61}\RG{19.726}{0.004}{10.9091}{rm}{100}\RG{25.661}{0.004}{10.9091}{rm}{101}\RG{30.406}{0.004}{10.9091}{rm}{103}\RG{37.715}{0.004}{10.9091}{mi}{103}%
}\hskip.1em\allowbreak{} 、\allowbreak{}\hskip.1em \RB{43.818}{8.075}{2.733}{%
\RG{0.000}{0.004}{10.9091}{mi}{113}\RG{8.158}{0.004}{10.9091}{rm}{61}\RG{19.493}{0.004}{10.9091}{rm}{100}\RG{25.427}{0.004}{10.9091}{rm}{101}\RG{30.173}{0.004}{10.9091}{rm}{103}\RG{37.481}{0.004}{10.9091}{mi}{104}%
}\hskip.1em\allowbreak{} ，于是\allowbreak{}\hskip.1em \RB{44.574}{6.787}{2.613}{%
\RG{0.000}{0.004}{10.9091}{mi}{110}\RG{9.406}{0.004}{10.9091}{rm}{61}\RG{20.745}{0.004}{10.9091}{mi}{112}\RG{28.531}{0.004}{10.9091}{rm}{43}\RG{39.263}{0.004}{10.9091}{mi}{113}%
}\hskip.1em\allowbreak{} 。若

\[
\RB{140.875}{9.844}{3.544}{%
\RG{35.898}{4.508}{7.9701}{mi}{112}\RG{76.073}{4.508}{7.9701}{mi}{112}\RG{79.963}{4.508}{7.9701}{sy}{0}\RG{85.972}{4.508}{7.9701}{rm}{49}\RG{0.000}{0.004}{10.9091}{mi}{103}\RG{5.507}{0.004}{10.9091}{rm}{40}\RG{9.661}{0.004}{10.9091}{mi}{116}\RG{13.518}{0.004}{10.9091}{rm}{41}\RG{20.696}{0.004}{10.9091}{rm}{61}\RG{32.041}{0.004}{10.9091}{mi}{116}\RG{42.667}{0.004}{10.9091}{rm}{43}\RG{53.399}{0.004}{10.9091}{mi}{99}\RG{72.215}{0.004}{10.9091}{mi}{116}\RG{92.715}{0.004}{10.9091}{rm}{43}\RG{103.447}{0.004}{10.9091}{sy}{1}\RG{108.291}{0.004}{10.9091}{sy}{1}\RG{113.145}{0.004}{10.9091}{sy}{1}\RG{118.599}{0.004}{10.9091}{rm}{43}\RG{129.331}{0.004}{10.9091}{mi}{99}\RG{138.248}{0.004}{10.9091}{mi}{59}\RG{57.997}{-1.633}{7.9701}{mi}{112}\RG{61.887}{-1.633}{7.9701}{sy}{0}\RG{67.896}{-1.633}{7.9701}{rm}{49}\RG{133.929}{-1.633}{7.9701}{rm}{48}%
}
\]
则所用Frobenius块为

\[
\RB{160.857}{58.030}{3.545}{%
\RG{36.582}{57.175}{9.0910}{ex}{48}\RG{148.457}{57.175}{9.0910}{ex}{49}\RG{50.566}{48.957}{10.9091}{rm}{48}\RG{77.054}{48.957}{10.9091}{rm}{49}\RG{132.135}{48.957}{10.9091}{rm}{48}\RG{36.582}{41.174}{9.0910}{ex}{66}\RG{148.457}{41.174}{9.0910}{ex}{67}\RT{74.065}{36.949}{10.9091}{.}\RT{99.129}{36.949}{10.9091}{.}\RG{36.582}{35.720}{9.0910}{ex}{66}\RG{148.457}{35.720}{9.0910}{ex}{67}\RT{78.244}{34.221}{10.9091}{.}\RT{103.308}{34.221}{10.9091}{.}\RT{82.424}{31.494}{10.9091}{.}\RT{107.488}{31.494}{10.9091}{.}\RG{36.582}{30.265}{9.0910}{ex}{66}\RG{148.457}{30.265}{9.0910}{ex}{67}\RG{0.000}{24.810}{10.9091}{mi}{70}\RG{8.399}{24.810}{10.9091}{rm}{40}\RG{12.553}{24.810}{10.9091}{mi}{103}\RG{18.059}{24.810}{10.9091}{rm}{41}\RG{25.237}{24.810}{10.9091}{rm}{61}\RG{36.582}{24.810}{9.0910}{ex}{66}\RG{148.457}{24.810}{9.0910}{ex}{67}\RG{158.230}{24.810}{10.9091}{mi}{59}\RG{36.582}{19.356}{9.0910}{ex}{66}\RG{148.457}{19.356}{9.0910}{ex}{67}\RG{50.566}{15.293}{10.9091}{rm}{48}\RG{102.122}{15.293}{10.9091}{rm}{48}\RG{132.133}{15.293}{10.9091}{rm}{49}\RG{36.582}{13.538}{9.0910}{ex}{64}\RG{148.457}{13.538}{9.0910}{ex}{65}\RG{44.536}{0.002}{10.9091}{sy}{0}\RG{53.022}{0.002}{10.9091}{mi}{99}\RG{71.028}{0.002}{10.9091}{sy}{0}\RG{79.513}{0.002}{10.9091}{mi}{99}\RG{97.519}{0.002}{10.9091}{sy}{1}\RG{102.363}{0.002}{10.9091}{sy}{1}\RG{107.218}{0.002}{10.9091}{sy}{1}\RG{121.160}{0.002}{10.9091}{sy}{0}\RG{129.641}{0.002}{10.9091}{mi}{99}\RG{57.619}{-1.634}{7.9701}{rm}{48}\RG{84.110}{-1.634}{7.9701}{rm}{49}\RG{134.238}{-1.634}{7.9701}{mi}{112}\RG{138.128}{-1.634}{7.9701}{sy}{0}\RG{144.139}{-1.634}{7.9701}{rm}{49}%
}
\]
其中\allowbreak{}\hskip.1em \RB{25.482}{7.752}{2.619}{%
\RG{0.000}{-0.003}{10.9091}{mi}{112}\RG{8.782}{-0.003}{10.9091}{rm}{61}\RG{20.509}{-0.003}{10.9091}{rm}{49}%
}\hskip.1em\allowbreak{} 时理解为\allowbreak{}\hskip.1em \RB{25.131}{8.679}{3.241}{%
\RG{0.000}{-0.003}{10.9091}{rm}{40}\RG{4.154}{-0.003}{10.9091}{sy}{0}\RG{12.638}{-0.003}{10.9091}{mi}{99}\RG{17.236}{-1.639}{7.9701}{rm}{48}\RG{21.554}{-0.003}{10.9091}{rm}{41}%
}\hskip.1em\allowbreak{} 。按行展开\allowbreak{}\hskip.1em \RB{67.854}{8.679}{3.241}{%
\RG{0.000}{-0.003}{10.9091}{rm}{100}\RG{5.935}{-0.003}{10.9091}{rm}{101}\RG{10.680}{-0.003}{10.9091}{rm}{116}\RG{14.837}{-0.003}{10.9091}{rm}{40}\RG{18.990}{-0.003}{10.9091}{mi}{116}\RG{22.852}{-0.003}{10.9091}{mi}{73}\RG{30.997}{-0.003}{10.9091}{sy}{0}\RG{42.062}{-0.003}{10.9091}{mi}{70}\RG{50.461}{-0.003}{10.9091}{rm}{40}\RG{54.615}{-0.003}{10.9091}{mi}{103}\RG{60.121}{-0.003}{10.9091}{rm}{41}\RG{64.277}{-0.003}{10.9091}{rm}{41}%
}\hskip.1em\allowbreak{} 得\allowbreak{}\hskip.1em \RB{17.094}{8.679}{3.241}{%
\RG{0.000}{-0.003}{10.9091}{mi}{103}\RG{5.506}{-0.003}{10.9091}{rm}{40}\RG{9.661}{-0.003}{10.9091}{mi}{116}\RG{13.518}{-0.003}{10.9091}{rm}{41}%
}\hskip.1em\allowbreak{} ；\allowbreak{}\hskip.1em \RB{22.278}{8.679}{3.241}{%
\RG{0.000}{-0.003}{10.9091}{mi}{70}\RG{8.399}{-0.003}{10.9091}{rm}{40}\RG{12.553}{-0.003}{10.9091}{mi}{104}\RG{18.702}{-0.003}{10.9091}{rm}{41}%
}\hskip.1em\allowbreak{} 同理。因此块上三角阵\allowbreak{}\hskip.1em \RB{11.606}{7.951}{0.500}{%
\RG{0.000}{0.000}{10.9091}{mi}{77}%
}\hskip.1em\allowbreak{} 的特征多项式为\allowbreak{}\hskip.1em \RB{11.835}{8.071}{2.725}{%
\RG{0.000}{0.000}{10.9091}{mi}{103}\RG{5.498}{0.000}{10.9091}{mi}{104}%
}\hskip.1em\allowbreak{} 。

矩阵\allowbreak{}\hskip.1em \RB{11.606}{7.952}{0.500}{%
\RG{0.000}{0.001}{10.9091}{mi}{77}%
}\hskip.1em\allowbreak{} 的超对角元全为\allowbreak{}\hskip.1em \RB{4.973}{7.756}{0.500}{%
\RG{0.000}{0.001}{10.9091}{rm}{49}%
}\hskip.1em\allowbreak{} ，超对角线以上的元素全为\allowbreak{}\hskip.1em \RB{5.387}{7.756}{0.728}{%
\RG{0.000}{0.001}{10.9091}{rm}{48}%
}\hskip.1em\allowbreak{} 。故归纳得到

\[
\RB{257.708}{9.820}{3.620}{%
\RG{7.744}{4.505}{7.9701}{mi}{106}\RG{0.000}{0.001}{10.9091}{mi}{84}\RG{11.743}{0.001}{10.9091}{mi}{101}\RG{24.262}{0.001}{10.9091}{sy}{0}\RG{35.171}{0.001}{10.9091}{mi}{101}\RG{57.850}{0.001}{10.9091}{sy}{50}\RG{68.153}{0.001}{10.9091}{rm}{115}\RG{72.364}{0.001}{10.9091}{rm}{112}\RG{78.298}{0.001}{10.9091}{rm}{97}\RG{83.644}{0.001}{10.9091}{rm}{110}\RG{89.629}{0.001}{10.9091}{sy}{102}\RG{95.083}{0.001}{10.9091}{mi}{101}\RG{124.605}{0.001}{10.9091}{mi}{59}\RG{129.384}{0.001}{10.9091}{mi}{58}\RG{134.173}{0.001}{10.9091}{mi}{58}\RG{138.951}{0.001}{10.9091}{mi}{58}\RG{143.740}{0.001}{10.9091}{mi}{59}\RG{148.529}{0.001}{10.9091}{mi}{101}\RG{158.626}{0.001}{10.9091}{sy}{103}\RG{164.081}{0.001}{10.9091}{mi}{59}\RG{190.684}{0.001}{10.9091}{rm}{48}\RG{199.056}{0.001}{10.9091}{sy}{20}\RG{210.571}{0.001}{10.9091}{mi}{106}\RG{218.634}{0.001}{10.9091}{sy}{20}\RG{230.149}{0.001}{10.9091}{mi}{110}\RG{238.951}{0.001}{10.9091}{sy}{0}\RG{249.860}{0.001}{10.9091}{rm}{49}\RG{255.201}{0.001}{10.9091}{mi}{58}\RG{16.694}{-1.636}{7.9701}{mi}{110}\RG{40.121}{-1.636}{7.9701}{mi}{110}\RG{44.810}{-1.636}{7.9701}{sy}{0}\RG{50.821}{-1.636}{7.9701}{mi}{106}\RG{100.034}{-1.636}{7.9701}{mi}{110}\RG{104.723}{-1.636}{7.9701}{sy}{0}\RG{110.733}{-1.636}{7.9701}{mi}{106}\RG{114.277}{-1.636}{7.9701}{rm}{43}\RG{120.284}{-1.636}{7.9701}{rm}{49}\RG{153.483}{-1.636}{7.9701}{mi}{110}%
}
\]
于是\allowbreak{}\hskip.1em \RB{89.308}{9.296}{2.605}{%
\RG{0.000}{0.000}{10.9091}{mi}{101}\RG{4.951}{-1.636}{7.9701}{mi}{110}\RG{10.094}{0.000}{10.9091}{mi}{59}\RG{14.872}{0.000}{10.9091}{mi}{84}\RG{22.618}{0.000}{10.9091}{mi}{101}\RG{27.574}{-1.636}{7.9701}{mi}{110}\RG{32.718}{0.000}{10.9091}{mi}{59}\RG{37.496}{0.000}{10.9091}{mi}{58}\RG{42.286}{0.000}{10.9091}{mi}{58}\RG{47.064}{0.000}{10.9091}{mi}{58}\RG{51.853}{0.000}{10.9091}{mi}{59}\RG{56.642}{0.000}{10.9091}{mi}{84}\RG{64.389}{3.959}{7.9701}{mi}{110}\RG{69.077}{3.959}{7.9701}{sy}{0}\RG{75.088}{3.959}{7.9701}{rm}{49}\RG{79.406}{0.000}{10.9091}{mi}{101}\RG{84.357}{-1.636}{7.9701}{mi}{110}%
}\hskip.1em\allowbreak{} 线性无关，\allowbreak{}\hskip.1em \RB{8.049}{7.897}{0.500}{%
\RG{0.000}{0.000}{10.9091}{mi}{84}%
}\hskip.1em\allowbreak{} 的极小多项式次数为\allowbreak{}\hskip.1em \RB{6.555}{5.311}{0.609}{%
\RG{0.000}{0.000}{10.9091}{mi}{110}%
}\hskip.1em\allowbreak{} 。结合Cayley–Hamilton定理，得到\allowbreak{}\hskip.1em \RB{56.346}{8.183}{2.723}{%
\RG{0.000}{0.003}{10.9091}{mi}{109}\RG{12.374}{0.003}{10.9091}{rm}{61}\RG{23.713}{0.003}{10.9091}{mi}{102}\RG{33.173}{0.003}{10.9091}{rm}{61}\RG{44.511}{0.003}{10.9091}{mi}{103}\RG{50.010}{0.003}{10.9091}{mi}{104}%
}\hskip.1em\allowbreak{} 。由式(2)：

•若\allowbreak{}\hskip.1em \RB{26.212}{8.316}{2.852}{%
\RG{0.000}{0.005}{10.9091}{mi}{103}\RG{8.537}{0.005}{10.9091}{sy}{54}\RG{8.537}{0.005}{10.9091}{rm}{61}\RG{19.875}{0.005}{10.9091}{mi}{104}%
}\hskip.1em\allowbreak{} ，则所有不变子空间为

\[
\RB{181.272}{8.687}{3.234}{%
\RG{0.000}{0.005}{10.9091}{sy}{102}\RG{5.454}{0.005}{10.9091}{rm}{48}\RG{10.795}{0.005}{10.9091}{sy}{103}\RG{16.250}{0.005}{10.9091}{mi}{59}\RG{42.854}{0.005}{10.9091}{rm}{107}\RG{48.200}{0.005}{10.9091}{rm}{101}\RG{52.945}{0.005}{10.9091}{rm}{114}\RG{58.915}{0.005}{10.9091}{mi}{103}\RG{64.421}{0.005}{10.9091}{rm}{40}\RG{68.574}{0.005}{10.9091}{mi}{84}\RG{76.319}{0.005}{10.9091}{rm}{41}\RG{80.474}{0.005}{10.9091}{mi}{59}\RG{107.077}{0.005}{10.9091}{rm}{107}\RG{112.423}{0.005}{10.9091}{rm}{101}\RG{117.168}{0.005}{10.9091}{rm}{114}\RG{123.137}{0.005}{10.9091}{mi}{104}\RG{129.287}{0.005}{10.9091}{rm}{40}\RG{133.440}{0.005}{10.9091}{mi}{84}\RG{141.185}{0.005}{10.9091}{rm}{41}\RG{145.339}{0.005}{10.9091}{mi}{59}\RG{171.935}{0.005}{10.9091}{mi}{86}\RG{178.765}{0.005}{10.9091}{mi}{58}%
}
\]
•若\allowbreak{}\hskip.1em \RB{26.212}{8.067}{2.729}{%
\RG{0.000}{-0.004}{10.9091}{mi}{103}\RG{8.537}{-0.004}{10.9091}{rm}{61}\RG{19.875}{-0.004}{10.9091}{mi}{104}%
}\hskip.1em\allowbreak{} ，则所有不变子空间为

\[
\RB{116.407}{8.678}{3.242}{%
\RG{0.000}{-0.004}{10.9091}{sy}{102}\RG{5.455}{-0.004}{10.9091}{rm}{48}\RG{10.796}{-0.004}{10.9091}{sy}{103}\RG{16.251}{-0.004}{10.9091}{mi}{59}\RG{42.854}{-0.004}{10.9091}{rm}{107}\RG{48.200}{-0.004}{10.9091}{rm}{101}\RG{52.945}{-0.004}{10.9091}{rm}{114}\RG{58.915}{-0.004}{10.9091}{mi}{103}\RG{64.422}{-0.004}{10.9091}{rm}{40}\RG{68.575}{-0.004}{10.9091}{mi}{84}\RG{76.320}{-0.004}{10.9091}{rm}{41}\RG{80.474}{-0.004}{10.9091}{mi}{59}\RG{107.070}{-0.004}{10.9091}{mi}{86}\RG{113.900}{-0.004}{10.9091}{mi}{58}%
}
\]
两种情形均有\allowbreak{}\hskip.1em \RB{126.567}{8.679}{3.550}{%
\RG{0.000}{-0.003}{10.9091}{rm}{107}\RG{5.346}{-0.003}{10.9091}{rm}{101}\RG{10.091}{-0.003}{10.9091}{rm}{114}\RG{16.060}{-0.003}{10.9091}{mi}{103}\RG{21.566}{-0.003}{10.9091}{rm}{40}\RG{25.721}{-0.003}{10.9091}{mi}{84}\RG{33.465}{-0.003}{10.9091}{rm}{41}\RG{40.643}{-0.003}{10.9091}{rm}{61}\RG{51.989}{-0.003}{10.9091}{rm}{115}\RG{56.200}{-0.003}{10.9091}{rm}{112}\RG{62.134}{-0.003}{10.9091}{rm}{97}\RG{67.480}{-0.003}{10.9091}{rm}{110}\RG{73.464}{-0.003}{10.9091}{sy}{102}\RG{78.918}{-0.003}{10.9091}{mi}{101}\RG{83.869}{-1.639}{7.9701}{rm}{49}\RG{88.187}{-0.003}{10.9091}{mi}{59}\RG{92.965}{-0.003}{10.9091}{mi}{58}\RG{97.754}{-0.003}{10.9091}{mi}{58}\RG{102.544}{-0.003}{10.9091}{mi}{58}\RG{107.322}{-0.003}{10.9091}{mi}{59}\RG{112.111}{-0.003}{10.9091}{mi}{101}\RG{117.064}{-1.639}{7.9701}{mi}{112}\RG{121.409}{-0.003}{10.9091}{sy}{103}%
}\hskip.1em\allowbreak{} ：右侧被\allowbreak{}\hskip.1em \RB{20.981}{8.679}{3.241}{%
\RG{0.000}{-0.003}{10.9091}{mi}{103}\RG{5.506}{-0.003}{10.9091}{rm}{40}\RG{9.661}{-0.003}{10.9091}{mi}{84}\RG{17.404}{-0.003}{10.9091}{rm}{41}%
}\hskip.1em\allowbreak{} 零化，而两侧的维数都为\allowbreak{}\hskip.1em \RB{5.715}{5.308}{2.619}{%
\RG{0.000}{-0.003}{10.9091}{mi}{112}%
}\hskip.1em\allowbreak{} 。

% Original page 4.
(v)所有不变子空间恰为

\[
\RB{297.783}{17.369}{44.631}{%
\RG{80.269}{16.004}{9.0910}{ex}{32}\RG{133.942}{16.004}{9.0910}{ex}{33}\RG{91.610}{11.822}{7.9701}{mi}{107}\RG{160.268}{11.822}{7.9701}{mi}{107}\RG{87.911}{9.094}{9.0910}{ex}{89}\RG{155.509}{9.094}{9.0910}{ex}{77}\RG{127.812}{4.507}{7.9701}{mi}{116}\RG{213.147}{4.507}{7.9701}{mi}{116}\RG{130.602}{3.398}{5.9776}{mi}{105}\RG{215.938}{3.398}{5.9776}{mi}{105}\RG{19.674}{0.003}{10.9091}{mi}{85}\RG{54.688}{0.003}{10.9091}{rm}{61}\RG{66.023}{0.003}{10.9091}{rm}{107}\RG{71.379}{0.003}{10.9091}{rm}{101}\RG{76.125}{0.003}{10.9091}{rm}{114}\RG{101.789}{0.003}{10.9091}{mi}{80}\RG{111.760}{0.003}{10.9091}{rm}{40}\RG{115.914}{0.003}{10.9091}{mi}{84}\RG{123.658}{0.003}{10.9091}{rm}{41}\RG{144.170}{0.003}{10.9091}{rm}{61}\RG{171.064}{0.003}{10.9091}{rm}{107}\RG{176.410}{0.003}{10.9091}{rm}{101}\RG{181.155}{0.003}{10.9091}{rm}{114}\RG{187.124}{0.003}{10.9091}{mi}{80}\RG{197.095}{0.003}{10.9091}{rm}{40}\RG{201.249}{0.003}{10.9091}{mi}{84}\RG{208.994}{0.003}{10.9091}{rm}{41}\RG{219.278}{0.003}{10.9091}{mi}{59}\RG{245.882}{0.003}{10.9091}{rm}{48}\RG{254.253}{0.003}{10.9091}{sy}{20}\RG{265.768}{0.003}{10.9091}{mi}{116}\RG{275.741}{0.003}{10.9091}{sy}{20}\RG{287.256}{0.003}{10.9091}{mi}{114}\RG{295.156}{0.003}{10.9091}{mi}{59}\RG{26.929}{-1.633}{7.9701}{mi}{116}\RG{33.507}{-1.633}{7.9701}{mi}{59}\RG{35.653}{-1.633}{7.9701}{mi}{58}\RG{37.798}{-1.633}{7.9701}{mi}{58}\RG{39.944}{-1.633}{7.9701}{mi}{58}\RG{42.089}{-1.633}{7.9701}{mi}{59}\RG{44.235}{-1.633}{7.9701}{mi}{116}\RG{108.674}{-1.633}{7.9701}{mi}{105}\RG{194.010}{-1.633}{7.9701}{mi}{105}\RG{269.626}{-1.633}{7.9701}{mi}{105}\RG{292.071}{-1.633}{7.9701}{mi}{105}\RG{29.720}{-2.644}{5.9776}{rm}{49}\RG{47.030}{-2.916}{5.9776}{mi}{107}\RG{87.466}{-10.865}{7.9701}{mi}{105}\RG{90.097}{-10.865}{7.9701}{rm}{61}\RG{96.104}{-10.865}{7.9701}{rm}{49}\RG{156.125}{-10.865}{7.9701}{mi}{105}\RG{158.755}{-10.865}{7.9701}{rm}{61}\RG{164.763}{-10.865}{7.9701}{rm}{49}\RG{70.484}{-21.371}{7.9701}{mi}{107}\RG{66.027}{-24.099}{9.0910}{ex}{88}\RG{0.000}{-33.189}{10.9091}{rm}{100}\RG{5.935}{-33.189}{10.9091}{rm}{105}\RG{8.902}{-33.189}{10.9091}{rm}{109}\RG{19.674}{-33.189}{10.9091}{mi}{85}\RG{54.688}{-33.189}{10.9091}{rm}{61}\RG{80.976}{-33.189}{10.9091}{mi}{116}\RG{89.737}{-33.189}{10.9091}{rm}{100}\RG{95.672}{-33.189}{10.9091}{rm}{101}\RG{100.418}{-33.189}{10.9091}{rm}{103}\RG{107.722}{-33.189}{10.9091}{mi}{80}\RG{117.693}{-33.189}{10.9091}{mi}{58}\RG{26.929}{-34.826}{7.9701}{mi}{116}\RG{33.507}{-34.826}{7.9701}{mi}{59}\RG{35.653}{-34.826}{7.9701}{mi}{58}\RG{37.798}{-34.826}{7.9701}{mi}{58}\RG{39.944}{-34.826}{7.9701}{mi}{58}\RG{42.089}{-34.826}{7.9701}{mi}{59}\RG{44.235}{-34.826}{7.9701}{mi}{116}\RG{84.834}{-34.826}{7.9701}{mi}{105}\RG{114.608}{-34.826}{7.9701}{mi}{105}\RG{29.720}{-35.836}{5.9776}{rm}{49}\RG{47.030}{-36.108}{5.9776}{mi}{107}\RG{66.341}{-44.058}{7.9701}{mi}{105}\RG{68.971}{-44.058}{7.9701}{rm}{61}\RG{74.979}{-44.058}{7.9701}{rm}{49}%
}
\]
式(2)给出第一式及维数公式；直和表达式由式(1)和与(iii)相同的互素性论证得到。它们互不相同，共有\allowbreak{}\hskip.1em \RB{45.806}{8.678}{3.242}{%
\RG{0.000}{7.269}{9.0910}{ex}{81}\RG{8.586}{-2.277}{7.9701}{mi}{105}\RG{11.671}{-0.004}{10.9091}{rm}{40}\RG{15.826}{-0.004}{10.9091}{mi}{114}\RG{20.639}{-1.640}{7.9701}{mi}{105}\RG{26.149}{-0.004}{10.9091}{rm}{43}\RG{36.884}{-0.004}{10.9091}{rm}{49}\RG{42.229}{-0.004}{10.9091}{rm}{41}%
}\hskip.1em\allowbreak{} 个。\allowbreak{}\hskip.1em \RB{8.366}{7.991}{0.504}{%
\RG{0.000}{-0.004}{10.9091}{am}{3}%
}\hskip.1em\allowbreak{}

\section*{2017S10}
设\allowbreak{}\hskip.1em \RB{8.715}{7.955}{0.725}{%
\RG{0.000}{0.004}{10.9091}{mi}{86}%
}\hskip.1em\allowbreak{} 是数域\allowbreak{}\hskip.1em \RB{8.889}{7.976}{0.500}{%
\RG{0.000}{0.004}{10.9091}{bb}{75}%
}\hskip.1em\allowbreak{} 上的\allowbreak{}\hskip.1em \RB{6.555}{5.315}{0.605}{%
\RG{0.000}{0.004}{10.9091}{mi}{110}%
}\hskip.1em\allowbreak{} 维线性空间，\allowbreak{}\hskip.1em \RB{7.122}{5.315}{2.853}{%
\RG{0.000}{0.004}{10.9091}{mi}{39}%
}\hskip.1em\allowbreak{} 是\allowbreak{}\hskip.1em \RB{8.715}{7.955}{0.725}{%
\RG{0.000}{0.004}{10.9091}{mi}{86}%
}\hskip.1em\allowbreak{} 上的线性变换。证明：\allowbreak{}\hskip.1em \RB{7.122}{5.315}{2.853}{%
\RG{0.000}{0.004}{10.9091}{mi}{39}%
}\hskip.1em\allowbreak{} 的极小多项式在\allowbreak{}\hskip.1em \RB{8.889}{7.970}{0.503}{%
\RG{0.000}{-0.003}{10.9091}{bb}{75}%
}\hskip.1em\allowbreak{} 上无重因式，当且仅当对\allowbreak{}\hskip.1em \RB{8.715}{7.948}{0.732}{%
\RG{0.000}{-0.003}{10.9091}{mi}{86}%
}\hskip.1em\allowbreak{} 的任一\allowbreak{}\hskip.1em \RB{7.122}{5.308}{2.859}{%
\RG{0.000}{-0.003}{10.9091}{mi}{39}%
}\hskip.1em\allowbreak{} -不变子空间\allowbreak{}\hskip.1em \RB{8.562}{7.948}{0.732}{%
\RG{0.000}{-0.003}{10.9091}{mi}{85}%
}\hskip.1em\allowbreak{} ，均存在\allowbreak{}\hskip.1em \RB{7.122}{5.308}{2.859}{%
\RG{0.000}{-0.003}{10.9091}{mi}{39}%
}\hskip.1em\allowbreak{} -不变子空间\allowbreak{}\hskip.1em \RB{11.693}{7.948}{0.732}{%
\RG{0.000}{-0.003}{10.9091}{mi}{87}%
}\hskip.1em\allowbreak{} ，使

\[
\RB{67.676}{9.226}{2.683}{%
\RG{0.000}{-0.001}{10.9091}{mi}{86}\RG{11.632}{-0.001}{10.9091}{rm}{61}\RG{22.970}{-0.001}{10.9091}{mi}{85}\RG{33.789}{-0.001}{10.9091}{sy}{8}\RG{44.698}{-0.001}{10.9091}{mi}{87}\RT{56.266}{-0.001}{10.9091}{。}%
}
\]
证明：记\allowbreak{}\hskip.1em \RB{29.236}{7.894}{2.859}{%
\RG{0.000}{-0.003}{10.9091}{mi}{84}\RG{10.775}{-0.003}{10.9091}{rm}{61}\RG{22.114}{-0.003}{10.9091}{mi}{39}%
}\hskip.1em\allowbreak{} ，其极小多项式为\allowbreak{}\hskip.1em \RB{9.522}{5.308}{0.612}{%
\RG{0.000}{-0.003}{10.9091}{mi}{109}%
}\hskip.1em\allowbreak{} 。先设

\[
\RB{61.773}{5.316}{2.612}{%
\RG{0.000}{0.005}{10.9091}{mi}{109}\RG{12.375}{0.005}{10.9091}{rm}{61}\RG{23.713}{0.005}{10.9091}{mi}{112}\RG{35.210}{0.005}{10.9091}{sy}{1}\RG{40.054}{0.005}{10.9091}{sy}{1}\RG{44.909}{0.005}{10.9091}{sy}{1}\RG{49.756}{0.005}{10.9091}{mi}{112}\RG{59.145}{0.005}{10.9091}{mi}{59}\RG{29.074}{-1.632}{7.9701}{rm}{49}\RG{55.117}{-1.632}{7.9701}{mi}{115}%
}
\]
其中\allowbreak{}\hskip.1em \RB{8.166}{5.313}{2.615}{%
\RG{0.000}{0.002}{10.9091}{mi}{112}\RG{5.361}{-1.634}{7.9701}{mi}{105}%
}\hskip.1em\allowbreak{} 为互异的首一不可约多项式。由问题2017S09中已证明的式(1)，

\[
\RB{286.686}{15.526}{11.441}{%
\RG{28.052}{11.819}{7.9701}{mi}{115}\RG{102.287}{11.819}{7.9701}{mi}{115}\RG{22.970}{9.092}{9.0910}{ex}{77}\RG{97.205}{9.092}{9.0910}{ex}{77}\RG{0.000}{0.001}{10.9091}{mi}{86}\RG{11.631}{0.001}{10.9091}{rm}{61}\RG{38.525}{0.001}{10.9091}{mi}{86}\RG{47.838}{0.001}{10.9091}{mi}{59}\RG{74.435}{0.001}{10.9091}{mi}{85}\RG{85.867}{0.001}{10.9091}{rm}{61}\RG{112.761}{0.001}{10.9091}{mi}{85}\RG{123.102}{0.001}{10.9091}{mi}{59}\RG{149.698}{0.001}{10.9091}{mi}{86}\RG{162.049}{0.001}{10.9091}{rm}{61}\RG{173.383}{0.001}{10.9091}{rm}{107}\RG{178.740}{0.001}{10.9091}{rm}{101}\RG{183.485}{0.001}{10.9091}{rm}{114}\RG{189.448}{0.001}{10.9091}{mi}{112}\RG{197.894}{0.001}{10.9091}{rm}{40}\RG{202.049}{0.001}{10.9091}{mi}{84}\RG{209.792}{0.001}{10.9091}{rm}{41}\RG{213.947}{0.001}{10.9091}{mi}{59}\RG{229.634}{0.001}{10.9091}{mi}{85}\RG{243.011}{0.001}{10.9091}{rm}{61}\RG{254.350}{0.001}{10.9091}{mi}{85}\RG{265.169}{0.001}{10.9091}{sy}{92}\RG{274.866}{0.001}{10.9091}{mi}{86}\RG{284.178}{0.001}{10.9091}{mi}{58}\RG{44.753}{-1.635}{7.9701}{mi}{105}\RG{120.016}{-1.635}{7.9701}{mi}{105}\RG{155.933}{-1.635}{7.9701}{mi}{105}\RG{194.809}{-1.635}{7.9701}{mi}{105}\RG{236.896}{-1.635}{7.9701}{mi}{105}\RG{281.093}{-1.635}{7.9701}{mi}{105}\RG{23.587}{-10.868}{7.9701}{mi}{105}\RG{26.217}{-10.868}{7.9701}{rm}{61}\RG{32.225}{-10.868}{7.9701}{rm}{49}\RG{97.822}{-10.868}{7.9701}{mi}{105}\RG{100.453}{-10.868}{7.9701}{rm}{61}\RG{106.460}{-10.868}{7.9701}{rm}{49}%
}
\]
令\allowbreak{}\hskip.1em \RB{64.580}{8.685}{3.236}{%
\RG{0.000}{0.003}{10.9091}{mi}{76}\RG{7.267}{-1.633}{7.9701}{mi}{105}\RG{13.383}{0.003}{10.9091}{rm}{61}\RG{24.722}{0.003}{10.9091}{bb}{75}\RG{33.207}{0.003}{10.9091}{rm}{91}\RG{36.237}{0.003}{10.9091}{mi}{116}\RG{40.095}{0.003}{10.9091}{rm}{93}\RG{43.061}{0.003}{10.9091}{mi}{61}\RG{48.402}{0.003}{10.9091}{rm}{40}\RG{52.557}{0.003}{10.9091}{mi}{112}\RG{57.918}{-1.633}{7.9701}{mi}{105}\RG{61.003}{0.003}{10.9091}{rm}{41}%
}\hskip.1em\allowbreak{} 。它是一个域：任意非零剩余类\allowbreak{}\hskip.1em \RB{10.829}{8.685}{3.225}{%
\RG{0.000}{0.003}{10.9091}{rm}{91}\RG{3.030}{0.003}{10.9091}{mi}{97}\RG{8.638}{0.003}{10.9091}{rm}{93}%
}\hskip.1em\allowbreak{} 均满足\allowbreak{}\hskip.1em \RB{62.513}{8.685}{3.236}{%
\RG{0.000}{0.003}{10.9091}{rm}{103}\RG{5.346}{0.003}{10.9091}{rm}{99}\RG{10.091}{0.003}{10.9091}{rm}{100}\RG{16.026}{0.003}{10.9091}{rm}{40}\RG{20.177}{0.003}{10.9091}{mi}{97}\RG{25.784}{0.003}{10.9091}{mi}{59}\RG{30.562}{0.003}{10.9091}{mi}{112}\RG{35.930}{-1.633}{7.9701}{mi}{105}\RG{39.016}{0.003}{10.9091}{rm}{41}\RG{46.195}{0.003}{10.9091}{rm}{61}\RG{57.540}{0.003}{10.9091}{rm}{49}%
}\hskip.1em\allowbreak{} ，Bézout等式给出其乘法逆元。在\allowbreak{}\hskip.1em \RB{9.032}{7.947}{2.213}{%
\RG{0.000}{-0.004}{10.9091}{mi}{86}\RG{6.227}{-1.640}{7.9701}{mi}{105}%
}\hskip.1em\allowbreak{} 上定义

\[
\RB{69.117}{8.681}{3.239}{%
\RG{0.000}{-0.001}{10.9091}{rm}{91}\RG{3.030}{-0.001}{10.9091}{mi}{97}\RG{8.638}{-0.001}{10.9091}{rm}{93}\RG{14.030}{-0.001}{10.9091}{sy}{1}\RG{19.484}{-0.001}{10.9091}{mi}{118}\RG{28.070}{-0.001}{10.9091}{rm}{61}\RG{39.409}{-0.001}{10.9091}{mi}{97}\RG{45.016}{-0.001}{10.9091}{rm}{40}\RG{49.170}{-0.001}{10.9091}{mi}{84}\RG{56.914}{-0.001}{10.9091}{rm}{41}\RG{61.068}{-0.001}{10.9091}{mi}{118}\RG{66.610}{-0.001}{10.9091}{mi}{58}%
}
\]
因\allowbreak{}\hskip.1em \RB{50.892}{8.678}{3.242}{%
\RG{0.000}{-0.004}{10.9091}{mi}{112}\RG{5.361}{-1.640}{7.9701}{mi}{105}\RG{8.447}{-0.004}{10.9091}{rm}{40}\RG{12.601}{-0.004}{10.9091}{mi}{84}\RG{20.345}{-0.004}{10.9091}{rm}{41}\RG{24.499}{-0.004}{10.9091}{mi}{118}\RG{33.624}{-0.004}{10.9091}{rm}{61}\RG{45.504}{-0.004}{10.9091}{rm}{48}%
}\hskip.1em\allowbreak{} ，此定义与代表元的选择无关，使\allowbreak{}\hskip.1em \RB{9.032}{7.947}{2.213}{%
\RG{0.000}{-0.004}{10.9091}{mi}{86}\RG{6.227}{-1.640}{7.9701}{mi}{105}%
}\hskip.1em\allowbreak{} 成为\allowbreak{}\hskip.1em \RB{10.074}{7.947}{2.213}{%
\RG{0.000}{-0.004}{10.9091}{mi}{76}\RG{7.268}{-1.640}{7.9701}{mi}{105}%
}\hskip.1em\allowbreak{} 上的向量空间。\allowbreak{}\hskip.1em \RB{10.061}{7.947}{2.213}{%
\RG{0.000}{-0.004}{10.9091}{mi}{85}\RG{7.255}{-1.640}{7.9701}{mi}{105}%
}\hskip.1em\allowbreak{} 对\allowbreak{}\hskip.1em \RB{8.049}{7.893}{0.504}{%
\RG{0.000}{-0.004}{10.9091}{mi}{84}%
}\hskip.1em\allowbreak{} 的每个多项式封闭，故也是\allowbreak{}\hskip.1em \RB{10.073}{7.950}{2.210}{%
\RG{0.000}{-0.001}{10.9091}{mi}{76}\RG{7.267}{-1.637}{7.9701}{mi}{105}%
}\hskip.1em\allowbreak{} -子空间。将\allowbreak{}\hskip.1em \RB{10.060}{7.950}{2.210}{%
\RG{0.000}{-0.001}{10.9091}{mi}{85}\RG{7.254}{-1.637}{7.9701}{mi}{105}%
}\hskip.1em\allowbreak{} 的一组\allowbreak{}\hskip.1em \RB{10.074}{7.950}{2.210}{%
\RG{0.000}{-0.001}{10.9091}{mi}{76}\RG{7.268}{-1.637}{7.9701}{mi}{105}%
}\hskip.1em\allowbreak{} -基扩充为\allowbreak{}\hskip.1em \RB{9.032}{7.950}{2.210}{%
\RG{0.000}{-0.001}{10.9091}{mi}{86}\RG{6.227}{-1.637}{7.9701}{mi}{105}%
}\hskip.1em\allowbreak{} 的基，以新增基向量张成\allowbreak{}\hskip.1em \RB{12.891}{7.953}{2.208}{%
\RG{0.000}{0.002}{10.9091}{mi}{87}\RG{10.085}{-1.635}{7.9701}{mi}{105}%
}\hskip.1em\allowbreak{} ，得

\[
\RB{63.032}{7.955}{2.205}{%
\RG{0.000}{0.004}{10.9091}{mi}{86}\RG{12.343}{0.004}{10.9091}{rm}{61}\RG{23.682}{0.004}{10.9091}{mi}{85}\RG{36.446}{0.004}{10.9091}{sy}{8}\RG{47.355}{0.004}{10.9091}{mi}{87}\RG{60.525}{0.004}{10.9091}{mi}{58}\RG{6.227}{-1.632}{7.9701}{mi}{105}\RG{30.936}{-1.632}{7.9701}{mi}{105}\RG{57.440}{-1.632}{7.9701}{mi}{105}%
}
\]
由于\allowbreak{}\hskip.1em \RB{50.889}{8.684}{3.225}{%
\RG{0.000}{0.002}{10.9091}{mi}{84}\RG{7.735}{0.002}{10.9091}{mi}{118}\RG{16.330}{0.002}{10.9091}{rm}{61}\RG{27.665}{0.002}{10.9091}{rm}{91}\RG{30.699}{0.002}{10.9091}{mi}{116}\RG{34.556}{0.002}{10.9091}{rm}{93}\RG{39.948}{0.002}{10.9091}{sy}{1}\RG{45.403}{0.002}{10.9091}{mi}{118}%
}\hskip.1em\allowbreak{} ，\allowbreak{}\hskip.1em \RB{12.890}{7.953}{2.208}{%
\RG{0.000}{0.002}{10.9091}{mi}{87}\RG{10.084}{-1.635}{7.9701}{mi}{105}%
}\hskip.1em\allowbreak{} 为\allowbreak{}\hskip.1em \RB{8.049}{7.898}{0.500}{%
\RG{0.000}{0.002}{10.9091}{mi}{84}%
}\hskip.1em\allowbreak{} -不变子空间。取\allowbreak{}\hskip.1em \RB{53.832}{7.953}{2.844}{%
\RG{0.000}{0.002}{10.9091}{mi}{87}\RG{14.599}{0.002}{10.9091}{rm}{61}\RG{25.937}{7.275}{9.0910}{ex}{76}\RG{36.039}{-2.271}{7.9701}{mi}{105}\RG{40.943}{0.002}{10.9091}{mi}{87}\RG{51.027}{-1.635}{7.9701}{mi}{105}%
}\hskip.1em\allowbreak{} ，便有\allowbreak{}\hskip.1em \RB{56.391}{7.953}{1.404}{%
\RG{0.000}{0.002}{10.9091}{mi}{86}\RG{11.632}{0.002}{10.9091}{rm}{61}\RG{22.971}{0.002}{10.9091}{mi}{85}\RG{33.789}{0.002}{10.9091}{sy}{8}\RG{44.699}{0.002}{10.9091}{mi}{87}%
}\hskip.1em\allowbreak{} 。

反之，设每个\allowbreak{}\hskip.1em \RB{8.049}{7.896}{0.501}{%
\RG{0.000}{-0.001}{10.9091}{mi}{84}%
}\hskip.1em\allowbreak{} -不变子空间都有\allowbreak{}\hskip.1em \RB{8.049}{7.896}{0.501}{%
\RG{0.000}{-0.001}{10.9091}{mi}{84}%
}\hskip.1em\allowbreak{} -不变的直和补空间。若\allowbreak{}\hskip.1em \RB{9.522}{5.310}{0.610}{%
\RG{0.000}{-0.001}{10.9091}{mi}{109}%
}\hskip.1em\allowbreak{} 含重不可约因子，可写

\[
\RB{167.475}{8.679}{3.241}{%
\RG{29.074}{4.501}{7.9701}{mi}{114}\RG{0.000}{-0.003}{10.9091}{mi}{109}\RG{12.374}{-0.003}{10.9091}{rm}{61}\RG{23.713}{-0.003}{10.9091}{mi}{112}\RG{33.230}{-0.003}{10.9091}{mi}{113}\RG{38.346}{-0.003}{10.9091}{mi}{59}\RG{64.954}{-0.003}{10.9091}{mi}{114}\RG{73.102}{-0.003}{10.9091}{sy}{21}\RG{84.617}{-0.003}{10.9091}{rm}{50}\RG{89.958}{-0.003}{10.9091}{mi}{59}\RG{105.653}{-0.003}{10.9091}{rm}{103}\RG{110.998}{-0.003}{10.9091}{rm}{99}\RG{115.744}{-0.003}{10.9091}{rm}{100}\RG{121.678}{-0.003}{10.9091}{rm}{40}\RG{125.830}{-0.003}{10.9091}{mi}{112}\RG{131.186}{-0.003}{10.9091}{mi}{59}\RG{135.965}{-0.003}{10.9091}{mi}{113}\RG{141.104}{-0.003}{10.9091}{rm}{41}\RG{148.282}{-0.003}{10.9091}{rm}{61}\RG{159.628}{-0.003}{10.9091}{rm}{49}\RG{164.968}{-0.003}{10.9091}{mi}{58}%
}
\]
对\allowbreak{}\hskip.1em \RB{59.660}{8.679}{3.241}{%
\RG{0.000}{-0.002}{10.9091}{mi}{85}\RG{11.425}{-0.003}{10.9091}{rm}{61}\RG{22.759}{-0.003}{10.9091}{rm}{107}\RG{28.116}{-0.003}{10.9091}{rm}{101}\RG{32.861}{-0.003}{10.9091}{rm}{114}\RG{38.824}{-0.002}{10.9091}{mi}{112}\RG{44.186}{-0.003}{10.9091}{rm}{40}\RG{48.339}{-0.002}{10.9091}{mi}{84}\RG{56.084}{-0.003}{10.9091}{rm}{41}%
}\hskip.1em\allowbreak{} ，取不变补空间\allowbreak{}\hskip.1em \RB{11.693}{7.948}{0.732}{%
\RG{0.000}{-0.003}{10.9091}{mi}{87}%
}\hskip.1em\allowbreak{} ，使\allowbreak{}\hskip.1em \RB{54.311}{7.948}{1.408}{%
\RG{0.000}{-0.003}{10.9091}{mi}{86}\RG{11.632}{-0.002}{10.9091}{rm}{61}\RG{22.971}{-0.003}{10.9091}{mi}{85}\RG{32.749}{-0.002}{10.9091}{sy}{8}\RG{42.618}{-0.003}{10.9091}{mi}{87}%
}\hskip.1em\allowbreak{} 。因为\allowbreak{}\hskip.1em \RB{89.406}{8.680}{3.241}{%
\RG{0.000}{-0.003}{10.9091}{mi}{87}\RG{12.952}{-0.002}{10.9091}{sy}{92}\RG{21.609}{-0.003}{10.9091}{rm}{107}\RG{26.955}{-0.003}{10.9091}{rm}{101}\RG{31.700}{-0.003}{10.9091}{rm}{114}\RG{37.670}{-0.002}{10.9091}{mi}{112}\RG{43.031}{-0.003}{10.9091}{rm}{40}\RG{47.185}{-0.002}{10.9091}{mi}{84}\RG{54.930}{-0.003}{10.9091}{rm}{41}\RG{62.108}{-0.003}{10.9091}{rm}{61}\RG{73.453}{-0.002}{10.9091}{sy}{102}\RG{78.907}{-0.003}{10.9091}{rm}{48}\RG{84.248}{-0.002}{10.9091}{sy}{103}%
}\hskip.1em\allowbreak{} ，算子\allowbreak{}\hskip.1em \RB{33.017}{8.680}{3.241}{%
\RG{0.000}{-0.002}{10.9091}{mi}{112}\RG{5.361}{-0.003}{10.9091}{rm}{40}\RG{9.516}{-0.002}{10.9091}{mi}{84}\RG{17.259}{-0.003}{10.9091}{rm}{41}\RG{21.414}{-0.002}{10.9091}{sy}{106}\RG{24.444}{-1.639}{7.9701}{mi}{87}%
}\hskip.1em\allowbreak{}为单射，故可逆。从

\[
\RB{78.350}{8.682}{3.238}{%
\RG{21.413}{4.504}{7.9701}{mi}{114}\RG{0.000}{0.000}{10.9091}{mi}{112}\RG{5.361}{0.000}{10.9091}{rm}{40}\RG{9.515}{0.000}{10.9091}{mi}{84}\RG{17.259}{0.000}{10.9091}{rm}{41}\RG{25.570}{0.000}{10.9091}{mi}{113}\RG{30.697}{0.000}{10.9091}{rm}{40}\RG{34.851}{0.000}{10.9091}{mi}{84}\RG{42.595}{0.000}{10.9091}{rm}{41}\RG{46.749}{0.000}{10.9091}{sy}{106}\RG{61.628}{0.000}{10.9091}{rm}{61}\RG{72.963}{0.000}{10.9091}{rm}{48}\RG{49.779}{-1.636}{7.9701}{mi}{87}%
}
\]
得到\allowbreak{}\hskip.1em \RB{52.781}{8.679}{3.241}{%
\RG{0.000}{-0.003}{10.9091}{mi}{113}\RG{5.127}{-0.003}{10.9091}{rm}{40}\RG{9.281}{-0.003}{10.9091}{mi}{84}\RG{17.026}{-0.003}{10.9091}{rm}{41}\RG{21.179}{-0.003}{10.9091}{sy}{106}\RG{24.210}{-1.639}{7.9701}{mi}{87}\RG{36.059}{-0.003}{10.9091}{rm}{61}\RG{47.394}{-0.003}{10.9091}{rm}{48}%
}\hskip.1em\allowbreak{} 。另一方面，\allowbreak{}\hskip.1em \RB{50.654}{8.679}{3.241}{%
\RG{0.000}{-0.003}{10.9091}{mi}{112}\RG{5.362}{-0.003}{10.9091}{rm}{40}\RG{9.516}{-0.003}{10.9091}{mi}{84}\RG{17.260}{-0.003}{10.9091}{rm}{41}\RG{21.414}{-0.003}{10.9091}{sy}{106}\RG{24.444}{-1.639}{7.9701}{mi}{85}\RG{33.932}{-0.003}{10.9091}{rm}{61}\RG{45.267}{-0.003}{10.9091}{rm}{48}%
}\hskip.1em\allowbreak{} ，所以\allowbreak{}\hskip.1em \RB{34.272}{8.679}{3.241}{%
\RG{0.000}{-0.003}{10.9091}{rm}{40}\RG{4.154}{-0.003}{10.9091}{mi}{112}\RG{9.510}{-0.003}{10.9091}{mi}{113}\RG{14.643}{-0.003}{10.9091}{rm}{41}\RG{18.800}{-0.003}{10.9091}{rm}{40}\RG{22.951}{-0.003}{10.9091}{mi}{84}\RG{30.695}{-0.003}{10.9091}{rm}{41}%
}\hskip.1em\allowbreak{} 同时零化\allowbreak{}\hskip.1em \RB{8.562}{7.948}{0.732}{%
\RG{0.000}{-0.003}{10.9091}{mi}{85}%
}\hskip.1em\allowbreak{} 与\allowbreak{}\hskip.1em \RB{11.693}{7.948}{0.732}{%
\RG{0.000}{-0.003}{10.9091}{mi}{87}%
}\hskip.1em\allowbreak{} ，即\allowbreak{}\hskip.1em \RB{54.605}{8.679}{3.241}{%
\RG{0.000}{-0.003}{10.9091}{rm}{40}\RG{4.154}{-0.003}{10.9091}{mi}{112}\RG{9.510}{-0.003}{10.9091}{mi}{113}\RG{14.642}{-0.003}{10.9091}{rm}{41}\RG{18.799}{-0.003}{10.9091}{rm}{40}\RG{22.951}{-0.003}{10.9091}{mi}{84}\RG{30.694}{-0.003}{10.9091}{rm}{41}\RG{37.872}{-0.003}{10.9091}{rm}{61}\RG{49.218}{-0.003}{10.9091}{rm}{48}%
}\hskip.1em\allowbreak{} 。但

\[
\RB{132.782}{8.685}{3.235}{%
\RG{75.014}{4.507}{7.9701}{mi}{114}\RG{0.000}{0.003}{10.9091}{rm}{100}\RG{5.935}{0.003}{10.9091}{rm}{101}\RG{10.680}{0.003}{10.9091}{rm}{103}\RG{16.168}{0.003}{10.9091}{rm}{40}\RG{20.321}{0.003}{10.9091}{mi}{112}\RG{25.677}{0.003}{10.9091}{mi}{113}\RG{30.809}{0.003}{10.9091}{rm}{41}\RG{37.994}{0.003}{10.9091}{mi}{60}\RG{49.332}{0.003}{10.9091}{rm}{100}\RG{55.267}{0.003}{10.9091}{rm}{101}\RG{60.012}{0.003}{10.9091}{rm}{103}\RG{65.500}{0.003}{10.9091}{rm}{40}\RG{69.653}{0.003}{10.9091}{mi}{112}\RG{79.171}{0.003}{10.9091}{mi}{113}\RG{84.298}{0.003}{10.9091}{rm}{41}\RG{91.476}{0.003}{10.9091}{rm}{61}\RG{102.822}{0.003}{10.9091}{rm}{100}\RG{108.756}{0.003}{10.9091}{rm}{101}\RG{113.502}{0.003}{10.9091}{rm}{103}\RG{120.806}{0.003}{10.9091}{mi}{109}\RG{130.155}{0.003}{10.9091}{mi}{59}%
}
\]
与\allowbreak{}\hskip.1em \RB{9.522}{5.311}{0.609}{%
\RG{0.000}{0.000}{10.9091}{mi}{109}%
}\hskip.1em\allowbreak{} 的极小性矛盾。因此\allowbreak{}\hskip.1em \RB{9.522}{5.311}{0.609}{%
\RG{0.000}{0.000}{10.9091}{mi}{109}%
}\hskip.1em\allowbreak{} 无重因式。\allowbreak{}\hskip.1em \RB{8.366}{7.995}{0.500}{%
\RG{0.000}{0.000}{10.9091}{am}{3}%
}\hskip.1em\allowbreak{}

% Original page 5.
\section*{2017S11}
设\allowbreak{}\hskip.1em \RB{19.608}{8.683}{3.237}{%
\RG{0.000}{0.001}{10.9091}{mi}{102}\RG{6.430}{0.001}{10.9091}{rm}{40}\RG{10.584}{0.001}{10.9091}{mi}{122}\RG{16.032}{0.001}{10.9091}{rm}{41}%
}\hskip.1em\allowbreak{} 是收敛半径为\allowbreak{}\hskip.1em \RB{19.096}{6.785}{1.339}{%
\RG{0.000}{0.001}{10.9091}{rm}{43}\RG{8.308}{0.001}{10.9091}{sy}{49}%
}\hskip.1em\allowbreak{} 的复幂级数，\allowbreak{}\hskip.1em \RB{52.428}{8.683}{3.237}{%
\RG{0.000}{0.001}{10.9091}{mi}{65}\RG{11.038}{0.001}{10.9091}{sy}{50}\RG{21.341}{0.001}{10.9091}{mi}{77}\RG{31.675}{-1.636}{7.9701}{mi}{110}\RG{36.818}{0.001}{10.9091}{rm}{40}\RG{40.973}{0.001}{10.9091}{bb}{67}\RG{48.851}{0.001}{10.9091}{rm}{41}%
}\hskip.1em\allowbreak{} ，并令

\[
\RB{127.953}{11.408}{5.947}{%
\RG{48.467}{10.544}{9.0910}{ex}{16}\RG{120.016}{10.544}{9.0910}{ex}{17}\RG{0.000}{-0.002}{10.9091}{mi}{103}\RG{5.506}{-0.002}{10.9091}{rm}{40}\RG{9.660}{-0.002}{10.9091}{mi}{122}\RG{15.108}{-0.002}{10.9091}{rm}{41}\RG{22.287}{-0.002}{10.9091}{rm}{61}\RG{33.632}{-0.002}{10.9091}{rm}{100}\RG{39.567}{-0.002}{10.9091}{rm}{101}\RG{44.313}{-0.002}{10.9091}{rm}{116}\RG{53.897}{-0.002}{10.9091}{mi}{102}\RG{60.326}{-0.002}{10.9091}{rm}{40}\RG{64.480}{-0.002}{10.9091}{mi}{122}\RG{69.929}{-0.002}{10.9091}{rm}{41}\RG{74.082}{-0.002}{10.9091}{mi}{73}\RG{86.361}{-0.002}{10.9091}{sy}{0}\RG{97.270}{-0.002}{10.9091}{mi}{102}\RG{103.701}{-0.002}{10.9091}{rm}{40}\RG{107.854}{-0.002}{10.9091}{mi}{65}\RG{115.862}{-0.002}{10.9091}{rm}{41}\RG{125.446}{-0.002}{10.9091}{mi}{58}\RG{78.794}{-1.638}{7.9701}{mi}{110}%
}
\]
证明：\allowbreak{}\hskip.1em \RB{41.579}{8.678}{3.242}{%
\RG{0.000}{-0.004}{10.9091}{mi}{103}\RG{5.506}{-0.004}{10.9091}{rm}{40}\RG{9.660}{-0.004}{10.9091}{mi}{65}\RG{17.668}{-0.004}{10.9091}{rm}{41}\RG{24.846}{-0.004}{10.9091}{rm}{61}\RG{36.192}{-0.004}{10.9091}{rm}{48}%
}\hskip.1em\allowbreak{} 。证明：对任意整函数\allowbreak{}\hskip.1em \RB{79.359}{9.277}{4.253}{%
\RG{0.000}{0.004}{10.9091}{mi}{70}\RG{8.399}{0.004}{10.9091}{rm}{40}\RG{12.553}{0.004}{10.9091}{mi}{122}\RG{18.001}{0.004}{10.9091}{rm}{41}\RG{25.179}{0.004}{10.9091}{rm}{61}\RG{36.523}{7.276}{9.0910}{ex}{80}\RG{46.120}{-2.269}{7.9701}{mi}{106}\RG{49.664}{-2.269}{7.9701}{sy}{21}\RG{55.674}{-2.269}{7.9701}{rm}{48}\RG{61.811}{0.004}{10.9091}{mi}{99}\RG{66.408}{-1.633}{7.9701}{mi}{106}\RG{70.407}{0.004}{10.9091}{mi}{122}\RG{75.855}{3.962}{7.9701}{mi}{106}%
}\hskip.1em\allowbreak{} ，定义

\[
\RB{77.428}{9.820}{12.906}{%
\RG{39.226}{9.092}{9.0910}{ex}{88}\RG{70.922}{4.505}{7.9701}{mi}{106}\RG{0.000}{0.001}{10.9091}{mi}{70}\RG{8.398}{0.001}{10.9091}{rm}{40}\RG{12.552}{0.001}{10.9091}{mi}{65}\RG{20.560}{0.001}{10.9091}{rm}{41}\RG{27.738}{0.001}{10.9091}{rm}{61}\RG{54.319}{0.001}{10.9091}{mi}{99}\RG{62.915}{0.001}{10.9091}{mi}{65}\RG{74.921}{0.001}{10.9091}{mi}{58}\RG{58.916}{-1.635}{7.9701}{mi}{106}\RG{39.083}{-10.922}{7.9701}{mi}{106}\RG{42.627}{-10.922}{7.9701}{sy}{21}\RG{48.637}{-10.922}{7.9701}{rm}{48}%
}
\]
取次乘性矩阵范数，因\allowbreak{}\hskip.1em \RB{79.960}{9.269}{4.260}{%
\RG{0.000}{7.269}{9.0910}{ex}{80}\RG{9.596}{-2.277}{7.9701}{mi}{106}\RG{15.413}{-0.004}{10.9091}{sy}{106}\RG{18.444}{-0.004}{10.9091}{mi}{99}\RG{23.041}{-1.641}{7.9701}{mi}{106}\RG{27.040}{-0.004}{10.9091}{sy}{106}\RG{31.883}{-0.004}{10.9091}{sy}{107}\RG{37.343}{-0.004}{10.9091}{mi}{65}\RG{45.350}{-0.004}{10.9091}{sy}{107}\RG{50.805}{3.955}{7.9701}{mi}{106}\RG{57.834}{-0.004}{10.9091}{mi}{60}\RG{69.172}{-0.004}{10.9091}{sy}{49}%
}\hskip.1em\allowbreak{} ，上述矩阵级数绝对收敛。

若\allowbreak{}\hskip.1em \RB{82.558}{9.509}{4.006}{%
\RG{0.000}{0.003}{10.9091}{mi}{71}\RG{8.426}{0.003}{10.9091}{rm}{40}\RG{12.580}{0.003}{10.9091}{mi}{122}\RG{18.029}{0.003}{10.9091}{rm}{41}\RG{25.207}{0.003}{10.9091}{rm}{61}\RG{36.551}{7.276}{9.0910}{ex}{80}\RG{46.147}{-2.269}{7.9701}{mi}{107}\RG{50.364}{-2.269}{7.9701}{sy}{21}\RG{56.375}{-2.269}{7.9701}{rm}{48}\RG{62.511}{0.003}{10.9091}{mi}{100}\RG{68.062}{-1.633}{7.9701}{mi}{107}\RG{72.734}{0.003}{10.9091}{mi}{122}\RG{78.181}{3.962}{7.9701}{mi}{107}%
}\hskip.1em\allowbreak{} 也是整函数，则

\[
\RB{269.029}{19.578}{14.131}{%
\RG{95.380}{18.724}{9.0910}{ex}{48}\RG{157.961}{18.724}{9.0910}{ex}{49}\RG{167.734}{18.724}{9.0910}{ex}{48}\RG{233.288}{18.724}{9.0910}{ex}{49}\RG{3.111}{9.087}{9.0910}{ex}{88}\RG{103.478}{9.087}{9.0910}{ex}{88}\RG{176.169}{9.087}{9.0910}{ex}{88}\RG{66.785}{4.501}{7.9701}{mi}{106}\RG{70.329}{4.501}{7.9701}{rm}{43}\RG{76.339}{4.501}{7.9701}{mi}{107}\RG{153.963}{4.500}{7.9701}{mi}{106}\RG{228.616}{4.501}{7.9701}{mi}{107}\RG{95.380}{2.360}{9.0910}{ex}{64}\RG{157.961}{2.360}{9.0910}{ex}{65}\RG{167.734}{2.360}{9.0910}{ex}{64}\RG{233.288}{2.360}{9.0910}{ex}{65}\RG{21.171}{-0.004}{10.9091}{sy}{106}\RG{24.202}{-0.004}{10.9091}{mi}{99}\RG{32.798}{-0.004}{10.9091}{mi}{100}\RG{43.019}{-0.004}{10.9091}{sy}{106}\RG{47.863}{-0.004}{10.9091}{sy}{107}\RG{53.323}{-0.004}{10.9091}{mi}{65}\RG{61.330}{-0.004}{10.9091}{sy}{107}\RG{84.042}{-0.004}{10.9091}{rm}{61}\RG{118.571}{-0.004}{10.9091}{sy}{106}\RG{121.601}{-0.004}{10.9091}{mi}{99}\RG{130.197}{-0.004}{10.9091}{sy}{106}\RG{135.041}{-0.004}{10.9091}{sy}{107}\RG{140.500}{-0.004}{10.9091}{mi}{65}\RG{148.509}{-0.004}{10.9091}{sy}{107}\RG{191.598}{-0.004}{10.9091}{sy}{106}\RG{194.629}{-0.004}{10.9091}{mi}{100}\RG{204.851}{-0.004}{10.9091}{sy}{106}\RG{209.695}{-0.004}{10.9091}{sy}{107}\RG{215.154}{-0.004}{10.9091}{mi}{65}\RG{223.162}{-0.004}{10.9091}{sy}{107}\RG{244.273}{-0.004}{10.9091}{mi}{60}\RG{255.612}{-0.004}{10.9091}{sy}{49}\RG{266.521}{-0.004}{10.9091}{mi}{58}\RG{28.799}{-1.641}{7.9701}{mi}{106}\RG{38.347}{-1.641}{7.9701}{mi}{107}\RG{126.199}{-1.641}{7.9701}{mi}{106}\RG{200.180}{-1.641}{7.9701}{mi}{107}\RG{103.335}{-10.927}{7.9701}{mi}{106}\RG{106.879}{-10.927}{7.9701}{sy}{21}\RG{112.890}{-10.927}{7.9701}{rm}{48}\RG{0.000}{-11.115}{7.9701}{mi}{106}\RG{3.120}{-11.115}{7.9701}{mi}{59}\RG{5.266}{-11.115}{7.9701}{mi}{107}\RG{9.478}{-11.115}{7.9701}{sy}{21}\RG{15.489}{-11.115}{7.9701}{rm}{48}\RG{175.689}{-11.115}{7.9701}{mi}{107}\RG{179.906}{-11.115}{7.9701}{sy}{21}\RG{185.916}{-11.115}{7.9701}{rm}{48}%
}
\]
故可按总次数重新排列乘积级数，得到

\[
\RB{107.779}{8.685}{3.235}{%
\RG{0.000}{0.003}{10.9091}{rm}{40}\RG{4.155}{0.003}{10.9091}{mi}{70}\RG{12.544}{0.003}{10.9091}{mi}{71}\RG{20.979}{0.003}{10.9091}{rm}{41}\RG{25.135}{0.003}{10.9091}{rm}{40}\RG{29.286}{0.003}{10.9091}{mi}{65}\RG{37.294}{0.003}{10.9091}{rm}{41}\RG{44.473}{0.003}{10.9091}{rm}{61}\RG{55.817}{0.003}{10.9091}{mi}{70}\RG{64.215}{0.003}{10.9091}{rm}{40}\RG{68.369}{0.003}{10.9091}{mi}{65}\RG{76.377}{0.003}{10.9091}{rm}{41}\RG{80.531}{0.003}{10.9091}{mi}{71}\RG{88.957}{0.003}{10.9091}{rm}{40}\RG{93.111}{0.003}{10.9091}{mi}{65}\RG{101.118}{0.003}{10.9091}{rm}{41}\RG{105.272}{0.003}{10.9091}{mi}{58}%
}
\]
连同加法及常数函数的情形，这说明\allowbreak{}\hskip.1em \RB{49.505}{8.682}{3.238}{%
\RG{0.000}{0.000}{10.9091}{mi}{70}\RG{11.429}{0.000}{10.9091}{sy}{55}\RG{11.429}{0.000}{10.9091}{sy}{33}\RG{25.369}{0.000}{10.9091}{mi}{70}\RG{33.767}{0.000}{10.9091}{rm}{40}\RG{37.921}{0.000}{10.9091}{mi}{65}\RG{45.929}{0.000}{10.9091}{rm}{41}%
}\hskip.1em\allowbreak{} 保持有限次多项式运算。

置\allowbreak{}\hskip.1em \RB{45.203}{8.680}{3.240}{%
\RG{0.000}{-0.002}{10.9091}{mi}{66}\RG{11.696}{-0.002}{10.9091}{rm}{61}\RG{23.035}{-0.002}{10.9091}{mi}{102}\RG{29.464}{-0.002}{10.9091}{rm}{40}\RG{33.618}{-0.002}{10.9091}{mi}{65}\RG{41.626}{-0.002}{10.9091}{rm}{41}%
}\hskip.1em\allowbreak{} ，令\allowbreak{}\hskip.1em \RB{90.960}{8.680}{3.240}{%
\RG{0.000}{-0.002}{10.9091}{mi}{104}\RG{6.149}{-0.002}{10.9091}{rm}{40}\RG{10.303}{-0.002}{10.9091}{mi}{116}\RG{14.160}{-0.002}{10.9091}{rm}{41}\RG{21.339}{-0.002}{10.9091}{rm}{61}\RG{32.684}{-0.002}{10.9091}{rm}{100}\RG{38.619}{-0.002}{10.9091}{rm}{101}\RG{43.364}{-0.002}{10.9091}{rm}{116}\RG{47.521}{-0.002}{10.9091}{rm}{40}\RG{51.673}{-0.002}{10.9091}{mi}{116}\RG{55.535}{-0.002}{10.9091}{mi}{73}\RG{60.241}{-1.638}{7.9701}{mi}{110}\RG{67.809}{-0.002}{10.9091}{sy}{0}\RG{78.718}{-0.002}{10.9091}{mi}{66}\RG{87.384}{-0.002}{10.9091}{rm}{41}%
}\hskip.1em\allowbreak{} 。由定义，\allowbreak{}\hskip.1em \RB{67.700}{8.680}{3.240}{%
\RG{0.000}{-0.002}{10.9091}{mi}{103}\RG{5.506}{-0.002}{10.9091}{rm}{40}\RG{9.660}{-0.002}{10.9091}{mi}{122}\RG{15.108}{-0.002}{10.9091}{rm}{41}\RG{22.287}{-0.002}{10.9091}{rm}{61}\RG{33.632}{-0.002}{10.9091}{mi}{104}\RG{39.780}{-0.002}{10.9091}{rm}{40}\RG{43.935}{-0.002}{10.9091}{mi}{102}\RG{50.364}{-0.002}{10.9091}{rm}{40}\RG{54.518}{-0.002}{10.9091}{mi}{122}\RG{59.967}{-0.002}{10.9091}{rm}{41}\RG{64.123}{-0.002}{10.9091}{rm}{41}%
}\hskip.1em\allowbreak{} 。因此

\[
\RB{81.651}{11.405}{23.037}{%
\RG{42.340}{10.541}{9.0910}{ex}{16}\RG{70.515}{10.541}{9.0910}{ex}{17}\RG{0.000}{-0.004}{10.9091}{mi}{103}\RG{5.506}{-0.004}{10.9091}{rm}{40}\RG{9.661}{-0.004}{10.9091}{mi}{65}\RG{17.668}{-0.004}{10.9091}{rm}{41}\RG{24.846}{-0.004}{10.9091}{rm}{61}\RG{36.191}{-0.004}{10.9091}{mi}{104}\RG{47.769}{-0.004}{10.9091}{mi}{102}\RG{54.200}{-0.004}{10.9091}{rm}{40}\RG{58.353}{-0.004}{10.9091}{mi}{65}\RG{66.361}{-0.004}{10.9091}{rm}{41}\RG{24.852}{-19.799}{10.9091}{rm}{61}\RG{36.191}{-19.799}{10.9091}{mi}{104}\RG{42.340}{-19.799}{10.9091}{rm}{40}\RG{46.494}{-19.799}{10.9091}{mi}{66}\RG{55.160}{-19.799}{10.9091}{rm}{41}\RG{62.338}{-19.799}{10.9091}{rm}{61}\RG{73.683}{-19.799}{10.9091}{rm}{48}\RG{79.024}{-19.799}{10.9091}{mi}{59}%
}
\]
最后一步为Cayley–Hamilton定理，简证见文末。\allowbreak{}\hskip.1em \RB{8.366}{7.996}{0.500}{%
\RG{0.000}{0.001}{10.9091}{am}{3}%
}\hskip.1em\allowbreak{}

\section*{2017S12}
设\allowbreak{}\hskip.1em \RB{8.202}{8.287}{0.502}{%
\RG{0.000}{-0.002}{10.9091}{mi}{65}%
}\hskip.1em\allowbreak{} 为\allowbreak{}\hskip.1em \RB{6.555}{5.309}{0.611}{%
\RG{0.000}{-0.002}{10.9091}{mi}{110}%
}\hskip.1em\allowbreak{} 阶正定实对称阵，\allowbreak{}\hskip.1em \RB{8.606}{7.949}{0.502}{%
\RG{0.000}{-0.002}{10.9091}{mi}{66}%
}\hskip.1em\allowbreak{} 为\allowbreak{}\hskip.1em \RB{6.555}{5.309}{0.611}{%
\RG{0.000}{-0.002}{10.9091}{mi}{110}%
}\hskip.1em\allowbreak{} 阶实矩阵，且

\[
\RB{75.689}{0.856}{32.854}{%
\RG{25.842}{0.001}{9.0910}{ex}{48}\RG{69.889}{0.001}{9.0910}{ex}{49}\RG{62.188}{-7.453}{7.9701}{rm}{84}\RG{34.126}{-11.412}{10.9091}{mi}{65}\RG{53.512}{-11.412}{10.9091}{mi}{66}\RG{25.842}{-16.363}{9.0910}{ex}{64}\RG{69.889}{-16.363}{9.0910}{ex}{65}\RG{0.000}{-18.727}{10.9091}{mi}{77}\RG{14.503}{-18.727}{10.9091}{rm}{61}\RG{59.562}{-22.744}{7.9701}{sy}{0}\RG{65.571}{-22.744}{7.9701}{rm}{49}\RG{33.797}{-26.703}{10.9091}{mi}{66}\RG{51.546}{-26.703}{10.9091}{mi}{65}%
}
\]
为半正定阵。证明：\allowbreak{}\hskip.1em \RB{8.606}{7.950}{0.501}{%
\RG{0.000}{-0.001}{10.9091}{mi}{66}%
}\hskip.1em\allowbreak{} 的特征值均落在复平面的闭单位圆盘内。证明：对每个实列向量\allowbreak{}\hskip.1em \RB{6.107}{5.308}{0.612}{%
\RG{0.000}{-0.003}{10.9091}{mi}{120}%
}\hskip.1em\allowbreak{} ，将\allowbreak{}\hskip.1em \RB{75.660}{9.453}{3.241}{%
\RG{0.000}{-0.003}{10.9091}{rm}{40}\RG{4.154}{-0.003}{10.9091}{mi}{120}\RG{10.224}{3.957}{7.9701}{rm}{84}\RG{15.956}{-0.003}{10.9091}{mi}{59}\RG{20.741}{-0.003}{10.9091}{sy}{0}\RG{29.227}{-0.003}{10.9091}{rm}{40}\RG{33.381}{-0.003}{10.9091}{mi}{65}\RG{41.388}{-0.003}{10.9091}{mi}{66}\RG{50.050}{-0.003}{10.9091}{mi}{120}\RG{56.124}{-0.003}{10.9091}{rm}{41}\RG{60.278}{3.957}{7.9701}{rm}{84}\RG{66.010}{-0.003}{10.9091}{rm}{41}\RG{70.164}{3.957}{7.9701}{rm}{84}%
}\hskip.1em\allowbreak{} 代入\allowbreak{}\hskip.1em \RB{11.606}{7.948}{0.503}{%
\RG{0.000}{-0.003}{10.9091}{mi}{77}%
}\hskip.1em\allowbreak{} 的二次型，得

\[
\RB{132.849}{49.682}{0.613}{%
\RG{19.886}{46.702}{9.0910}{ex}{48}\RG{59.070}{46.702}{9.0910}{ex}{49}\RG{87.867}{46.702}{9.0910}{ex}{48}\RG{127.049}{46.702}{9.0910}{ex}{49}\RG{67.024}{44.186}{7.9701}{rm}{84}\RG{40.421}{35.288}{10.9091}{mi}{120}\RG{108.400}{35.288}{10.9091}{mi}{120}\RG{19.886}{30.337}{9.0910}{ex}{64}\RG{59.070}{30.337}{9.0910}{ex}{65}\RG{87.867}{30.337}{9.0910}{ex}{64}\RG{127.049}{30.337}{9.0910}{ex}{65}\RG{0.000}{27.973}{10.9091}{rm}{48}\RG{8.371}{27.973}{10.9091}{sy}{20}\RG{74.574}{27.973}{10.9091}{mi}{77}\RG{27.842}{19.997}{10.9091}{sy}{0}\RG{36.326}{19.997}{10.9091}{mi}{65}\RG{44.334}{19.997}{10.9091}{mi}{66}\RG{52.996}{19.997}{10.9091}{mi}{120}\RG{95.821}{19.997}{10.9091}{sy}{0}\RG{104.306}{19.997}{10.9091}{mi}{65}\RG{112.313}{19.997}{10.9091}{mi}{66}\RG{120.975}{19.997}{10.9091}{mi}{120}\RG{25.780}{4.501}{7.9701}{rm}{84}\RG{64.994}{4.501}{7.9701}{rm}{84}\RG{79.391}{4.501}{7.9701}{rm}{84}\RG{8.371}{-0.004}{10.9091}{rm}{61}\RG{19.710}{-0.004}{10.9091}{mi}{120}\RG{31.513}{-0.004}{10.9091}{mi}{65}\RG{39.520}{-0.004}{10.9091}{mi}{120}\RG{48.015}{-0.004}{10.9091}{sy}{0}\RG{58.924}{-0.004}{10.9091}{mi}{120}\RG{70.726}{-0.004}{10.9091}{mi}{66}\RG{85.124}{-0.004}{10.9091}{mi}{65}\RG{93.131}{-0.004}{10.9091}{mi}{66}\RG{101.793}{-0.004}{10.9091}{mi}{120}\RG{107.859}{-0.004}{10.9091}{mi}{58}%
}
\]
故实对称阵\allowbreak{}\hskip.1em \RB{73.090}{9.459}{0.725}{%
\RG{0.000}{0.004}{10.9091}{mi}{83}\RG{10.235}{0.004}{10.9091}{rm}{61}\RG{21.573}{0.004}{10.9091}{mi}{65}\RG{31.588}{0.004}{10.9091}{sy}{0}\RG{42.079}{0.004}{10.9091}{mi}{66}\RG{50.745}{3.962}{7.9701}{rm}{84}\RG{56.477}{0.004}{10.9091}{mi}{65}\RG{64.485}{0.004}{10.9091}{mi}{66}%
}\hskip.1em\allowbreak{} 半正定。对任意\allowbreak{}\hskip.1em \RB{73.581}{8.184}{1.336}{%
\RG{0.000}{0.004}{10.9091}{mi}{122}\RG{8.478}{0.004}{10.9091}{rm}{61}\RG{19.817}{0.004}{10.9091}{mi}{117}\RG{27.903}{0.004}{10.9091}{rm}{43}\RG{38.218}{0.004}{10.9091}{mi}{105}\RG{41.862}{0.004}{10.9091}{mi}{118}\RG{50.448}{0.004}{10.9091}{sy}{50}\RG{60.751}{0.004}{10.9091}{bb}{67}\RG{68.630}{3.962}{7.9701}{mi}{110}%
}\hskip.1em\allowbreak{} ，其中\allowbreak{}\hskip.1em \RB{42.584}{7.977}{2.602}{%
\RG{0.000}{0.004}{10.9091}{mi}{117}\RG{6.076}{0.004}{10.9091}{mi}{59}\RG{10.855}{0.004}{10.9091}{mi}{118}\RG{19.451}{0.004}{10.9091}{sy}{50}\RG{29.754}{0.004}{10.9091}{bb}{82}\RG{37.633}{3.962}{7.9701}{mi}{110}%
}\hskip.1em\allowbreak{} ，对称性给出

\[
\RB{121.483}{10.003}{1.994}{%
\RG{5.448}{4.506}{7.9701}{sy}{3}\RG{42.867}{4.506}{7.9701}{rm}{84}\RG{80.597}{4.506}{7.9701}{rm}{84}\RG{0.000}{0.001}{10.9091}{mi}{122}\RG{9.767}{0.001}{10.9091}{mi}{83}\RG{16.967}{0.001}{10.9091}{mi}{122}\RG{25.449}{0.001}{10.9091}{rm}{61}\RG{36.788}{0.001}{10.9091}{mi}{117}\RG{48.600}{0.001}{10.9091}{mi}{83}\RG{55.800}{0.001}{10.9091}{mi}{117}\RG{64.308}{0.001}{10.9091}{rm}{43}\RG{75.040}{0.001}{10.9091}{mi}{118}\RG{86.329}{0.001}{10.9091}{mi}{83}\RG{93.529}{0.001}{10.9091}{mi}{118}\RG{102.119}{0.001}{10.9091}{sy}{21}\RG{113.634}{0.001}{10.9091}{rm}{48}\RG{118.975}{0.001}{10.9091}{mi}{58}%
}
\]
% Original page 6.
同理，当\allowbreak{}\hskip.1em \RB{25.201}{8.309}{2.858}{%
\RG{0.000}{-0.002}{10.9091}{mi}{122}\RG{8.479}{-0.002}{10.9091}{sy}{54}\RG{8.479}{-0.002}{10.9091}{rm}{61}\RG{19.814}{-0.002}{10.9091}{rm}{48}%
}\hskip.1em\allowbreak{} 时，\allowbreak{}\hskip.1em \RB{42.978}{8.287}{0.807}{%
\RG{0.000}{-0.002}{10.9091}{mi}{122}\RG{5.449}{3.957}{7.9701}{sy}{3}\RG{9.767}{-0.002}{10.9091}{mi}{65}\RG{17.774}{-0.002}{10.9091}{mi}{122}\RG{26.251}{-0.002}{10.9091}{mi}{62}\RG{37.591}{-0.002}{10.9091}{rm}{48}%
}\hskip.1em\allowbreak{} 。

设\allowbreak{}\hskip.1em \RB{40.231}{8.067}{0.624}{%
\RG{0.000}{-0.004}{10.9091}{mi}{66}\RG{8.662}{-0.004}{10.9091}{mi}{122}\RG{17.144}{-0.004}{10.9091}{rm}{61}\RG{28.483}{-0.004}{10.9091}{mi}{21}\RG{34.712}{-0.004}{10.9091}{mi}{122}%
}\hskip.1em\allowbreak{} ，\allowbreak{}\hskip.1em \RB{25.200}{8.307}{2.860}{%
\RG{0.000}{-0.004}{10.9091}{mi}{122}\RG{8.478}{-0.004}{10.9091}{sy}{54}\RG{8.478}{-0.004}{10.9091}{rm}{61}\RG{19.813}{-0.004}{10.9091}{rm}{48}%
}\hskip.1em\allowbreak{} 。于是

\[
\RB{112.858}{28.017}{21.259}{%
\RG{25.335}{22.521}{7.9701}{sy}{3}\RG{63.814}{22.521}{7.9701}{rm}{84}\RG{0.000}{18.016}{10.9091}{rm}{48}\RG{8.371}{18.016}{10.9091}{sy}{20}\RG{19.887}{18.016}{10.9091}{mi}{122}\RG{29.654}{18.016}{10.9091}{rm}{40}\RG{33.808}{18.016}{10.9091}{mi}{65}\RG{44.240}{18.016}{10.9091}{sy}{0}\RG{55.149}{18.016}{10.9091}{mi}{66}\RG{69.546}{18.016}{10.9091}{mi}{65}\RG{77.554}{18.016}{10.9091}{mi}{66}\RG{86.220}{18.016}{10.9091}{rm}{41}\RG{90.373}{18.016}{10.9091}{mi}{122}\RG{25.159}{4.502}{7.9701}{sy}{3}\RG{78.688}{4.502}{7.9701}{sy}{3}\RG{8.371}{-0.003}{10.9091}{rm}{61}\RG{19.710}{-0.003}{10.9091}{mi}{122}\RG{29.477}{-0.003}{10.9091}{mi}{65}\RG{37.484}{-0.003}{10.9091}{mi}{122}\RG{45.357}{-0.003}{10.9091}{sy}{0}\RG{56.266}{-0.003}{10.9091}{rm}{40}\RG{60.420}{-0.003}{10.9091}{mi}{66}\RG{69.082}{-0.003}{10.9091}{mi}{122}\RG{74.534}{-0.003}{10.9091}{rm}{41}\RG{83.006}{-0.003}{10.9091}{mi}{65}\RG{91.014}{-0.003}{10.9091}{rm}{40}\RG{95.168}{-0.003}{10.9091}{mi}{66}\RG{103.830}{-0.003}{10.9091}{mi}{122}\RG{109.282}{-0.003}{10.9091}{rm}{41}\RG{54.830}{-13.516}{7.9701}{rm}{50}\RG{68.751}{-13.516}{7.9701}{sy}{3}\RG{8.371}{-18.021}{10.9091}{rm}{61}\RG{19.706}{-18.021}{10.9091}{rm}{40}\RG{23.863}{-18.021}{10.9091}{rm}{49}\RG{31.630}{-18.021}{10.9091}{sy}{0}\RG{42.539}{-18.021}{10.9091}{sy}{106}\RG{45.569}{-18.021}{10.9091}{mi}{21}\RG{51.801}{-18.021}{10.9091}{sy}{106}\RG{59.149}{-18.021}{10.9091}{rm}{41}\RG{63.303}{-18.021}{10.9091}{mi}{122}\RG{73.069}{-18.021}{10.9091}{mi}{65}\RG{81.077}{-18.021}{10.9091}{mi}{122}\RG{86.521}{-18.021}{10.9091}{mi}{58}%
}
\]
由于\allowbreak{}\hskip.1em \RB{42.978}{8.288}{0.806}{%
\RG{0.000}{-0.001}{10.9091}{mi}{122}\RG{5.449}{3.958}{7.9701}{sy}{3}\RG{9.767}{-0.001}{10.9091}{mi}{65}\RG{17.774}{-0.001}{10.9091}{mi}{122}\RG{26.251}{-0.001}{10.9091}{mi}{62}\RG{37.591}{-0.001}{10.9091}{rm}{48}%
}\hskip.1em\allowbreak{} ，得\allowbreak{}\hskip.1em \RB{31.810}{8.681}{3.228}{%
\RG{0.000}{-0.001}{10.9091}{sy}{106}\RG{3.030}{-0.001}{10.9091}{mi}{21}\RG{9.261}{-0.001}{10.9091}{sy}{106}\RG{15.316}{-0.001}{10.9091}{sy}{20}\RG{26.837}{-0.001}{10.9091}{rm}{49}%
}\hskip.1em\allowbreak{} 。\allowbreak{}\hskip.1em \RB{8.366}{7.994}{0.501}{%
\RG{0.000}{-0.001}{10.9091}{am}{3}%
}\hskip.1em\allowbreak{}

\section*{2017S13}
设\allowbreak{}\hskip.1em \RB{21.391}{8.286}{2.608}{%
\RG{0.000}{-0.003}{10.9091}{mi}{65}\RG{8.007}{-0.003}{10.9091}{mi}{59}\RG{12.786}{-0.003}{10.9091}{mi}{66}%
}\hskip.1em\allowbreak{} 均为\allowbreak{}\hskip.1em \RB{6.555}{5.308}{0.612}{%
\RG{0.000}{-0.003}{10.9091}{mi}{110}%
}\hskip.1em\allowbreak{} 阶半正定实对称阵，且\allowbreak{}\hskip.1em \RB{54.781}{8.679}{3.241}{%
\RG{0.000}{-0.003}{10.9091}{rm}{116}\RG{4.156}{-0.003}{10.9091}{rm}{114}\RG{8.313}{-0.003}{10.9091}{rm}{40}\RG{12.461}{-0.003}{10.9091}{mi}{65}\RG{20.469}{-0.003}{10.9091}{mi}{66}\RG{29.135}{-0.003}{10.9091}{rm}{41}\RG{37.186}{-0.003}{10.9091}{rm}{61}\RG{49.393}{-0.003}{10.9091}{rm}{48}%
}\hskip.1em\allowbreak{} 。利用半正定阵的基本性质证明：

\[
\RB{47.792}{9.227}{2.682}{%
\RG{0.000}{0.000}{10.9091}{mi}{65}\RG{8.007}{0.000}{10.9091}{mi}{66}\RG{19.704}{0.000}{10.9091}{rm}{61}\RG{31.039}{0.000}{10.9091}{rm}{48}\RT{36.383}{0.000}{10.9091}{。}%
}
\]
证明：先证：若\allowbreak{}\hskip.1em \RB{8.606}{7.948}{0.503}{%
\RG{0.000}{-0.003}{10.9091}{mi}{66}%
}\hskip.1em\allowbreak{} 半正定且\allowbreak{}\hskip.1em \RB{46.309}{9.453}{0.732}{%
\RG{0.000}{-0.003}{10.9091}{mi}{117}\RG{6.079}{3.957}{7.9701}{rm}{84}\RG{11.812}{-0.003}{10.9091}{mi}{66}\RG{20.474}{-0.003}{10.9091}{mi}{117}\RG{29.586}{-0.003}{10.9091}{rm}{61}\RG{40.921}{-0.003}{10.9091}{rm}{48}%
}\hskip.1em\allowbreak{} ，则\allowbreak{}\hskip.1em \RB{34.497}{7.948}{0.732}{%
\RG{0.000}{-0.003}{10.9091}{mi}{66}\RG{8.662}{-0.003}{10.9091}{mi}{117}\RG{17.775}{-0.003}{10.9091}{rm}{61}\RG{29.110}{-0.003}{10.9091}{rm}{48}%
}\hskip.1em\allowbreak{} 。对任意实向量\allowbreak{}\hskip.1em \RB{5.715}{5.308}{2.728}{%
\RG{0.000}{-0.003}{10.9091}{mi}{121}%
}\hskip.1em\allowbreak{} 与实数\allowbreak{}\hskip.1em \RB{4.013}{7.316}{0.612}{%
\RG{0.000}{-0.003}{10.9091}{mi}{116}%
}\hskip.1em\allowbreak{} ，有

\[
\RB{209.180}{10.005}{3.234}{%
\RG{56.888}{4.509}{7.9701}{rm}{84}\RG{139.272}{4.509}{7.9701}{rm}{84}\RG{176.764}{4.509}{7.9701}{rm}{50}\RG{186.682}{4.509}{7.9701}{rm}{84}\RG{0.000}{0.005}{10.9091}{rm}{48}\RG{8.370}{0.005}{10.9091}{sy}{20}\RG{19.886}{0.005}{10.9091}{rm}{40}\RG{24.040}{0.005}{10.9091}{mi}{117}\RG{32.544}{0.005}{10.9091}{rm}{43}\RG{43.277}{0.005}{10.9091}{mi}{116}\RG{47.139}{0.005}{10.9091}{mi}{121}\RG{52.734}{0.005}{10.9091}{rm}{41}\RG{62.620}{0.005}{10.9091}{mi}{66}\RG{71.285}{0.005}{10.9091}{rm}{40}\RG{75.439}{0.005}{10.9091}{mi}{117}\RG{83.943}{0.005}{10.9091}{rm}{43}\RG{94.676}{0.005}{10.9091}{mi}{116}\RG{98.538}{0.005}{10.9091}{mi}{121}\RG{104.133}{0.005}{10.9091}{rm}{41}\RG{111.312}{0.005}{10.9091}{rm}{61}\RG{122.657}{0.005}{10.9091}{rm}{50}\RG{127.997}{0.005}{10.9091}{mi}{116}\RG{133.670}{0.005}{10.9091}{mi}{121}\RG{145.005}{0.005}{10.9091}{mi}{66}\RG{153.667}{0.005}{10.9091}{mi}{117}\RG{162.174}{0.005}{10.9091}{rm}{43}\RG{172.907}{0.005}{10.9091}{mi}{116}\RG{181.082}{0.005}{10.9091}{mi}{121}\RG{192.414}{0.005}{10.9091}{mi}{66}\RG{201.076}{0.005}{10.9091}{mi}{121}\RG{206.673}{0.005}{10.9091}{mi}{58}%
}
\]
若\allowbreak{}\hskip.1em \RB{45.829}{9.458}{2.855}{%
\RG{0.000}{0.002}{10.9091}{mi}{121}\RG{5.600}{3.961}{7.9701}{rm}{84}\RG{11.332}{0.002}{10.9091}{mi}{66}\RG{19.994}{0.002}{10.9091}{mi}{117}\RG{29.107}{0.002}{10.9091}{sy}{54}\RG{29.107}{0.002}{10.9091}{rm}{61}\RG{40.442}{0.002}{10.9091}{rm}{48}%
}\hskip.1em\allowbreak{} ，取绝对值充分小且与其符号相反的\allowbreak{}\hskip.1em \RB{4.013}{7.320}{0.607}{%
\RG{0.000}{0.002}{10.9091}{mi}{116}%
}\hskip.1em\allowbreak{} ，右端为负，产生矛盾。因此\allowbreak{}\hskip.1em \RB{45.829}{9.458}{2.724}{%
\RG{0.000}{0.002}{10.9091}{mi}{121}\RG{5.599}{3.961}{7.9701}{rm}{84}\RG{11.332}{0.002}{10.9091}{mi}{66}\RG{19.994}{0.002}{10.9091}{mi}{117}\RG{29.107}{0.002}{10.9091}{rm}{61}\RG{40.442}{0.002}{10.9091}{rm}{48}%
}\hskip.1em\allowbreak{}对所有\allowbreak{}\hskip.1em \RB{5.715}{5.316}{2.721}{%
\RG{0.000}{0.005}{10.9091}{mi}{121}%
}\hskip.1em\allowbreak{} 成立，故\allowbreak{}\hskip.1em \RB{34.498}{7.956}{0.725}{%
\RG{0.000}{0.005}{10.9091}{mi}{66}\RG{8.662}{0.005}{10.9091}{mi}{117}\RG{17.776}{0.005}{10.9091}{rm}{61}\RG{29.110}{0.005}{10.9091}{rm}{48}%
}\hskip.1em\allowbreak{} 。

由实对称阵的正交对角化定理（简证见文末），存在标准正交基\allowbreak{}\hskip.1em \RB{45.355}{5.314}{2.603}{%
\RG{0.000}{0.003}{10.9091}{mi}{117}\RG{6.080}{-1.633}{7.9701}{rm}{49}\RG{10.398}{0.003}{10.9091}{mi}{59}\RG{15.176}{0.003}{10.9091}{mi}{58}\RG{19.965}{0.003}{10.9091}{mi}{58}\RG{24.744}{0.003}{10.9091}{mi}{58}\RG{29.533}{0.003}{10.9091}{mi}{59}\RG{34.322}{0.003}{10.9091}{mi}{117}\RG{40.404}{-1.633}{7.9701}{mi}{110}%
}\hskip.1em\allowbreak{} ，使

\[
\RB{128.960}{15.527}{12.851}{%
\RG{26.741}{11.820}{7.9701}{mi}{110}\RG{22.519}{9.092}{9.0910}{ex}{88}\RG{64.000}{4.505}{7.9701}{rm}{84}\RG{0.000}{0.001}{10.9091}{mi}{65}\RG{11.038}{0.001}{10.9091}{rm}{61}\RG{37.612}{0.001}{10.9091}{mi}{21}\RG{47.842}{0.001}{10.9091}{mi}{117}\RG{57.921}{0.001}{10.9091}{mi}{117}\RG{69.732}{0.001}{10.9091}{mi}{59}\RG{96.329}{0.001}{10.9091}{mi}{21}\RG{109.596}{0.001}{10.9091}{sy}{21}\RG{121.111}{0.001}{10.9091}{rm}{48}\RG{126.452}{0.001}{10.9091}{mi}{58}\RG{43.843}{-1.635}{7.9701}{mi}{106}\RG{53.922}{-1.635}{7.9701}{mi}{106}\RG{102.567}{-1.635}{7.9701}{mi}{106}\RG{64.000}{-2.696}{7.9701}{mi}{106}\RG{22.376}{-10.867}{7.9701}{mi}{106}\RG{25.921}{-10.868}{7.9701}{rm}{61}\RG{31.928}{-10.868}{7.9701}{rm}{49}%
}
\]
其中\allowbreak{}\hskip.1em \RB{74.430}{9.455}{4.681}{%
\RG{0.000}{-0.000}{10.9091}{mi}{21}\RG{6.231}{-1.636}{7.9701}{mi}{106}\RG{13.260}{-0.000}{10.9091}{rm}{61}\RG{24.599}{-0.000}{10.9091}{mi}{117}\RG{30.679}{3.958}{7.9701}{rm}{84}\RG{30.679}{-2.697}{7.9701}{mi}{106}\RG{36.410}{-0.000}{10.9091}{mi}{65}\RG{44.418}{-0.000}{10.9091}{mi}{117}\RG{50.498}{-1.636}{7.9701}{mi}{106}\RG{57.527}{-0.000}{10.9091}{sy}{21}\RG{69.043}{-0.000}{10.9091}{rm}{48}%
}\hskip.1em\allowbreak{} 。直接按对角元求和可得

\[
\RB{94.604}{9.999}{4.683}{%
\RG{28.620}{4.502}{7.9701}{rm}{84}\RG{67.620}{4.502}{7.9701}{rm}{84}\RG{0.000}{-0.003}{10.9091}{rm}{116}\RG{4.156}{-0.003}{10.9091}{rm}{114}\RG{8.313}{-0.003}{10.9091}{rm}{40}\RG{12.462}{-0.002}{10.9091}{mi}{117}\RG{22.540}{-0.003}{10.9091}{mi}{117}\RG{34.352}{-0.003}{10.9091}{mi}{66}\RG{43.018}{-0.003}{10.9091}{rm}{41}\RG{50.196}{-0.003}{10.9091}{rm}{61}\RG{61.541}{-0.003}{10.9091}{mi}{117}\RG{73.352}{-0.003}{10.9091}{mi}{66}\RG{82.014}{-0.003}{10.9091}{mi}{117}\RG{92.096}{-0.003}{10.9091}{mi}{58}\RG{18.541}{-1.639}{7.9701}{mi}{106}\RG{88.098}{-1.639}{7.9701}{mi}{106}\RG{28.620}{-2.699}{7.9701}{mi}{106}\RG{67.620}{-2.699}{7.9701}{mi}{106}%
}
\]
因此

\[
\RB{125.897}{15.522}{12.855}{%
\RG{71.733}{11.815}{7.9701}{mi}{110}\RG{67.511}{9.088}{9.0910}{ex}{88}\RG{98.914}{4.501}{7.9701}{rm}{84}\RG{0.000}{-0.004}{10.9091}{rm}{48}\RG{8.367}{-0.004}{10.9091}{rm}{61}\RG{19.713}{-0.004}{10.9091}{rm}{116}\RG{23.869}{-0.004}{10.9091}{rm}{114}\RG{28.026}{-0.004}{10.9091}{rm}{40}\RG{32.172}{-0.004}{10.9091}{mi}{65}\RG{40.180}{-0.004}{10.9091}{mi}{66}\RG{48.845}{-0.004}{10.9091}{rm}{41}\RG{56.023}{-0.004}{10.9091}{rm}{61}\RG{82.604}{-0.004}{10.9091}{mi}{21}\RG{92.834}{-0.004}{10.9091}{mi}{117}\RG{104.646}{-0.004}{10.9091}{mi}{66}\RG{113.308}{-0.004}{10.9091}{mi}{117}\RG{123.390}{-0.004}{10.9091}{mi}{58}\RG{88.835}{-1.640}{7.9701}{mi}{106}\RG{119.391}{-1.640}{7.9701}{mi}{106}\RG{98.914}{-2.700}{7.9701}{mi}{106}\RG{67.368}{-10.872}{7.9701}{mi}{106}\RG{70.912}{-10.872}{7.9701}{rm}{61}\RG{76.920}{-10.872}{7.9701}{rm}{49}%
}
\]
每个加数均非负，故当\allowbreak{}\hskip.1em \RB{29.986}{8.076}{3.616}{%
\RG{0.000}{0.005}{10.9091}{mi}{21}\RG{6.232}{-1.633}{7.9701}{mi}{106}\RG{13.261}{0.005}{10.9091}{mi}{62}\RG{24.599}{0.005}{10.9091}{rm}{48}%
}\hskip.1em\allowbreak{} 时，\allowbreak{}\hskip.1em \RB{50.308}{9.459}{4.677}{%
\RG{0.000}{0.005}{10.9091}{mi}{117}\RG{6.079}{3.963}{7.9701}{rm}{84}\RG{6.079}{-2.693}{7.9701}{mi}{106}\RG{11.812}{0.005}{10.9091}{mi}{66}\RG{20.474}{0.005}{10.9091}{mi}{117}\RG{26.557}{-1.633}{7.9701}{mi}{106}\RG{33.586}{0.005}{10.9091}{rm}{61}\RG{44.921}{0.005}{10.9091}{rm}{48}%
}\hskip.1em\allowbreak{} ，由前述性质得\allowbreak{}\hskip.1em \RB{38.496}{7.956}{3.616}{%
\RG{0.000}{0.005}{10.9091}{mi}{66}\RG{8.662}{0.005}{10.9091}{mi}{117}\RG{14.745}{-1.633}{7.9701}{mi}{106}\RG{21.774}{0.005}{10.9091}{rm}{61}\RG{33.109}{0.005}{10.9091}{rm}{48}%
}\hskip.1em\allowbreak{} 。于是

\[
\RB{117.590}{15.524}{12.854}{%
\RG{35.407}{11.817}{7.9701}{mi}{110}\RG{31.186}{9.089}{9.0910}{ex}{88}\RG{89.640}{4.503}{7.9701}{rm}{84}\RG{0.000}{-0.001}{10.9091}{mi}{66}\RG{8.662}{-0.001}{10.9091}{mi}{65}\RG{19.704}{-0.001}{10.9091}{rm}{61}\RG{46.279}{-0.001}{10.9091}{mi}{21}\RG{56.508}{-0.001}{10.9091}{rm}{40}\RG{60.662}{-0.001}{10.9091}{mi}{66}\RG{69.324}{-0.001}{10.9091}{mi}{117}\RG{79.406}{-0.001}{10.9091}{rm}{41}\RG{83.561}{-0.001}{10.9091}{mi}{117}\RG{98.403}{-0.001}{10.9091}{rm}{61}\RG{109.737}{-0.001}{10.9091}{rm}{48}\RG{115.082}{-0.001}{10.9091}{mi}{58}\RG{52.510}{-1.638}{7.9701}{mi}{106}\RG{75.407}{-1.638}{7.9701}{mi}{106}\RG{89.640}{-2.698}{7.9701}{mi}{106}\RG{31.043}{-10.870}{7.9701}{mi}{106}\RG{34.587}{-10.870}{7.9701}{rm}{61}\RG{40.594}{-10.870}{7.9701}{rm}{49}%
}
\]
两边转置，并利用\allowbreak{}\hskip.1em \RB{21.391}{8.290}{2.605}{%
\RG{0.000}{0.001}{10.9091}{mi}{65}\RG{8.007}{0.001}{10.9091}{mi}{59}\RG{12.786}{0.001}{10.9091}{mi}{66}%
}\hskip.1em\allowbreak{} 的对称性，得到\allowbreak{}\hskip.1em \RB{36.425}{8.290}{0.728}{%
\RG{0.000}{0.001}{10.9091}{mi}{65}\RG{8.007}{0.001}{10.9091}{mi}{66}\RG{19.703}{0.001}{10.9091}{rm}{61}\RG{31.038}{0.001}{10.9091}{rm}{48}%
}\hskip.1em\allowbreak{} 。\allowbreak{}\hskip.1em \RB{8.366}{7.995}{0.500}{%
\RG{0.000}{0.001}{10.9091}{am}{3}%
}\hskip.1em\allowbreak{}

% Original page 7.
\section*{2017S14}
设\allowbreak{}\hskip.1em \RB{44.411}{5.312}{2.605}{%
\RG{0.000}{0.001}{10.9091}{mi}{97}\RG{5.607}{-1.636}{7.9701}{rm}{49}\RG{9.925}{0.001}{10.9091}{mi}{59}\RG{14.704}{0.001}{10.9091}{mi}{58}\RG{19.493}{0.001}{10.9091}{mi}{58}\RG{24.282}{0.001}{10.9091}{mi}{58}\RG{29.060}{0.001}{10.9091}{mi}{59}\RG{33.849}{0.001}{10.9091}{mi}{97}\RG{39.460}{-1.636}{7.9701}{mi}{110}%
}\hskip.1em\allowbreak{} 是互异的正实数。用两种方法证明实对称阵

\[
\RB{88.488}{28.558}{1.982}{%
\RG{22.376}{27.365}{9.0910}{ex}{32}\RG{63.211}{27.365}{9.0910}{ex}{33}\RG{70.408}{24.850}{7.9701}{mi}{110}\RG{43.721}{18.746}{10.9091}{rm}{49}\RG{0.000}{11.365}{10.9091}{mi}{65}\RG{11.038}{11.365}{10.9091}{rm}{61}\RG{30.664}{3.882}{10.9091}{mi}{97}\RG{41.781}{3.882}{10.9091}{rm}{43}\RG{52.513}{3.882}{10.9091}{mi}{97}\RG{36.271}{2.246}{7.9701}{mi}{105}\RG{58.121}{2.246}{7.9701}{mi}{106}\RG{70.408}{0.002}{7.9701}{mi}{105}\RG{73.040}{0.002}{7.9701}{mi}{59}\RG{75.186}{0.002}{7.9701}{mi}{106}\RG{78.730}{0.002}{7.9701}{rm}{61}\RG{84.737}{0.002}{7.9701}{rm}{49}\RR{30.664}{14.093}{31.455}{0.364}%
}
\]
为正定阵。证明：对任意\allowbreak{}\hskip.1em \RB{106.171}{9.454}{3.239}{%
\RG{0.000}{-0.001}{10.9091}{mi}{120}\RG{9.100}{-0.001}{10.9091}{rm}{61}\RG{20.435}{-0.001}{10.9091}{rm}{40}\RG{24.592}{-0.001}{10.9091}{mi}{120}\RG{30.663}{-1.637}{7.9701}{rm}{49}\RG{34.981}{-0.001}{10.9091}{mi}{59}\RG{39.759}{-0.001}{10.9091}{mi}{58}\RG{44.549}{-0.001}{10.9091}{mi}{58}\RG{49.327}{-0.001}{10.9091}{mi}{58}\RG{54.116}{-0.001}{10.9091}{mi}{59}\RG{58.905}{-0.001}{10.9091}{mi}{120}\RG{64.978}{-1.637}{7.9701}{mi}{110}\RG{70.121}{-0.001}{10.9091}{rm}{41}\RG{74.275}{3.958}{7.9701}{rm}{84}\RG{83.038}{-0.001}{10.9091}{sy}{50}\RG{93.341}{-0.001}{10.9091}{bb}{82}\RG{101.220}{3.958}{7.9701}{mi}{110}%
}\hskip.1em\allowbreak{} ，有

\[
\RB{158.791}{15.525}{46.442}{%
\RG{81.305}{12.828}{9.0910}{ex}{90}\RG{47.003}{11.818}{7.9701}{mi}{110}\RG{90.396}{9.949}{7.9701}{sy}{49}\RG{42.780}{9.090}{9.0910}{ex}{88}\RG{6.070}{4.504}{7.9701}{rm}{84}\RG{105.347}{4.504}{7.9701}{sy}{0}\RG{111.358}{4.504}{7.9701}{rm}{40}\RG{114.362}{4.504}{7.9701}{mi}{97}\RG{121.352}{4.504}{7.9701}{rm}{43}\RG{127.362}{4.504}{7.9701}{mi}{97}\RG{135.028}{4.504}{7.9701}{rm}{41}\RG{138.033}{4.504}{7.9701}{mi}{116}\RG{118.467}{3.395}{5.9776}{mi}{105}\RG{131.467}{3.395}{5.9776}{mi}{106}\RG{0.000}{-0.001}{10.9091}{mi}{120}\RG{11.802}{-0.001}{10.9091}{mi}{65}\RG{19.810}{-0.001}{10.9091}{mi}{120}\RG{28.910}{-0.001}{10.9091}{rm}{61}\RG{60.262}{-0.001}{10.9091}{mi}{120}\RG{69.418}{-0.001}{10.9091}{mi}{120}\RG{100.396}{-0.001}{10.9091}{mi}{101}\RG{143.097}{-0.001}{10.9091}{mi}{100}\RG{148.649}{-0.001}{10.9091}{mi}{116}\RG{66.332}{-1.637}{7.9701}{mi}{105}\RG{75.489}{-1.637}{7.9701}{mi}{106}\RG{86.356}{-7.829}{7.9701}{rm}{48}\RG{40.249}{-10.869}{7.9701}{mi}{105}\RG{42.882}{-10.869}{7.9701}{mi}{59}\RG{45.027}{-10.869}{7.9701}{mi}{106}\RG{48.570}{-10.869}{7.9701}{rm}{61}\RG{54.578}{-10.869}{7.9701}{rm}{49}\RG{59.340}{-19.000}{9.0910}{ex}{32}\RG{111.838}{-19.000}{9.0910}{ex}{33}\RG{119.035}{-21.516}{7.9701}{rm}{50}\RG{40.249}{-22.172}{9.0910}{ex}{90}\RG{70.759}{-23.182}{7.9701}{mi}{110}\RG{49.340}{-25.051}{7.9701}{sy}{49}\RG{66.538}{-25.910}{9.0910}{ex}{88}\RG{95.593}{-30.497}{7.9701}{sy}{0}\RG{101.604}{-30.497}{7.9701}{mi}{97}\RG{108.594}{-30.497}{7.9701}{mi}{116}\RG{105.708}{-31.606}{5.9776}{mi}{105}\RG{28.910}{-35.001}{10.9091}{rm}{61}\RG{81.487}{-35.001}{10.9091}{mi}{120}\RG{90.643}{-35.001}{10.9091}{mi}{101}\RG{126.990}{-35.001}{10.9091}{mi}{100}\RG{132.543}{-35.001}{10.9091}{mi}{116}\RG{139.428}{-35.001}{10.9091}{sy}{21}\RG{150.943}{-35.001}{10.9091}{rm}{48}\RG{156.284}{-35.001}{10.9091}{mi}{58}\RG{87.557}{-36.637}{7.9701}{mi}{105}\RG{45.300}{-42.830}{7.9701}{rm}{48}\RG{66.851}{-45.869}{7.9701}{mi}{105}\RG{69.481}{-45.869}{7.9701}{rm}{61}\RG{75.489}{-45.869}{7.9701}{rm}{49}%
}
\]
若上式等于零，由被积函数的非负性与连续性，得

\[
\RB{121.389}{15.523}{11.443}{%
\RG{4.222}{11.816}{7.9701}{mi}{110}\RG{0.000}{9.090}{9.0910}{ex}{88}\RG{29.056}{4.502}{7.9701}{sy}{0}\RG{35.067}{4.502}{7.9701}{mi}{97}\RG{42.057}{4.502}{7.9701}{mi}{116}\RG{39.171}{3.394}{5.9776}{mi}{105}\RG{14.950}{-0.002}{10.9091}{mi}{120}\RG{24.106}{-0.002}{10.9091}{mi}{101}\RG{48.332}{-0.002}{10.9091}{rm}{61}\RG{59.667}{-0.002}{10.9091}{rm}{48}\RG{86.831}{-0.002}{10.9091}{rm}{40}\RG{90.984}{-0.002}{10.9091}{mi}{116}\RG{97.872}{-0.002}{10.9091}{sy}{21}\RG{109.387}{-0.002}{10.9091}{rm}{48}\RG{114.732}{-0.002}{10.9091}{rm}{41}\RG{118.882}{-0.002}{10.9091}{mi}{58}\RG{21.021}{-1.638}{7.9701}{mi}{105}\RG{0.314}{-10.871}{7.9701}{mi}{105}\RG{2.945}{-10.871}{7.9701}{rm}{61}\RG{8.952}{-10.871}{7.9701}{rm}{49}%
}
\]
重新编号使\allowbreak{}\hskip.1em \RB{63.315}{6.246}{2.212}{%
\RG{0.000}{-0.003}{10.9091}{mi}{97}\RG{5.607}{-1.640}{7.9701}{rm}{49}\RG{13.297}{-0.003}{10.9091}{mi}{60}\RG{24.977}{-0.003}{10.9091}{sy}{1}\RG{29.821}{-0.003}{10.9091}{sy}{1}\RG{34.676}{-0.003}{10.9091}{sy}{1}\RG{41.077}{-0.003}{10.9091}{mi}{60}\RG{52.761}{-0.003}{10.9091}{mi}{97}\RG{58.364}{-1.640}{7.9701}{mi}{110}%
}\hskip.1em\allowbreak{} 。两边乘\allowbreak{}\hskip.1em \RB{15.889}{9.001}{0.612}{%
\RG{0.000}{-0.003}{10.9091}{mi}{101}\RG{4.951}{3.956}{7.9701}{mi}{97}\RG{9.056}{2.946}{5.9776}{rm}{49}\RG{12.844}{3.956}{7.9701}{mi}{116}%
}\hskip.1em\allowbreak{} 并令\allowbreak{}\hskip.1em \RB{32.297}{7.316}{0.623}{%
\RG{0.000}{-0.003}{10.9091}{mi}{116}\RG{7.229}{-0.003}{10.9091}{sy}{33}\RG{21.509}{-0.003}{10.9091}{sy}{49}%
}\hskip.1em\allowbreak{} ，得\allowbreak{}\hskip.1em \RB{30.832}{7.752}{2.140}{%
\RG{0.000}{-0.003}{10.9091}{mi}{120}\RG{6.070}{-1.640}{7.9701}{rm}{49}\RG{13.761}{-0.003}{10.9091}{rm}{61}\RG{25.445}{-0.003}{10.9091}{rm}{48}%
}\hskip.1em\allowbreak{} 。删去首项后重复此论证，依次得到\allowbreak{}\hskip.1em \RB{82.820}{7.755}{2.210}{%
\RG{0.000}{0.000}{10.9091}{mi}{120}\RG{6.070}{-1.637}{7.9701}{rm}{50}\RG{13.419}{0.000}{10.9091}{rm}{61}\RG{24.757}{0.000}{10.9091}{sy}{1}\RG{29.601}{0.000}{10.9091}{sy}{1}\RG{34.456}{0.000}{10.9091}{sy}{1}\RG{40.516}{0.000}{10.9091}{rm}{61}\RG{51.854}{0.000}{10.9091}{mi}{120}\RG{57.924}{-1.637}{7.9701}{mi}{110}\RG{66.098}{0.000}{10.9091}{rm}{61}\RG{77.433}{0.000}{10.9091}{rm}{48}%
}\hskip.1em\allowbreak{} 。因此\allowbreak{}\hskip.1em \RB{25.823}{8.311}{2.856}{%
\RG{0.000}{0.000}{10.9091}{mi}{120}\RG{9.100}{0.000}{10.9091}{sy}{54}\RG{9.100}{0.000}{10.9091}{rm}{61}\RG{20.435}{0.000}{10.9091}{rm}{48}%
}\hskip.1em\allowbreak{} 时\allowbreak{}\hskip.1em \RB{45.636}{9.455}{0.805}{%
\RG{0.000}{0.000}{10.9091}{mi}{120}\RG{6.070}{3.959}{7.9701}{rm}{84}\RG{11.802}{0.000}{10.9091}{mi}{65}\RG{19.810}{0.000}{10.9091}{mi}{120}\RG{28.897}{0.000}{10.9091}{mi}{62}\RG{40.249}{0.000}{10.9091}{rm}{48}%
}\hskip.1em\allowbreak{} ，即\allowbreak{}\hskip.1em \RB{8.202}{8.289}{0.500}{%
\RG{0.000}{0.000}{10.9091}{mi}{65}%
}\hskip.1em\allowbreak{} 正定。\allowbreak{}\hskip.1em \RB{8.366}{7.995}{0.500}{%
\RG{0.000}{0.000}{10.9091}{am}{3}%
}\hskip.1em\allowbreak{}另一证明：对\allowbreak{}\hskip.1em \RB{6.555}{5.308}{0.612}{%
\RG{0.000}{-0.003}{10.9091}{mi}{110}%
}\hskip.1em\allowbreak{} 归纳。\allowbreak{}\hskip.1em \RB{25.714}{7.752}{0.612}{%
\RG{0.000}{-0.003}{10.9091}{mi}{110}\RG{9.406}{-0.003}{10.9091}{rm}{61}\RG{20.741}{-0.003}{10.9091}{rm}{49}%
}\hskip.1em\allowbreak{} 时，\allowbreak{}\hskip.1em \RB{64.366}{8.679}{3.241}{%
\RG{0.000}{-0.003}{10.9091}{mi}{65}\RG{11.038}{-0.003}{10.9091}{rm}{61}\RG{22.372}{-0.003}{10.9091}{rm}{40}\RG{26.529}{-0.003}{10.9091}{rm}{49}\RG{31.871}{-0.003}{10.9091}{mi}{61}\RG{37.212}{-0.003}{10.9091}{rm}{40}\RG{41.369}{-0.003}{10.9091}{rm}{50}\RG{46.708}{-0.003}{10.9091}{mi}{97}\RG{52.315}{-1.639}{7.9701}{rm}{49}\RG{56.633}{-0.003}{10.9091}{rm}{41}\RG{60.790}{-0.003}{10.9091}{rm}{41}%
}\hskip.1em\allowbreak{} 正定。设结论对\allowbreak{}\hskip.1em \RB{24.419}{7.752}{0.612}{%
\RG{0.000}{-0.003}{10.9091}{mi}{110}\RG{8.669}{-0.003}{10.9091}{sy}{0}\RG{19.447}{-0.003}{10.9091}{rm}{49}%
}\hskip.1em\allowbreak{} 成立，将\allowbreak{}\hskip.1em \RB{8.202}{8.287}{0.503}{%
\RG{0.000}{-0.003}{10.9091}{mi}{65}%
}\hskip.1em\allowbreak{} 分块为

\[
\RB{337.838}{19.588}{14.122}{%
\RG{22.377}{18.733}{9.0910}{ex}{48}\RG{56.560}{18.733}{9.0910}{ex}{49}\RG{266.831}{16.005}{9.0910}{ex}{32}\RG{307.665}{16.005}{9.0910}{ex}{33}\RG{314.863}{13.490}{7.9701}{mi}{110}\RG{166.222}{13.278}{9.0910}{ex}{18}\RG{206.872}{13.278}{9.0910}{ex}{19}\RG{213.564}{11.376}{7.9701}{mi}{110}\RG{50.828}{11.279}{7.9701}{rm}{84}\RG{120.223}{7.385}{10.9091}{rm}{49}\RG{187.223}{7.385}{10.9091}{rm}{49}\RG{288.176}{7.385}{10.9091}{rm}{49}\RG{30.331}{7.319}{10.9091}{mi}{11}\RG{46.280}{7.319}{10.9091}{mi}{98}\RG{22.377}{2.369}{9.0910}{ex}{64}\RG{56.560}{2.369}{9.0910}{ex}{65}\RG{0.000}{0.005}{10.9091}{mi}{65}\RG{11.038}{0.005}{10.9091}{rm}{61}\RG{66.332}{0.005}{10.9091}{mi}{59}\RG{92.929}{0.005}{10.9091}{mi}{11}\RG{102.830}{0.005}{10.9091}{rm}{61}\RG{131.617}{0.005}{10.9091}{mi}{59}\RG{147.305}{0.005}{10.9091}{mi}{98}\RG{154.884}{0.005}{10.9091}{rm}{61}\RG{228.341}{0.005}{10.9091}{mi}{59}\RG{244.028}{0.005}{10.9091}{mi}{67}\RG{255.492}{0.005}{10.9091}{rm}{61}\RG{335.330}{0.005}{10.9091}{mi}{58}\RG{115.260}{-7.478}{10.9091}{rm}{50}\RG{120.601}{-7.478}{10.9091}{mi}{97}\RG{174.005}{-7.478}{10.9091}{mi}{97}\RG{186.355}{-7.478}{10.9091}{rm}{43}\RG{197.088}{-7.478}{10.9091}{mi}{97}\RG{275.118}{-7.478}{10.9091}{mi}{97}\RG{286.235}{-7.478}{10.9091}{rm}{43}\RG{296.968}{-7.478}{10.9091}{mi}{97}\RG{31.492}{-7.972}{10.9091}{mi}{98}\RG{47.201}{-7.972}{10.9091}{mi}{67}\RG{126.209}{-9.115}{7.9701}{rm}{49}\RG{179.613}{-9.115}{7.9701}{rm}{49}\RG{202.696}{-9.115}{7.9701}{mi}{105}\RG{280.725}{-9.115}{7.9701}{mi}{105}\RG{302.576}{-9.115}{7.9701}{mi}{106}\RG{213.564}{-9.570}{7.9701}{mi}{105}\RG{216.194}{-9.570}{7.9701}{rm}{61}\RG{222.202}{-9.570}{7.9701}{rm}{50}\RG{314.863}{-11.359}{7.9701}{mi}{105}\RG{317.495}{-11.359}{7.9701}{mi}{59}\RG{319.641}{-11.359}{7.9701}{mi}{106}\RG{323.184}{-11.359}{7.9701}{rm}{61}\RG{329.191}{-11.359}{7.9701}{rm}{50}\RR{115.260}{2.732}{15.266}{0.364}\RR{174.005}{2.732}{31.776}{0.364}\RR{275.118}{2.732}{31.456}{0.364}%
}
\]
归纳假设给出\allowbreak{}\hskip.1em \RB{8.649}{8.176}{0.734}{%
\RG{0.000}{-0.005}{10.9091}{mi}{67}%
}\hskip.1em\allowbreak{} 正定。对于\allowbreak{}\hskip.1em \RB{54.449}{7.750}{2.730}{%
\RG{0.000}{-0.005}{10.9091}{rm}{50}\RG{8.371}{-0.005}{10.9091}{sy}{20}\RG{19.886}{-0.005}{10.9091}{mi}{105}\RG{23.530}{-0.005}{10.9091}{mi}{59}\RG{28.308}{-0.005}{10.9091}{mi}{106}\RG{36.379}{-0.005}{10.9091}{sy}{20}\RG{47.894}{-0.005}{10.9091}{mi}{110}%
}\hskip.1em\allowbreak{} ，

\[
\RB{269.043}{23.550}{3.616}{%
\RG{13.057}{14.868}{10.9091}{rm}{49}\RG{79.878}{14.868}{10.9091}{rm}{50}\RG{85.219}{14.868}{10.9091}{mi}{97}\RG{164.306}{14.868}{10.9091}{rm}{40}\RG{168.461}{14.868}{10.9091}{mi}{97}\RG{179.577}{14.868}{10.9091}{sy}{0}\RG{190.487}{14.868}{10.9091}{mi}{97}\RG{200.412}{14.868}{10.9091}{rm}{41}\RG{204.569}{14.868}{10.9091}{rm}{40}\RG{208.720}{14.868}{10.9091}{mi}{97}\RG{220.751}{14.868}{10.9091}{sy}{0}\RG{231.660}{14.868}{10.9091}{mi}{97}\RG{241.585}{14.868}{10.9091}{rm}{41}\RG{90.826}{13.231}{7.9701}{rm}{49}\RG{174.067}{13.231}{7.9701}{mi}{105}\RG{196.094}{13.231}{7.9701}{rm}{49}\RG{214.327}{13.231}{7.9701}{mi}{106}\RG{237.267}{13.231}{7.9701}{rm}{49}\RG{34.971}{7.487}{10.9091}{sy}{0}\RG{132.172}{7.487}{10.9091}{rm}{61}\RG{266.536}{7.487}{10.9091}{mi}{58}\RG{0.000}{0.004}{10.9091}{mi}{97}\RG{11.117}{0.004}{10.9091}{rm}{43}\RG{21.849}{0.004}{10.9091}{mi}{97}\RG{46.971}{0.004}{10.9091}{rm}{40}\RG{51.124}{0.004}{10.9091}{mi}{97}\RG{62.242}{0.004}{10.9091}{rm}{43}\RG{72.975}{0.004}{10.9091}{mi}{97}\RG{82.900}{0.004}{10.9091}{rm}{41}\RG{87.056}{0.004}{10.9091}{rm}{40}\RG{91.208}{0.004}{10.9091}{mi}{97}\RG{103.238}{0.004}{10.9091}{rm}{43}\RG{113.971}{0.004}{10.9091}{mi}{97}\RG{123.896}{0.004}{10.9091}{rm}{41}\RG{144.601}{0.004}{10.9091}{rm}{40}\RG{148.755}{0.004}{10.9091}{mi}{97}\RG{159.872}{0.004}{10.9091}{rm}{43}\RG{170.605}{0.004}{10.9091}{mi}{97}\RG{180.211}{0.004}{10.9091}{rm}{41}\RG{184.368}{0.004}{10.9091}{rm}{40}\RG{188.519}{0.004}{10.9091}{mi}{97}\RG{199.636}{0.004}{10.9091}{rm}{43}\RG{210.369}{0.004}{10.9091}{mi}{97}\RG{220.294}{0.004}{10.9091}{rm}{41}\RG{224.451}{0.004}{10.9091}{rm}{40}\RG{228.602}{0.004}{10.9091}{mi}{97}\RG{240.632}{0.004}{10.9091}{rm}{43}\RG{251.365}{0.004}{10.9091}{mi}{97}\RG{261.291}{0.004}{10.9091}{rm}{41}\RG{5.607}{-1.633}{7.9701}{mi}{105}\RG{27.457}{-1.633}{7.9701}{mi}{106}\RG{56.732}{-1.633}{7.9701}{mi}{105}\RG{78.582}{-1.633}{7.9701}{rm}{49}\RG{96.816}{-1.633}{7.9701}{mi}{106}\RG{119.578}{-1.633}{7.9701}{rm}{49}\RG{154.362}{-1.633}{7.9701}{mi}{105}\RG{176.212}{-1.633}{7.9701}{mi}{106}\RG{194.126}{-1.633}{7.9701}{mi}{105}\RG{215.976}{-1.633}{7.9701}{rm}{49}\RG{234.209}{-1.633}{7.9701}{mi}{106}\RG{256.973}{-1.633}{7.9701}{rm}{49}\RR{0.000}{10.215}{31.455}{0.364}\RR{46.971}{10.215}{81.079}{0.364}\RR{144.601}{10.215}{120.843}{0.364}%
}
\]
令

\[
\RB{158.275}{14.132}{9.697}{%
\RG{45.081}{13.268}{9.0910}{ex}{18}\RG{147.258}{13.268}{9.0910}{ex}{19}\RG{52.864}{7.375}{10.9091}{mi}{97}\RG{65.214}{7.375}{10.9091}{sy}{0}\RG{76.123}{7.375}{10.9091}{mi}{97}\RG{112.157}{7.375}{10.9091}{mi}{97}\RG{125.332}{7.375}{10.9091}{sy}{0}\RG{136.241}{7.375}{10.9091}{mi}{97}\RG{58.471}{5.739}{7.9701}{rm}{50}\RG{81.730}{5.739}{7.9701}{rm}{49}\RG{117.764}{5.739}{7.9701}{mi}{110}\RG{141.848}{5.739}{7.9701}{rm}{49}\RG{0.000}{-0.005}{10.9091}{mi}{68}\RG{12.196}{-0.005}{10.9091}{rm}{61}\RG{23.530}{-0.005}{10.9091}{rm}{100}\RG{29.465}{-0.005}{10.9091}{rm}{105}\RG{32.432}{-0.005}{10.9091}{rm}{97}\RG{37.778}{-0.005}{10.9091}{rm}{103}\RG{87.140}{-0.005}{10.9091}{mi}{59}\RG{91.918}{-0.005}{10.9091}{mi}{58}\RG{96.707}{-0.005}{10.9091}{mi}{58}\RG{101.485}{-0.005}{10.9091}{mi}{58}\RG{106.274}{-0.005}{10.9091}{mi}{59}\RG{155.768}{-0.005}{10.9091}{mi}{58}\RG{52.952}{-7.488}{10.9091}{mi}{97}\RG{65.302}{-7.488}{10.9091}{rm}{43}\RG{76.034}{-7.488}{10.9091}{mi}{97}\RG{112.246}{-7.488}{10.9091}{mi}{97}\RG{125.420}{-7.488}{10.9091}{rm}{43}\RG{136.152}{-7.488}{10.9091}{mi}{97}\RG{58.560}{-9.124}{7.9701}{rm}{50}\RG{81.642}{-9.124}{7.9701}{rm}{49}\RG{117.853}{-9.124}{7.9701}{mi}{110}\RG{141.760}{-9.124}{7.9701}{rm}{49}\RR{52.864}{2.722}{33.184}{0.364}\RR{112.157}{2.722}{34.009}{0.364}%
}
\]
则\allowbreak{}\hskip.1em \RB{91.198}{9.451}{0.733}{%
\RG{0.000}{-0.004}{10.9091}{mi}{67}\RG{9.028}{-0.004}{10.9091}{sy}{0}\RG{18.117}{-0.004}{10.9091}{mi}{11}\RG{24.981}{3.955}{7.9701}{sy}{0}\RG{30.991}{3.955}{7.9701}{rm}{49}\RG{35.309}{-0.004}{10.9091}{mi}{98}\RG{39.847}{-0.004}{10.9091}{mi}{98}\RG{44.393}{3.955}{7.9701}{rm}{84}\RG{53.155}{-0.004}{10.9091}{rm}{61}\RG{64.494}{-0.004}{10.9091}{mi}{68}\RG{73.657}{-0.004}{10.9091}{mi}{67}\RG{82.090}{-0.004}{10.9091}{mi}{68}%
}\hskip.1em\allowbreak{} 。\allowbreak{}\hskip.1em \RB{9.107}{7.947}{0.504}{%
\RG{0.000}{-0.004}{10.9091}{mi}{68}%
}\hskip.1em\allowbreak{} 的所有对角元均非零，因而\allowbreak{}\hskip.1em \RB{9.107}{7.947}{0.504}{%
\RG{0.000}{-0.004}{10.9091}{mi}{68}%
}\hskip.1em\allowbreak{} 可逆，\allowbreak{}\hskip.1em \RB{26.704}{8.176}{0.733}{%
\RG{0.000}{-0.004}{10.9091}{mi}{68}\RG{9.164}{-0.004}{10.9091}{mi}{67}\RG{17.597}{-0.004}{10.9091}{mi}{68}%
}\hskip.1em\allowbreak{} 正定。对\allowbreak{}\hskip.1em \RB{60.749}{9.451}{3.242}{%
\RG{0.000}{-0.004}{10.9091}{mi}{120}\RG{9.100}{-0.004}{10.9091}{rm}{61}\RG{20.435}{-0.004}{10.9091}{rm}{40}\RG{24.593}{-0.004}{10.9091}{mi}{120}\RG{30.663}{-1.641}{7.9701}{rm}{49}\RG{34.981}{-0.004}{10.9091}{mi}{59}\RG{39.759}{-0.004}{10.9091}{mi}{121}\RG{45.366}{3.955}{7.9701}{rm}{84}\RG{51.099}{-0.004}{10.9091}{rm}{41}\RG{55.253}{3.955}{7.9701}{rm}{84}%
}\hskip.1em\allowbreak{} ，配方得

\[
\RB{211.180}{13.363}{27.758}{%
\RG{47.112}{10.542}{9.0910}{ex}{16}\RG{109.153}{10.542}{9.0910}{ex}{17}\RG{114.582}{8.026}{7.9701}{rm}{50}\RG{6.070}{4.500}{7.9701}{rm}{84}\RG{82.951}{4.500}{7.9701}{sy}{0}\RG{88.961}{4.500}{7.9701}{rm}{49}\RG{97.820}{4.500}{7.9701}{rm}{84}\RG{137.657}{4.500}{7.9701}{rm}{84}\RG{176.166}{4.500}{7.9701}{sy}{0}\RG{182.177}{4.500}{7.9701}{rm}{49}\RG{195.578}{4.500}{7.9701}{rm}{84}\RG{0.000}{-0.004}{10.9091}{mi}{120}\RG{11.802}{-0.004}{10.9091}{mi}{65}\RG{19.810}{-0.004}{10.9091}{mi}{120}\RG{28.910}{-0.004}{10.9091}{rm}{61}\RG{40.249}{-0.004}{10.9091}{mi}{11}\RG{52.542}{-0.004}{10.9091}{mi}{120}\RG{65.354}{-0.004}{10.9091}{rm}{43}\RG{76.087}{-0.004}{10.9091}{mi}{11}\RG{93.279}{-0.004}{10.9091}{mi}{98}\RG{103.553}{-0.004}{10.9091}{mi}{121}\RG{121.325}{-0.004}{10.9091}{rm}{43}\RG{132.057}{-0.004}{10.9091}{mi}{121}\RG{143.389}{-0.004}{10.9091}{rm}{40}\RG{147.543}{-0.004}{10.9091}{mi}{67}\RG{158.393}{-0.004}{10.9091}{sy}{0}\RG{169.303}{-0.004}{10.9091}{mi}{11}\RG{186.495}{-0.004}{10.9091}{mi}{98}\RG{191.033}{-0.004}{10.9091}{mi}{98}\RG{201.310}{-0.004}{10.9091}{rm}{41}\RG{205.465}{-0.004}{10.9091}{mi}{121}\RG{58.612}{-1.640}{7.9701}{rm}{49}\RG{47.112}{-11.267}{9.0910}{ex}{16}\RG{109.153}{-11.267}{9.0910}{ex}{17}\RG{114.582}{-13.782}{7.9701}{rm}{50}\RG{82.951}{-17.308}{7.9701}{sy}{0}\RG{88.961}{-17.308}{7.9701}{rm}{49}\RG{97.820}{-17.308}{7.9701}{rm}{84}\RG{155.131}{-17.308}{7.9701}{rm}{84}\RG{28.910}{-21.812}{10.9091}{rm}{61}\RG{40.249}{-21.812}{10.9091}{mi}{11}\RG{52.542}{-21.812}{10.9091}{mi}{120}\RG{65.354}{-21.812}{10.9091}{rm}{43}\RG{76.087}{-21.812}{10.9091}{mi}{11}\RG{93.279}{-21.812}{10.9091}{mi}{98}\RG{103.553}{-21.812}{10.9091}{mi}{121}\RG{121.325}{-21.812}{10.9091}{rm}{43}\RG{132.060}{-21.812}{10.9091}{rm}{40}\RG{136.212}{-21.812}{10.9091}{mi}{68}\RG{145.375}{-21.812}{10.9091}{mi}{121}\RG{150.977}{-21.812}{10.9091}{rm}{41}\RG{160.863}{-21.812}{10.9091}{mi}{67}\RG{169.289}{-21.812}{10.9091}{rm}{40}\RG{173.443}{-21.812}{10.9091}{mi}{68}\RG{182.607}{-21.812}{10.9091}{mi}{121}\RG{188.208}{-21.812}{10.9091}{rm}{41}\RG{192.362}{-21.812}{10.9091}{mi}{58}\RG{58.612}{-23.449}{7.9701}{rm}{49}%
}
\]
若\allowbreak{}\hskip.1em \RB{25.352}{8.315}{2.853}{%
\RG{0.000}{0.004}{10.9091}{mi}{121}\RG{8.630}{0.004}{10.9091}{sy}{54}\RG{8.630}{0.004}{10.9091}{rm}{61}\RG{19.964}{0.004}{10.9091}{rm}{48}%
}\hskip.1em\allowbreak{} ，第二项为正；若\allowbreak{}\hskip.1em \RB{25.352}{7.758}{2.722}{%
\RG{0.000}{0.004}{10.9091}{mi}{121}\RG{8.630}{0.004}{10.9091}{rm}{61}\RG{19.964}{0.004}{10.9091}{rm}{48}%
}\hskip.1em\allowbreak{} 且\allowbreak{}\hskip.1em \RB{25.823}{8.315}{2.853}{%
\RG{0.000}{0.004}{10.9091}{mi}{120}\RG{9.100}{0.004}{10.9091}{sy}{54}\RG{9.100}{0.004}{10.9091}{rm}{61}\RG{20.435}{0.004}{10.9091}{rm}{48}%
}\hskip.1em\allowbreak{} ，第一项为正。因此\allowbreak{}\hskip.1em \RB{8.202}{8.293}{0.500}{%
\RG{0.000}{0.004}{10.9091}{mi}{65}%
}\hskip.1em\allowbreak{} 正定，归纳完成。\allowbreak{}\hskip.1em \RB{8.366}{7.998}{0.500}{%
\RG{0.000}{0.004}{10.9091}{am}{3}%
}\hskip.1em\allowbreak{}

% Original page 8.
\section*{2017S15}
设\allowbreak{}\hskip.1em \RB{8.202}{8.290}{0.500}{%
\RG{0.000}{0.001}{10.9091}{mi}{65}%
}\hskip.1em\allowbreak{} 为\allowbreak{}\hskip.1em \RB{6.555}{5.312}{0.608}{%
\RG{0.000}{0.001}{10.9091}{mi}{110}%
}\hskip.1em\allowbreak{} 阶正定实对称阵，\allowbreak{}\hskip.1em \RB{79.772}{9.456}{3.237}{%
\RG{0.000}{0.001}{10.9091}{mi}{120}\RG{9.100}{0.001}{10.9091}{rm}{61}\RG{20.435}{0.001}{10.9091}{rm}{40}\RG{24.593}{0.001}{10.9091}{mi}{120}\RG{30.663}{-1.636}{7.9701}{rm}{49}\RG{34.981}{0.001}{10.9091}{mi}{59}\RG{39.759}{0.001}{10.9091}{mi}{58}\RG{44.549}{0.001}{10.9091}{mi}{58}\RG{49.338}{0.001}{10.9091}{mi}{58}\RG{54.116}{0.001}{10.9091}{mi}{59}\RG{58.905}{0.001}{10.9091}{mi}{120}\RG{64.978}{-1.636}{7.9701}{mi}{110}\RG{70.121}{0.001}{10.9091}{rm}{41}\RG{74.276}{3.959}{7.9701}{rm}{84}%
}\hskip.1em\allowbreak{} ，\allowbreak{}\hskip.1em \RB{61.095}{9.456}{3.237}{%
\RG{0.000}{0.001}{10.9091}{mi}{102}\RG{6.430}{0.001}{10.9091}{rm}{40}\RG{10.584}{0.001}{10.9091}{mi}{120}\RG{16.654}{0.001}{10.9091}{rm}{41}\RG{23.833}{0.001}{10.9091}{rm}{61}\RG{35.177}{0.001}{10.9091}{mi}{120}\RG{41.247}{3.959}{7.9701}{rm}{84}\RG{46.980}{0.001}{10.9091}{mi}{65}\RG{54.987}{0.001}{10.9091}{mi}{120}%
}\hskip.1em\allowbreak{} 。删去\allowbreak{}\hskip.1em \RB{8.202}{8.290}{0.500}{%
\RG{0.000}{0.001}{10.9091}{mi}{65}%
}\hskip.1em\allowbreak{} 的第\allowbreak{}\hskip.1em \RB{3.587}{7.668}{0.608}{%
\RG{0.000}{0.001}{10.9091}{mi}{105}%
}\hskip.1em\allowbreak{} 行与第\allowbreak{}\hskip.1em \RB{3.587}{7.668}{0.608}{%
\RG{0.000}{0.001}{10.9091}{mi}{105}%
}\hskip.1em\allowbreak{}列所得的主子阵记为\allowbreak{}\hskip.1em \RB{10.813}{8.292}{2.206}{%
\RG{0.000}{0.003}{10.9091}{mi}{65}\RG{8.007}{-1.633}{7.9701}{mi}{105}%
}\hskip.1em\allowbreak{} 。证明：在约束\allowbreak{}\hskip.1em \RB{28.494}{7.757}{2.206}{%
\RG{0.000}{0.003}{10.9091}{mi}{120}\RG{6.070}{-1.633}{7.9701}{mi}{105}\RG{12.187}{0.003}{10.9091}{rm}{61}\RG{23.521}{0.003}{10.9091}{rm}{49}%
}\hskip.1em\allowbreak{} 下，\allowbreak{}\hskip.1em \RB{20.231}{8.685}{3.235}{%
\RG{0.000}{0.003}{10.9091}{mi}{102}\RG{6.430}{0.003}{10.9091}{rm}{40}\RG{10.584}{0.003}{10.9091}{mi}{120}\RG{16.654}{0.003}{10.9091}{rm}{41}%
}\hskip.1em\allowbreak{} 的最小值为

\[
\RB{102.407}{15.672}{9.690}{%
\RG{1.543}{7.383}{10.9091}{rm}{100}\RG{7.478}{7.383}{10.9091}{rm}{101}\RG{12.223}{7.383}{10.9091}{rm}{116}\RG{18.196}{7.383}{10.9091}{mi}{65}\RG{28.838}{0.003}{10.9091}{mi}{59}\RG{55.442}{0.003}{10.9091}{rm}{49}\RG{63.813}{0.003}{10.9091}{sy}{20}\RG{75.328}{0.003}{10.9091}{mi}{105}\RG{82.003}{0.003}{10.9091}{sy}{20}\RG{93.518}{0.003}{10.9091}{mi}{110}\RG{99.900}{0.003}{10.9091}{mi}{58}\RG{0.000}{-7.481}{10.9091}{rm}{100}\RG{5.935}{-7.481}{10.9091}{rm}{101}\RG{10.680}{-7.481}{10.9091}{rm}{116}\RG{16.654}{-7.481}{10.9091}{mi}{65}\RG{24.662}{-9.117}{7.9701}{mi}{105}\RR{0.000}{2.730}{27.747}{0.364}%
}
\]
证明：先设\allowbreak{}\hskip.1em \RB{26.222}{7.752}{1.997}{%
\RG{0.000}{-0.003}{10.9091}{mi}{110}\RG{9.407}{-0.003}{10.9091}{sy}{21}\RG{20.922}{-0.003}{10.9091}{rm}{50}%
}\hskip.1em\allowbreak{} 。同步置换行列及坐标，将第\allowbreak{}\hskip.1em \RB{3.587}{7.665}{0.612}{%
\RG{0.000}{-0.003}{10.9091}{mi}{105}%
}\hskip.1em\allowbreak{} 个坐标移至末位，可写为

\[
\RB{138.073}{19.579}{14.130}{%
\RG{22.377}{18.725}{9.0910}{ex}{48}\RG{55.303}{18.725}{9.0910}{ex}{49}\RG{112.119}{18.725}{9.0910}{ex}{48}\RG{125.673}{18.725}{9.0910}{ex}{49}\RG{31.255}{7.312}{10.9091}{mi}{67}\RG{50.227}{7.312}{10.9091}{mi}{98}\RG{120.074}{7.312}{10.9091}{mi}{121}\RG{2.962}{3.538}{9.0910}{ex}{101}\RG{22.377}{2.361}{9.0910}{ex}{64}\RG{55.304}{2.361}{9.0910}{ex}{65}\RG{112.119}{2.361}{9.0910}{ex}{64}\RG{125.673}{2.361}{9.0910}{ex}{65}\RG{92.487}{0.780}{9.0910}{ex}{101}\RG{0.000}{-0.003}{10.9091}{mi}{65}\RG{11.038}{-0.003}{10.9091}{rm}{61}\RG{65.076}{-0.003}{10.9091}{mi}{59}\RG{91.680}{-0.003}{10.9091}{mi}{120}\RG{100.781}{-0.003}{10.9091}{rm}{61}\RG{135.446}{-0.003}{10.9091}{mi}{59}\RG{34.873}{-4.021}{7.9701}{rm}{84}\RG{30.331}{-7.979}{10.9091}{mi}{98}\RG{49.696}{-7.979}{10.9091}{mi}{97}\RG{120.203}{-7.979}{10.9091}{rm}{49}%
}
\]
其中其余坐标保持原有次序，故\allowbreak{}\hskip.1em \RB{33.607}{8.294}{2.205}{%
\RG{0.000}{0.005}{10.9091}{mi}{67}\RG{11.456}{0.005}{10.9091}{rm}{61}\RG{22.795}{0.005}{10.9091}{mi}{65}\RG{30.802}{-1.633}{7.9701}{mi}{105}%
}\hskip.1em\allowbreak{} 、\allowbreak{}\hskip.1em \RB{31.021}{5.316}{2.205}{%
\RG{0.000}{0.005}{10.9091}{mi}{97}\RG{8.638}{0.005}{10.9091}{rm}{61}\RG{19.976}{0.005}{10.9091}{mi}{97}\RG{25.583}{-1.633}{7.9701}{mi}{105}\RG{28.216}{-1.633}{7.9701}{mi}{105}%
}\hskip.1em\allowbreak{} 。置换不改变二次型的取值，也不改变\allowbreak{}\hskip.1em \RB{8.202}{8.294}{0.500}{%
\RG{0.000}{0.005}{10.9091}{mi}{65}%
}\hskip.1em\allowbreak{}的行列式。

对任意\allowbreak{}\hskip.1em \RB{25.353}{8.316}{2.852}{%
\RG{0.000}{0.005}{10.9091}{mi}{121}\RG{8.631}{0.005}{10.9091}{sy}{54}\RG{8.631}{0.005}{10.9091}{rm}{61}\RG{19.965}{0.005}{10.9091}{rm}{48}%
}\hskip.1em\allowbreak{} ，向末位补零并利用\allowbreak{}\hskip.1em \RB{8.508}{10.609}{0.500}{%
\RG{0.000}{0.005}{10.9091}{mi}{65}\RG{2.962}{3.545}{9.0910}{ex}{101}%
}\hskip.1em\allowbreak{} 的正定性，得\allowbreak{}\hskip.1em \RB{45.114}{9.459}{2.721}{%
\RG{0.000}{0.005}{10.9091}{mi}{121}\RG{5.599}{3.963}{7.9701}{rm}{84}\RG{11.332}{0.005}{10.9091}{mi}{67}\RG{19.754}{0.005}{10.9091}{mi}{121}\RG{28.383}{0.005}{10.9091}{mi}{62}\RG{39.726}{0.005}{10.9091}{rm}{48}%
}\hskip.1em\allowbreak{} 。所以\allowbreak{}\hskip.1em \RB{8.649}{8.185}{0.725}{%
\RG{0.000}{0.005}{10.9091}{mi}{67}%
}\hskip.1em\allowbreak{} 正定且可逆。配方得到

\[
\RB{218.227}{28.019}{3.239}{%
\RG{40.777}{22.523}{7.9701}{rm}{84}\RG{83.574}{22.523}{7.9701}{rm}{84}\RG{0.000}{18.018}{10.9091}{mi}{102}\RG{6.430}{18.018}{10.9091}{rm}{40}\RG{10.584}{18.018}{10.9091}{mi}{120}\RG{16.654}{18.018}{10.9091}{rm}{41}\RG{23.833}{18.018}{10.9091}{rm}{61}\RG{35.176}{18.018}{10.9091}{mi}{121}\RG{46.509}{18.018}{10.9091}{mi}{67}\RG{54.931}{18.018}{10.9091}{mi}{121}\RG{62.959}{18.018}{10.9091}{rm}{43}\RG{73.694}{18.018}{10.9091}{rm}{50}\RG{79.033}{18.018}{10.9091}{mi}{98}\RG{89.307}{18.018}{10.9091}{mi}{121}\RG{97.330}{18.018}{10.9091}{rm}{43}\RG{108.063}{18.018}{10.9091}{mi}{97}\RG{66.513}{4.504}{7.9701}{sy}{0}\RG{72.524}{4.504}{7.9701}{rm}{49}\RG{85.537}{4.504}{7.9701}{rm}{84}\RG{131.032}{4.504}{7.9701}{sy}{0}\RG{137.042}{4.504}{7.9701}{rm}{49}\RG{186.695}{4.504}{7.9701}{rm}{84}\RG{200.853}{4.504}{7.9701}{sy}{0}\RG{206.864}{4.504}{7.9701}{rm}{49}\RG{23.838}{-0.001}{10.9091}{rm}{61}\RG{35.173}{-0.001}{10.9091}{rm}{40}\RG{39.331}{-0.001}{10.9091}{mi}{121}\RG{47.355}{-0.001}{10.9091}{rm}{43}\RG{58.088}{-0.001}{10.9091}{mi}{67}\RG{76.842}{-0.001}{10.9091}{mi}{98}\RG{81.383}{-0.001}{10.9091}{rm}{41}\RG{91.270}{-0.001}{10.9091}{mi}{67}\RG{99.696}{-0.001}{10.9091}{rm}{40}\RG{103.849}{-0.001}{10.9091}{mi}{121}\RG{111.873}{-0.001}{10.9091}{rm}{43}\RG{122.606}{-0.001}{10.9091}{mi}{67}\RG{141.360}{-0.001}{10.9091}{mi}{98}\RG{145.902}{-0.001}{10.9091}{rm}{41}\RG{152.480}{-0.001}{10.9091}{rm}{43}\RG{163.213}{-0.001}{10.9091}{mi}{97}\RG{171.244}{-0.001}{10.9091}{sy}{0}\RG{182.154}{-0.001}{10.9091}{mi}{98}\RG{192.428}{-0.001}{10.9091}{mi}{67}\RG{211.182}{-0.001}{10.9091}{mi}{98}\RG{215.720}{-0.001}{10.9091}{mi}{58}%
}
\]
因此最小值为\allowbreak{}\hskip.1em \RB{52.865}{9.453}{0.732}{%
\RG{0.000}{-0.003}{10.9091}{mi}{97}\RG{8.031}{-0.003}{10.9091}{sy}{0}\RG{18.940}{-0.003}{10.9091}{mi}{98}\RG{23.482}{3.957}{7.9701}{rm}{84}\RG{29.214}{-0.003}{10.9091}{mi}{67}\RG{37.640}{3.957}{7.9701}{sy}{0}\RG{43.651}{3.957}{7.9701}{rm}{49}\RG{47.969}{-0.003}{10.9091}{mi}{98}%
}\hskip.1em\allowbreak{} ，且仅在\allowbreak{}\hskip.1em \RB{52.104}{9.293}{2.728}{%
\RG{0.000}{-0.003}{10.9091}{mi}{121}\RG{8.630}{-0.003}{10.9091}{rm}{61}\RG{19.969}{-0.003}{10.9091}{sy}{0}\RG{28.454}{-0.003}{10.9091}{mi}{67}\RG{36.879}{3.957}{7.9701}{sy}{0}\RG{42.890}{3.957}{7.9701}{rm}{49}\RG{47.208}{-0.003}{10.9091}{mi}{98}%
}\hskip.1em\allowbreak{} 时取得。

另一方面，直接分块相乘可得

\[
\RB{215.187}{27.562}{6.148}{%
\RG{0.000}{26.707}{9.0910}{ex}{48}\RG{59.900}{26.707}{9.0910}{ex}{49}\RG{69.673}{26.707}{9.0910}{ex}{48}\RG{102.600}{26.707}{9.0910}{ex}{49}\RG{124.924}{26.707}{9.0910}{ex}{48}\RG{202.906}{26.707}{9.0910}{ex}{49}\RG{16.848}{15.294}{10.9091}{mi}{73}\RG{54.559}{15.294}{10.9091}{rm}{48}\RG{78.552}{15.294}{10.9091}{mi}{67}\RG{97.523}{15.294}{10.9091}{mi}{98}\RG{132.878}{15.294}{10.9091}{mi}{67}\RG{174.380}{15.294}{10.9091}{mi}{98}\RG{21.559}{13.658}{7.9701}{mi}{110}\RG{26.248}{13.658}{7.9701}{sy}{0}\RG{32.257}{13.658}{7.9701}{rm}{49}\RG{0.000}{10.343}{9.0910}{ex}{64}\RG{59.900}{10.343}{9.0910}{ex}{65}\RG{69.673}{10.343}{9.0910}{ex}{64}\RG{102.600}{10.343}{9.0910}{ex}{65}\RG{124.924}{10.343}{9.0910}{ex}{64}\RG{202.906}{10.343}{9.0910}{ex}{65}\RG{113.585}{7.979}{10.9091}{rm}{61}\RG{212.679}{7.979}{10.9091}{mi}{58}\RG{20.981}{3.962}{7.9701}{rm}{84}\RG{35.140}{3.962}{7.9701}{sy}{0}\RG{41.150}{3.962}{7.9701}{rm}{49}\RG{82.170}{3.962}{7.9701}{rm}{84}\RG{173.877}{3.962}{7.9701}{rm}{84}\RG{188.036}{3.962}{7.9701}{sy}{0}\RG{194.046}{3.962}{7.9701}{rm}{49}\RG{7.955}{0.003}{10.9091}{sy}{0}\RG{16.440}{0.003}{10.9091}{mi}{98}\RG{26.714}{0.003}{10.9091}{mi}{67}\RG{54.559}{0.003}{10.9091}{rm}{49}\RG{77.627}{0.003}{10.9091}{mi}{98}\RG{96.993}{0.003}{10.9091}{mi}{97}\RG{134.420}{0.003}{10.9091}{rm}{48}\RG{150.395}{0.003}{10.9091}{mi}{97}\RG{158.427}{0.003}{10.9091}{sy}{0}\RG{169.336}{0.003}{10.9091}{mi}{98}\RG{179.610}{0.003}{10.9091}{mi}{67}\RG{198.365}{0.003}{10.9091}{mi}{98}%
}
\]
左边第一个矩阵的行列式为\allowbreak{}\hskip.1em \RB{4.973}{7.755}{0.500}{%
\RG{0.000}{-0.000}{10.9091}{rm}{49}%
}\hskip.1em\allowbreak{} ，故

\[
\RB{131.805}{11.407}{5.948}{%
\RG{65.928}{10.543}{9.0910}{ex}{16}\RG{123.868}{10.543}{9.0910}{ex}{17}\RG{94.840}{4.501}{7.9701}{rm}{84}\RG{108.999}{4.501}{7.9701}{sy}{0}\RG{115.009}{4.501}{7.9701}{rm}{49}\RG{0.000}{-0.003}{10.9091}{rm}{100}\RG{5.935}{-0.003}{10.9091}{rm}{101}\RG{10.680}{-0.003}{10.9091}{rm}{116}\RG{16.653}{-0.003}{10.9091}{mi}{65}\RG{27.692}{-0.003}{10.9091}{rm}{61}\RG{39.027}{-0.003}{10.9091}{rm}{100}\RG{44.961}{-0.003}{10.9091}{rm}{101}\RG{49.707}{-0.003}{10.9091}{rm}{116}\RG{55.684}{-0.003}{10.9091}{mi}{67}\RG{71.358}{-0.003}{10.9091}{mi}{97}\RG{79.390}{-0.003}{10.9091}{sy}{0}\RG{90.299}{-0.003}{10.9091}{mi}{98}\RG{100.573}{-0.003}{10.9091}{mi}{67}\RG{119.327}{-0.003}{10.9091}{mi}{98}\RG{129.298}{-0.003}{10.9091}{mi}{58}%
}
\]
结合\allowbreak{}\hskip.1em \RB{34.018}{8.294}{2.204}{%
\RG{0.000}{0.004}{10.9091}{mi}{67}\RG{11.661}{0.004}{10.9091}{rm}{61}\RG{23.204}{0.004}{10.9091}{mi}{65}\RG{31.213}{-1.632}{7.9701}{mi}{105}%
}\hskip.1em\allowbreak{} ，所求最小值即为\allowbreak{}\hskip.1em \RB{59.288}{8.686}{3.223}{%
\RG{0.000}{0.004}{10.9091}{rm}{100}\RG{5.935}{0.004}{10.9091}{rm}{101}\RG{10.680}{0.004}{10.9091}{rm}{116}\RG{16.654}{0.004}{10.9091}{mi}{65}\RG{24.662}{0.004}{10.9091}{mi}{61}\RG{31.821}{0.004}{10.9091}{rm}{100}\RG{37.756}{0.004}{10.9091}{rm}{101}\RG{42.501}{0.004}{10.9091}{rm}{116}\RG{48.475}{0.004}{10.9091}{mi}{65}\RG{56.483}{-1.632}{7.9701}{mi}{105}%
}\hskip.1em\allowbreak{} 。若\allowbreak{}\hskip.1em \RB{26.127}{7.759}{0.605}{%
\RG{0.000}{0.004}{10.9091}{mi}{110}\RG{9.612}{0.004}{10.9091}{rm}{61}\RG{21.154}{0.004}{10.9091}{rm}{49}%
}\hskip.1em\allowbreak{} ，按空矩阵行列式为\allowbreak{}\hskip.1em \RB{4.973}{7.759}{0.500}{%
\RG{0.000}{0.004}{10.9091}{rm}{49}%
}\hskip.1em\allowbreak{} 的约定，结论为\allowbreak{}\hskip.1em \RB{89.907}{9.293}{3.199}{%
\RG{0.000}{-0.002}{10.9091}{rm}{109}\RG{8.902}{-0.002}{10.9091}{rm}{105}\RG{11.869}{-0.002}{10.9091}{rm}{110}\RG{17.856}{-1.639}{7.9701}{mi}{120}\RG{22.206}{-2.649}{5.9776}{rm}{49}\RG{25.994}{-1.639}{7.9701}{rm}{61}\RG{32.001}{-1.639}{7.9701}{rm}{49}\RG{38.140}{-0.002}{10.9091}{mi}{97}\RG{43.748}{-1.639}{7.9701}{rm}{49}\RG{47.610}{-1.639}{7.9701}{rm}{49}\RG{51.929}{-0.002}{10.9091}{mi}{120}\RG{57.999}{3.957}{7.9701}{rm}{50}\RG{57.999}{-2.699}{7.9701}{rm}{49}\RG{65.348}{-0.002}{10.9091}{rm}{61}\RG{76.687}{-0.002}{10.9091}{mi}{97}\RG{82.294}{-1.639}{7.9701}{rm}{49}\RG{86.156}{-1.639}{7.9701}{rm}{49}%
}\hskip.1em\allowbreak{} ，同样成立。\allowbreak{}\hskip.1em \RB{8.366}{7.993}{0.502}{%
\RG{0.000}{-0.002}{10.9091}{am}{3}%
}\hskip.1em\allowbreak{}

% Original page 9.
所用基本结论的简证

Cayley–Hamilton定理。任意\allowbreak{}\hskip.1em \RB{6.020}{8.072}{0.608}{%
\RG{0.000}{0.001}{10.9091}{mi}{100}%
}\hskip.1em\allowbreak{} 阶矩阵\allowbreak{}\hskip.1em \RB{8.049}{7.897}{0.500}{%
\RG{0.000}{0.001}{10.9091}{mi}{84}%
}\hskip.1em\allowbreak{} 都满足其特征多项式\allowbreak{}\hskip.1em \RB{48.528}{8.683}{3.237}{%
\RG{0.000}{0.001}{10.9091}{mi}{31}\RG{6.692}{-1.636}{7.9701}{mi}{84}\RG{12.719}{0.001}{10.9091}{rm}{40}\RG{16.872}{0.001}{10.9091}{mi}{84}\RG{24.617}{0.001}{10.9091}{rm}{41}\RG{31.795}{0.001}{10.9091}{rm}{61}\RG{43.141}{0.001}{10.9091}{rm}{48}%
}\hskip.1em\allowbreak{} 。

证明：写

\[
\RB{108.966}{17.363}{12.854}{%
\RG{71.294}{11.816}{7.9701}{mi}{100}\RG{75.270}{11.816}{7.9701}{sy}{0}\RG{81.280}{11.816}{7.9701}{rm}{49}\RG{71.653}{9.089}{9.0910}{ex}{88}\RG{102.461}{4.502}{7.9701}{mi}{106}\RG{0.000}{-0.003}{10.9091}{rm}{97}\RG{5.346}{-0.003}{10.9091}{rm}{100}\RG{11.280}{-0.003}{10.9091}{rm}{106}\RG{14.542}{-0.003}{10.9091}{rm}{40}\RG{18.694}{-0.003}{10.9091}{mi}{116}\RG{22.556}{-0.003}{10.9091}{mi}{73}\RG{34.117}{-0.003}{10.9091}{sy}{0}\RG{45.026}{-0.003}{10.9091}{mi}{84}\RG{52.771}{-0.003}{10.9091}{rm}{41}\RG{59.949}{-0.003}{10.9091}{rm}{61}\RG{86.962}{-0.003}{10.9091}{mi}{67}\RG{98.603}{-0.003}{10.9091}{mi}{116}\RG{106.459}{-0.003}{10.9091}{mi}{58}\RG{27.262}{-1.639}{7.9701}{mi}{100}\RG{94.604}{-1.639}{7.9701}{mi}{106}\RG{71.510}{-10.871}{7.9701}{mi}{106}\RG{75.054}{-10.871}{7.9701}{rm}{61}\RG{81.062}{-10.871}{7.9701}{rm}{48}%
}
\]
由伴随矩阵恒等式，

\[
\RB{146.181}{19.583}{14.127}{%
\RG{48.395}{18.728}{9.0910}{ex}{48}\RG{91.515}{18.728}{9.0910}{ex}{49}\RG{56.349}{11.819}{7.9701}{mi}{100}\RG{60.326}{11.819}{7.9701}{sy}{0}\RG{66.336}{11.819}{7.9701}{rm}{49}\RG{56.709}{9.091}{9.0910}{ex}{88}\RG{87.516}{4.505}{7.9701}{mi}{106}\RG{48.395}{2.364}{9.0910}{ex}{64}\RG{91.515}{2.364}{9.0910}{ex}{65}\RG{0.000}{0.001}{10.9091}{mi}{31}\RG{12.719}{0.001}{10.9091}{rm}{40}\RG{16.873}{0.001}{10.9091}{mi}{116}\RG{20.730}{0.001}{10.9091}{rm}{41}\RG{24.884}{0.001}{10.9091}{mi}{73}\RG{37.056}{0.001}{10.9091}{rm}{61}\RG{72.018}{0.001}{10.9091}{mi}{67}\RG{83.659}{0.001}{10.9091}{mi}{116}\RG{101.288}{0.001}{10.9091}{rm}{40}\RG{105.442}{0.001}{10.9091}{mi}{116}\RG{109.304}{0.001}{10.9091}{mi}{73}\RG{120.865}{0.001}{10.9091}{sy}{0}\RG{131.775}{0.001}{10.9091}{mi}{84}\RG{139.519}{0.001}{10.9091}{rm}{41}\RG{143.673}{0.001}{10.9091}{mi}{58}\RG{6.692}{-1.636}{7.9701}{mi}{84}\RG{29.596}{-1.636}{7.9701}{mi}{100}\RG{79.661}{-1.636}{7.9701}{mi}{106}\RG{114.010}{-1.636}{7.9701}{mi}{100}\RG{56.566}{-10.868}{7.9701}{mi}{106}\RG{60.110}{-10.868}{7.9701}{rm}{61}\RG{66.117}{-10.868}{7.9701}{rm}{48}%
}
\]
比较多项式的矩阵系数，再将\allowbreak{}\hskip.1em \RB{8.234}{9.508}{0.608}{%
\RG{0.000}{0.001}{10.9091}{mi}{116}\RG{3.857}{3.960}{7.9701}{mi}{107}%
}\hskip.1em\allowbreak{} 的系数右乘\allowbreak{}\hskip.1em \RB{12.121}{9.508}{0.500}{%
\RG{0.000}{0.001}{10.9091}{mi}{84}\RG{7.744}{3.960}{7.9701}{mi}{107}%
}\hskip.1em\allowbreak{} 后相加，得到

\[
\RB{151.701}{17.365}{28.261}{%
\RG{43.140}{11.818}{7.9701}{mi}{100}\RG{47.116}{11.818}{7.9701}{sy}{0}\RG{53.126}{11.818}{7.9701}{rm}{49}\RG{105.399}{11.818}{7.9701}{mi}{100}\RG{109.376}{11.818}{7.9701}{sy}{0}\RG{115.385}{11.818}{7.9701}{rm}{49}\RG{43.499}{9.091}{9.0910}{ex}{88}\RG{105.758}{9.091}{9.0910}{ex}{88}\RG{78.193}{4.504}{7.9701}{mi}{106}\RG{81.737}{4.504}{7.9701}{rm}{43}\RG{87.745}{4.504}{7.9701}{rm}{49}\RG{148.198}{4.504}{7.9701}{mi}{106}\RG{0.000}{0.000}{10.9091}{mi}{31}\RG{12.719}{0.000}{10.9091}{rm}{40}\RG{16.872}{0.000}{10.9091}{mi}{84}\RG{24.617}{0.000}{10.9091}{rm}{41}\RG{31.795}{0.000}{10.9091}{rm}{61}\RG{58.808}{0.000}{10.9091}{mi}{67}\RG{70.449}{0.000}{10.9091}{mi}{84}\RG{94.490}{0.000}{10.9091}{sy}{0}\RG{121.068}{0.000}{10.9091}{mi}{67}\RG{132.709}{0.000}{10.9091}{mi}{84}\RG{140.454}{0.000}{10.9091}{mi}{84}\RG{6.691}{-1.636}{7.9701}{mi}{84}\RG{66.450}{-1.636}{7.9701}{mi}{106}\RG{128.710}{-1.636}{7.9701}{mi}{106}\RG{43.355}{-10.869}{7.9701}{mi}{106}\RG{46.900}{-10.869}{7.9701}{rm}{61}\RG{52.907}{-10.869}{7.9701}{rm}{48}\RG{105.615}{-10.869}{7.9701}{mi}{106}\RG{109.159}{-10.869}{7.9701}{rm}{61}\RG{115.167}{-10.869}{7.9701}{rm}{48}\RG{31.801}{-27.532}{10.9091}{rm}{61}\RG{43.136}{-27.532}{10.9091}{rm}{48}\RG{48.481}{-27.532}{10.9091}{mi}{58}%
}
\]
此论证不要求预先知道\allowbreak{}\hskip.1em \RB{11.146}{8.176}{3.624}{%
\RG{0.000}{-0.004}{10.9091}{mi}{67}\RG{7.642}{-1.641}{7.9701}{mi}{106}%
}\hskip.1em\allowbreak{} 与\allowbreak{}\hskip.1em \RB{8.049}{7.893}{0.504}{%
\RG{0.000}{-0.004}{10.9091}{mi}{84}%
}\hskip.1em\allowbreak{} 可交换。\allowbreak{}\hskip.1em \RB{8.366}{7.991}{0.504}{%
\RG{0.000}{-0.004}{10.9091}{am}{3}%
}\hskip.1em\allowbreak{}实对称阵的正交对角化。每个实对称阵都有一组标准正交特征向量。证明：对阶数归纳。设实对称阵为\allowbreak{}\hskip.1em \RB{7.427}{8.181}{0.728}{%
\RG{0.000}{0.001}{10.9091}{mi}{83}%
}\hskip.1em\allowbreak{} 。连续函数\allowbreak{}\hskip.1em \RB{48.150}{9.457}{0.728}{%
\RG{0.000}{0.001}{10.9091}{mi}{120}\RG{9.101}{0.001}{10.9091}{sy}{55}\RG{9.101}{0.001}{10.9091}{sy}{33}\RG{23.040}{0.001}{10.9091}{mi}{120}\RG{29.110}{3.960}{7.9701}{rm}{84}\RG{34.843}{0.001}{10.9091}{mi}{83}\RG{42.043}{0.001}{10.9091}{mi}{120}%
}\hskip.1em\allowbreak{} 在紧集\allowbreak{}\hskip.1em \RB{62.400}{8.683}{3.226}{%
\RG{0.000}{0.001}{10.9091}{sy}{102}\RG{5.455}{0.001}{10.9091}{mi}{120}\RG{14.555}{0.001}{10.9091}{rm}{58}\RG{20.553}{0.001}{10.9091}{sy}{107}\RG{26.008}{0.001}{10.9091}{mi}{120}\RG{32.078}{0.001}{10.9091}{sy}{107}\RG{40.562}{0.001}{10.9091}{rm}{61}\RG{51.897}{0.001}{10.9091}{rm}{49}\RG{57.242}{0.001}{10.9091}{sy}{103}%
}\hskip.1em\allowbreak{} 上达到最大值，设取到最大值的单位向量为\allowbreak{}\hskip.1em \RB{6.260}{5.315}{0.605}{%
\RG{0.000}{0.004}{10.9091}{mi}{117}%
}\hskip.1em\allowbreak{} 。若\allowbreak{}\hskip.1em \RB{26.362}{7.769}{0.605}{%
\RG{0.000}{0.004}{10.9091}{mi}{118}\RG{8.586}{0.004}{10.9091}{sy}{63}\RG{20.102}{0.004}{10.9091}{mi}{117}%
}\hskip.1em\allowbreak{} ，则函数

\[
\RB{119.009}{9.457}{18.088}{%
\RG{69.693}{3.960}{7.9701}{rm}{84}\RG{32.736}{0.002}{10.9091}{rm}{40}\RG{36.889}{0.002}{10.9091}{mi}{117}\RG{45.393}{0.002}{10.9091}{rm}{43}\RG{56.126}{0.002}{10.9091}{mi}{116}\RG{59.988}{0.002}{10.9091}{mi}{118}\RG{65.539}{0.002}{10.9091}{rm}{41}\RG{75.426}{0.002}{10.9091}{mi}{83}\RG{82.630}{0.002}{10.9091}{rm}{40}\RG{86.784}{0.002}{10.9091}{mi}{117}\RG{95.287}{0.002}{10.9091}{rm}{43}\RG{106.020}{0.002}{10.9091}{mi}{116}\RG{109.882}{0.002}{10.9091}{mi}{118}\RG{115.433}{0.002}{10.9091}{rm}{41}\RG{0.000}{-7.379}{10.9091}{mi}{114}\RG{5.110}{-7.379}{10.9091}{rm}{40}\RG{9.265}{-7.379}{10.9091}{mi}{116}\RG{13.122}{-7.379}{10.9091}{rm}{41}\RG{20.300}{-7.379}{10.9091}{rm}{61}\RG{74.788}{-11.710}{7.9701}{rm}{50}\RG{95.571}{-11.710}{7.9701}{rm}{50}\RG{52.433}{-14.861}{10.9091}{rm}{49}\RG{60.200}{-14.861}{10.9091}{rm}{43}\RG{70.930}{-14.861}{10.9091}{mi}{116}\RG{79.107}{-14.861}{10.9091}{sy}{107}\RG{84.561}{-14.861}{10.9091}{mi}{118}\RG{90.116}{-14.861}{10.9091}{sy}{107}\RR{32.735}{-4.651}{86.851}{0.364}%
}
\]
在\allowbreak{}\hskip.1em \RB{25.531}{7.755}{0.729}{%
\RG{0.000}{0.000}{10.9091}{mi}{116}\RG{7.849}{0.000}{10.9091}{rm}{61}\RG{20.143}{0.000}{10.9091}{rm}{48}%
}\hskip.1em\allowbreak{} 处取极大值。于是\allowbreak{}\hskip.1em \RB{89.322}{9.456}{3.238}{%
\RG{0.000}{0.000}{10.9091}{rm}{48}\RG{9.327}{0.000}{10.9091}{rm}{61}\RG{21.632}{0.000}{10.9091}{mi}{114}\RG{26.743}{3.959}{7.9701}{sy}{48}\RG{29.293}{0.000}{10.9091}{rm}{40}\RG{33.450}{0.000}{10.9091}{rm}{48}\RG{38.795}{0.000}{10.9091}{rm}{41}\RG{46.934}{0.000}{10.9091}{rm}{61}\RG{59.239}{0.000}{10.9091}{rm}{50}\RG{64.574}{0.000}{10.9091}{mi}{118}\RG{70.129}{3.959}{7.9701}{rm}{84}\RG{75.862}{0.000}{10.9091}{mi}{83}\RG{83.062}{0.000}{10.9091}{mi}{117}%
}\hskip.1em\allowbreak{} ，故\allowbreak{}\hskip.1em \RB{13.460}{8.180}{0.729}{%
\RG{0.000}{0.000}{10.9091}{mi}{83}\RG{7.200}{0.000}{10.9091}{mi}{117}%
}\hskip.1em\allowbreak{} 与\allowbreak{}\hskip.1em \RB{6.260}{5.311}{0.609}{%
\RG{0.000}{0.000}{10.9091}{mi}{117}%
}\hskip.1em\allowbreak{} 平行，\allowbreak{}\hskip.1em \RB{6.260}{5.311}{0.609}{%
\RG{0.000}{0.000}{10.9091}{mi}{117}%
}\hskip.1em\allowbreak{} 是特征向量。又对

\[
\RB{261.601}{24.515}{3.242}{%
\RG{0.000}{15.287}{10.9091}{mi}{118}\RG{8.586}{15.287}{10.9091}{sy}{63}\RG{20.102}{15.287}{10.9091}{mi}{117}\RT{26.181}{15.287}{10.9091}{，}\RG{173.523}{4.501}{7.9701}{rm}{84}\RG{227.976}{4.501}{7.9701}{rm}{84}\RG{167.443}{-0.004}{10.9091}{mi}{117}\RG{179.255}{-0.004}{10.9091}{mi}{83}\RG{186.455}{-0.004}{10.9091}{mi}{118}\RG{195.045}{-0.004}{10.9091}{rm}{61}\RG{206.380}{-0.004}{10.9091}{rm}{40}\RG{210.538}{-0.004}{10.9091}{mi}{83}\RG{217.738}{-0.004}{10.9091}{mi}{117}\RG{223.822}{-0.004}{10.9091}{rm}{41}\RG{233.708}{-0.004}{10.9091}{mi}{118}\RG{242.294}{-0.004}{10.9091}{rm}{61}\RG{253.629}{-0.004}{10.9091}{rm}{48}\RG{258.974}{-0.004}{10.9091}{mi}{59}%
}
\]
所以\allowbreak{}\hskip.1em \RB{12.157}{9.509}{0.606}{%
\RG{0.000}{0.003}{10.9091}{mi}{117}\RG{6.079}{3.961}{7.9701}{sy}{63}%
}\hskip.1em\allowbreak{} 在\allowbreak{}\hskip.1em \RB{7.427}{8.183}{0.726}{%
\RG{0.000}{0.003}{10.9091}{mi}{83}%
}\hskip.1em\allowbreak{} 下不变，且限制算子仍为实对称算子。在\allowbreak{}\hskip.1em \RB{12.157}{9.509}{0.606}{%
\RG{0.000}{0.003}{10.9091}{mi}{117}\RG{6.079}{3.961}{7.9701}{sy}{63}%
}\hskip.1em\allowbreak{} 中应用归纳假设，并添入\allowbreak{}\hskip.1em \RB{6.260}{5.314}{0.606}{%
\RG{0.000}{0.003}{10.9091}{mi}{117}%
}\hskip.1em\allowbreak{} ，即得所需标准正交特征向量基。\allowbreak{}\hskip.1em \RB{8.366}{7.991}{0.504}{%
\RG{0.000}{-0.004}{10.9091}{am}{3}%
}\hskip.1em\allowbreak{}


\end{document}
